\input{preamble}

\begin{document}

\title{Seminars}
\maketitle

\phantomsection
\label{section-phantom}

\tableofcontents

\section{A primer on symplectic groupoids}
\label{section-prime-sympl-group}

\noindent
Camilo Angulo, Jilin University.
Geometric Structures Seminar, IMPA.
February 13, 2025.

\medskip\noindent
{\bf Abstract.}
In the late 17th century, Simeon Denis
Poisson discovered an operation that
helped encoding and producing conserved
quantities. This operation is what we
now know as a Lie bracket, an infinitesimal
symmetry, but what is its global
counterpart? Symplectic groupoids are one
possible answer to this question. In
this talk, we will introduce all the basic
concepts to define symplectic
groupoids, and their role in Poisson
geometry. We will discuss key examples, and
applications. The talk will be accessible to
those familiar with differential
geometry, but no prior knowledge of groupoids
will be assumed.


\medskip\noindent
{\bf Part 1. Poisson geometry.}

Hamiltonian formalism. Recall that being a
conserved quantity \(f \in
C^\infty(X)\) is the same thing as
\(\{H,f\}=0\).

\begin{itemize} \item We have seen that it is
always possible to take quotient
of a symplectic manifold with a group action
to obtain a  Poisson
manifold.  \item Then we have found a way to
produce a symplectic foliation
from a 2-vector \(\pi \in
\mathfrak{X}^2(M):=\Lambda^{2}(TM)\).  \item

\begin{remark}
\[\left\{ \text{Lie algebra on
\(\mathfrak{g}\)}
\right\} \xrightarrow{1-1}\left\{
\text{Linear Poisson bracket on
\(\mathfrak{g}^*\)}  \right\}  \]
\end{remark}

\item We saw very nice examples
of foliation that have to do with Lie
algebras. So \(\mathfrak{b}^*_3\) which
gives the ``open book foliation", and
\(\mathfrak{e}^*\) that gives a foliation
by cylinders.
\end{itemize}

\medskip\noindent
{\bf Part 2. Symplectic realizations.}

Consider
\[(\Sigma,\omega) \xrightarrow{\mu}(M,\pi)\]


\iffalse\begin{tikzcd}
	T_p\Sigma
	\arrow[r,"d_p\mu"]\arrow[d,"\omega^\flat"]&
	T_xM\\
	T_p^*\Sigma &
	T^*_pM\arrow[l,"(d\mu)^*",swap]\arrow[u,"\pi
	^\sharp"]
\end{tikzcd}\]\fi
So that
\[\pi ^\sharp=d_p\mu \circ \omega^{-1}
\circ(d\mu)^*\]



\begin{lemma}
\(\dim (\Sigma) \geq  2 \dim (M) -
\operatorname{rk}(\pi_x)\) for all \(x \in
M\).
\end{lemma}

\begin{proof}
Done in seminar.
\end{proof}

\begin{example}
\((\mathbb{R}^2,0)\). So the map
$$
\begin{aligned}
	(\mathbb{R}^4,dx \wedge du + dy \wedge dv)
	&\longrightarrow  \mathbb{R}^2\\
	(x,y,u,v) &\longmapsto (x,y)
\end{aligned}
$$

\end{example}

\begin{exercise}
Find the symplectic realization \(\omega\) in
\((\mathbb{R}^4, \omega )
\xrightarrow{\mu}(\mathbb{R}^2,(x^2+y^2)
\partial_x
\wedge \partial_y\)

$$
\begin{aligned}
	(\mathbb{R}^4, \omega) &\longrightarrow
	(\mathbb{R}^2,(x^2+y^2) \partial_x \wedge
	\partial _y) \\
	(x,y,z,w) &\longmapsto (x,y)
\end{aligned}
$$

Also find the symplectic realization of
\(\operatorname{ a f f }^*\) with
\(\{x,y\}=x\).
\end{exercise}

\medskip\noindent
{\bf Part 3: Groupoids}

\medskip\noindent
{\bf Motivation}
\begin{enumerate}
\item Fundamental grupoid: objects are points
in the manifold and arrows are paths.
\item The action of \(S^1\) on \(S^2\) by
rotation does not give a nice quotient
because there are two singular points.
Consider the groupoid  \(S^1 \times S^2\) of
orbits. These are the arrows. The points are
just the points of \(S^2\).
\item Consider a foliation (like Möbius
foliation of circles; where there is a
singular circle, the soul). You can do the
same thing as in fundamental groupid
leafwise. Arrows then are equivalence classes
of paths that live inside leaves. This is
called monodromy of a foliation. Again
objects are points.
\item You can take a quotient of monodromy
using a connection given by the foliation.
This allows to identify certain paths between
the leaves. This is called {\it holonomy (of
a foliatiation)}. (So you can make this
notion match the usual holonomy given by
riemmanian connection.)

{\bf Upshot.}
	So the point is taking some sort of function
	space on these groupoids you can gather the
	information given by the non-smooth quotient
	(like in the case of the sphere rotating).
	So this groupoid motivation says how to get
	some structure that resembles a non smooth
	quotient.
\item Last motivation: the grid of squares
has a tone of symmetries. If you restrict to
just a few squares you loose so many
symmetries. But there's a grupoid hidden in
there that tells you what you intuition knows
about this finite grid of squares.
\end{enumerate}

\begin{definition}
A {\it groupoid} is a category where all
morphisms are invertible.
\end{definition}

So there is a kind of product among the
objects, given by composition but:
 not every two pair of objects can be
 multiplied!---only those whose
source and target match. So that's the lance
about grupoids.

Just so you make sure you understand: the
groupoid \(G\) is the morphisms of the
category. The objects are points (of a
manifold).

\begin{definition}
{\it Lie grupoid} is when the following
diagram is inside category of smooth
manifolds and \(s,t\) (source and target
maps) are submersions:
$$
\xymatrix{
G^{(2)}\ar[r]^m& G\ar[r]^ s_t & M\ar[r]^u&G
}
$$
\end{definition}

\begin{proposition}[Properties of Lie
groupoids]
\begin{itemize}
\item \(m\) is also a submersion.
\item \(i\) (inversion) is a diffeomorphism.
\item  \(u\) (unit=identity) is an embbeding.
\end{itemize}
\end{proposition}

\begin{definition}
Consider \(x \in M\) and the inverse image of
source map: \(s^{-1}(x)=\{\text{arrows that
start at \(x\)} \}\). Now if you act with
\(t\) on this set you get {\it the orbit} of
\(x\): \(\{\text{\(y \in M\) such that there
is an arrow from \(x\) to \(y\)} \}\).

And there also {\it an isotropy} \(G_x=\{g
\in G: \text{\(g\) goes from \(x\) to \(x\)}
\}\)
\end{definition}

\begin{example}
\begin{enumerate}
\item \(G=M\),  \(M=M\).
 \item Lie groups.
	\item Lie group bundles.
\item \(G=M \times M\), \(M =M\).
 \item Fundamental groupoid. Isotropy group
 is fundamental group! And orbit
 is…\clearpage universal cover!
\item Subgroupoids.
\item  Foliations.
\item If you have a normal group action of
\(G\) on \(M\) you construct a groupoid
action with groupoid \(G \times M\) and
objects \(M\), with product given on the
group part of the product. Orbits are orbits.
Isotropy group is isotropy group.
 \item Principal bundles.
\end{enumerate}
\end{example}

{\bf Back to Poisson.}

There's also a notion of Lie algebroid. Which
is strange. But the point is that to every
Poisson manifold there is a Lie algebroid.

So the question is whether there is a Lie
groupoid associated to that Lie algebroid.
Not always.

\medskip\noindent
{\bf Big question}[Fernandez and ?]
When a symplectic manifold is integrable?


(Remember that integrating means go from
algebra(oid) to group(oid).

And the point is that:

\medskip\noindent
When you {\it can} go back, you get a
\textit{ symplectic groupoid}.

\begin{remark}
Look for Kontsevich's notes on Weinstein!
\end{remark}

\begin{remark}
History: Weinstein did this intending to do
quantization (geometric?) on Poisson
manifolds. (That involves a \(C^*\)algebra
coming from the symplectic groupoid.)
\end{remark}

\begin{definition}
A {\it symplectic groupoid} is a groupoid
\(G, M\) together with  \(\omega \in
\Omega^2(M)\) such that \(\omega\) is
symplectic and multiplicative, meaning that
\(\partial  \omega =0\), that is, \(\iff
m^*\omega=\operatorname{pr}_1^*\omega+
\operatorname{pr}^*_2\omega\in
\Omega^{2}(G^{(2)})\iff\) take two vectors \(
X_k,Y_k \in TG\), and \(\omega(X_0 \star
Y_0,X_1\star Y_1) =  \omega(X_0,Y_0) + \omega
(X_1,Y_1)\)
\end{definition}

\begin{theorem}
If \((G, \omega)\) is a symplectic groupoid,
then
\begin{enumerate}
\item \(\exists !\) poisson structure on
\(M\)
\item for which \(t: G \to M\) is a
symplectic realization,
\item  Leaves are connected components of
orbits,
\item \(\mathsf{Lie}(G)\cong T^*M\) via \(X
\mapsto -u ^*(i_X \omega\).
\end{enumerate}
\end{theorem}

\begin{remark}
Look for Alejandro Cabrera, Kontsevich. There
are two things one is de Rham and the other…
from the future: simplicial?
\end{remark}

\medskip\noindent
{\bf Upshot.} The obstruction to knowing when
symplectic groupoid exists is ``variation of
symplectic form \(\omega=(1+t^2)
\omega_{S^2}\)". So how does the symplectic
group vary from leaf to leaf. So there are
two situations in which the thing doesn't
work.



\section{Neutrinos}
\label{section-neutrinos}

\noindent
Hiroshi Nunokawas, PUC-Rio.
Friday Seminar, Seminar Name.
June 27, 2025.

\medskip\noindent
{\bf Abstract.} Hiroshi will come and tell us
everything we (not) wanted to know
about these mysterious particles, and are not
going to be afraid to ask. In
particular, about the neutrino oscillation,
and the great matrices.

\bigskip\noindent

Protons and neutrons have very similar mass
of $m_p\approx 940$ MeV, while
electrons have mass of $m_e \approx 0.5$ MeV.
MeV is $10^{6}$ electronvolts,
where one eV is approximately $1.6\times
10^{-19}$ J. This is standard in high
energy physics, they use electronvolts
instead of Joules. Recall that
2J=1N$\times$1m.

Most of the things we see are protons since
they are so much larger than
electrons. But protons nor neutrons are
elementary particles.

Here's the standard model:

\medskip

\begin{tabular}{c c c c}
Quarks & $\begin{bmatrix} u\\d
\end{bmatrix}_L $ & $\begin{bmatrix} c\\s
\end{bmatrix}_L$ & $\begin{bmatrix} t\\b
\end{bmatrix}_L$\\
Leptons & $\begin{bmatrix} \nu_e\\e^-
\end{bmatrix}_L $ & $\begin{bmatrix}
\nu_\mu\\ \mu^-
\end{bmatrix}_L$ & $\begin{bmatrix}
\nu_\tau\\ \tau^- \end{bmatrix}_L$\\
Generation & 1st & 2nd & 3rd\\
Bosons & $g$, $\gamma$ , $\omega^\pm$, $z$ \\
Higgs Bosson & $H$
\end{tabular}

\medskip

It is very particular that nature repeats
itself three times. The $L$ in those
matrix actually means left-handed, and
accounts for chirality. Only left-handed
fermions have weak interaction. Right-handed
have electromagnetic interaction,
gravitational interaction, but not weak
interaction.

And then there's neutrinos. They have
negative helicity (chirality). Being
left-handed, mathematically, means to have
helicity $-1$. I think this means
that the spin is left-handed. But chirality
and helicity are not the same:
helicity is observer-dependent, and chirality
is not. Almost all neutrinos we
can see (\% 99.99999…) have negative
helicity, but not all of them.

Consider the following:
$$
n+\nu_e\leftrightarrow p+e^-
$$
But it's not completely correct: we'd better
put $d$ instead of $n$, and $u$
instead of $p$: the $d$ and $u$ quarks,
instead of the neutrons and protons.

Now consider the following reaction: a
neutron decays into a proton, an electron
and an antineutrino:
$$
n\to p+e^+\overline{\nu}_e
$$
Protons is very stable, that's why we are
here. But neutron decays in only 15
minutes.

By experimental data, we can conclude that
neutrinos' mass is consistent with
zero. But if they have mass, it should be
much smaller than the electron's $m_v
\leq 0.5$ eV. And the electron is already the
lightest fermion!

If the mass of the neutrino was zero, i.e.
$m_\nu=0$, then $v_\nu=c$ in vacuum,
which would imply that
$$
\xymatrix{
\nu_e\ar[r]^{L}&\nu_e\ar[r]&\nu_e\\
0:00&0:00&0:00
}
$$
meaning: time doesn't pass! And this means
the state of the particle cannot
change.



\section{Spheres with minimal equators}
\label{section-spheres-with-minimal-equators}

\noindent
Lucas Ambrozio, IMPA.
Differential Geometry Seminar, IMPA.
June 24, 2025.

\medskip\noindent
{\bf Abstract.} We will discuss the
connection between Riemannian metrics on the
sphere with respect to which all equators are
minimal hypersurfaces, and
algebraic curvature tensors with positive
sectional curvatures.

\bigskip\noindent

\begin{definition}
\label{definition-equator}
An {\it $(n-k)$-equator} orthogonal to $\Pi$
is
$$
\Sigma_\Pi:=\{p\in\mathbb{S}^n:\left\langle{}p,
x\right\rangle{}=0\forall
x\in\Pi\}
$$
for $\Pi$ a $k$-dimensional linear subspace
of $\mathbb{R}^{n+1}$.
\end{definition}

\begin{remark}
\label{remark-equat-total-geode-usual-metr}
Equators are totally geodesic hypersurfaces
with the usual sphere metric, which
implies they are minimal hypersurfaces.
\end{remark}

{\bf Problem.} Characterize the set
$\mathcal{M}_k(U)$ of metrics $g$ on an
open set $U\subset\mathbb{S}^n$ such that all
$k$-equators $\Sigma_\Pi$ with
 $\Sigma\cap U\neq\emptyset$ yield are
 minimal hypersurfaces $\Sigma\cap U$
on $(U,g)$.

\begin{remark}
\label{remark-why-not-Rn}
This problem can be thought of as a problem
of finding metrics on $\mathbb{R}^n$
such that $k$-planes are minimal. To see why
project the $k$-equators to
$T_p\mathbb{S}^n$ and pullback those metrics
to the sphere.
\end{remark}

\medskip\noindent

Let $g\in\mathcal{M}_k(U)$ for
$U\subset\mathbb{S}^n$ open and $n\geq 2$.

\begin{theorem}[Beltrami, Schäfli]
\label{theorem-Beltrami-Schafli}
If $k=1$ then $g$ has constant sectional
curvature.
\end{theorem}

\begin{theorem}[Hongan]
\label{theorem-Hongan}
If $1<k<n-1$ then $g$ has constant sectional
curvature.
\end{theorem}

Then Hongan also managed to produce a
classification of these metrics for
$k=n-1$.

\begin{remark}
\label{remark-linea-inver-prese-equat-mg}
If $T\in \text{GL}(n+1,\mathbb{R})$, then
$$
\begin{aligned}
\varphi: \mathbb{S}^n &\longrightarrow
\mathbb{S}^n \\
x &\longmapsto \frac{Tx}{|Tx|}
\end{aligned}
$$
is a diffeomorphism that maps $k$-equators
into $k$-equators. Thus if
$g\in\mathcal{M}_k(\mathbb{S}^n)$ then so is
$\varphi(T)^*g$.
\end{remark}

\begin{theorem}
\label{theorem-equivariant-bijection}
There exists a $\text{GL}(n+1,\mathbb{R})$
equivariant bijection
$$
\mathcal{M}_{n-1}(\mathbb{S}^n)
\leftrightarrow
\text{Curv}_+(\mathbb{R}^{n+1})
$$
where the set on the right-hand-side is the
set of algebraic curvature tensors
(also called curvature-like, i.e. with the
same symmetries as the Riemannian
curvature tensor) on $\mathbb{R}^{n+1}$ with
positive sectional curvature.

The group action is given as follows for
$T\in\text{GL}(n+1,\mathbb{R})$:
$$
(R\cdot
T)(x,y,z,w)=\frac{1}{|\det(T)|^{\frac{1}{n+1}
}}R(Tx,Ty,Tz,Tw)
$$
\end{theorem}

The point is that
$\text{Curv}_+(\mathbb{R}^{n+1})$ is an open
cone on a linear
space. Here are two simple corollaries:

\begin{lemma}
\label{lemma-corollaries}
\begin{enumerate}
\item $\mathcal{M}_{n+1}(\mathbb{S}^n)$ is in
bijection with an open positive
cone of an
$\frac{n(n+2)(n+1)^2}{12}$-dimensional real
vector space.
\item Every metric on
$\mathcal{M}_{n-1}(\mathbb{S}^n)$ is
invariant by the
antipodal map.
\end{enumerate}
\end{lemma}

{\bf Algorithm.} From any $R\in
\text{Curv}_p(\mathbb{R}^{n+1})$ we obtain a
symmetric positive definite
(positive-definitiness comes from the
positiveness
of the curvature of $R$) 2-tensor  $k_R$
satisfying
$$
(k_R)_p(v,v)=R(pv,pv)>0
$$
Also, $k_R$ has the  {\it Killing property},
i.e. that $\overline{\nabla}k(X,X,X)=0$ for
all
$X\in\mathfrak{X}(\mathbb{S}^n)$.

Then we define a positive function on
 $\mathbb{S}^n$ by
\begin{equation}
\label{equation-def-DR}
D_R:=\left(\frac{d\text{Vol}_{k_R}}{dV_g}
\right)^{\frac{4}{n-1}}
\end{equation}
and finally a Riemannian metric on
$\mathbb{S}^n$ in
$\mathcal{M}_{n-1}\mathbb{S}^n$ by
$$
g_R=\frac{1}{D_R}k_R
$$
And to go back, for $g \in
\mathcal{M}_{n-1}(\mathbb{S}^n)$ define a
positive
function on $\mathbb{S}^n$
$$
F_g:=\left(\frac{dV_g}{dV_{\overline{g}}}
\right)^{\frac{4}{n-1}}
$$
Then let $k_g:=\frac{1}{F_g}g>0$, which is a
positive definite Killing 2-tensor,
from which we may define
$R_g\in\text{Curv}_+(\mathbb{R}^{n+1})$ with
$R_g(pv,pv)=(k_g)_p(v,v)$ for all $p,v\in
T\mathbb{S}^n$.

More corollaries:
\begin{lemma}
\label{lemma-corollaries2}
\begin{enumerate}
\item $g\in\mathcal{M}_{n-1}(\mathbb{S}^n)$
is analytic because it is a Killing
tensor on $\mathbb{S}^n$, which are
well-known.
\item If $g$ is left-invatiant on
$\mathbb{S}^3$, seen as unit quaternions,
then
$g\in\mathcal{M}_2(\mathbb{S}^3$. Moreover,
for $a\geq b\geq c>0$,
$$
aL_i\odot L_i+bL_j\odot L_j+cL_k\odot L_k=k
$$
is Killing, $k>0$,  $D_k$ constant and thus
$g=\frac{1}{\text{const.}}k
\in\mathcal{M}_2(\mathbb{S}^3)$.
\item $R$ curvature tensor of
$(\mathbb{C}P^{2},g_{FS})$. We may not
remember
what's the curvature tensor, but we know the
sectional curvature is $1\leq
\text{sec}(R)\leq 4$,
$$
(k_R)_p(v,w)=\overline{g}(v,w)+3\overline{g}
(Jp,v)\overline{g}(Jp,w)
$$
and $D_R=4^{\frac{4}{3-1}}=4$, so that by
\ref{equation-def-DR} we obtain
$g_R=\frac{1}{4}k_R$, which is a Berger
metric on $\mathbb{S}^3$ with scalar
curvature 0.
\end{enumerate}
\end{lemma}

\bigskip
Now define
$$
\Sigma_V=\{p\in\mathbb{S}^n:\left\langle{}p,v\right\rangle{}=
0\}=V^{-1}(0)
$$
where $V(x):=\left\langle{}x,v\right\rangle{}$ for all
$x\in\mathbb{S}^n$. Then the normal
vector field is $\nabla V/|\nabla V|_g$, and
the second fundamental form is
given by
$$
A=\frac{1}{|\nabla V|_g}\text{Hess}_gV
$$
and its mean curvature by
\begin{equation}
\label{equation-mean-curva-level-sets-V}
H=\frac{1}{|\nabla
V|_g}\left(\Delta_gV-\text{Hess}
_gV\left(\frac{\nabla
V}{|\nabla V|},\frac{\nabla V}{|\nabla
V|}\right)\right)
\end{equation}
For every $v \in \mathbb{S}^n$ and
$p\in\Sigma_V$, we see that $H_{\Sigma_V}=0$
iff
$$
|\nabla
V|^2_g(p)\Delta_gV(p)-\text{Hess}_gV(\nabla
V(p),\nabla V(p))=0
$$
And for $\overline{g}$,
$$
\text{Hess}_{\overline{g}}V+V\overline{g}=
0\implies
\text{Hess}_{\overline{g}}V(X,X)=0
$$
for all $X\in T_p\mathbb{S}^n$ and
$p\in\Sigma_v$. Then
$$
J_g(X,Y,Z)=g(\nabla_XY-\overline{\nabla}_XY,
Z)
$$
$$
J_g(X,Y,\nabla
V)=\text{Hess}_{\overline{g}}-\text{Hess}
$$
{\bf Problems.}
\begin{enumerate}
\item Similar story for
$\mathbb{C}P^{n},\mathbb{H}P^n$?
\item Complete metrics on $\mathbb{R}^n$ with
minimal hyperplanes.
\item Find geometric invariants of metrics on
$\mathcal{M}_{n-1}(\mathbb{S}^n)$
(may be useful to study $(M^n,g)$, $n\geq 4$,
$\text{sec}>0$.
\end{enumerate}

\section{Smoothable compactified Jacobians of
nodal curves}
\label{section-smoot-compa-jacob-nodal}

\noindent
Nicola Pagani, University of Liverpool and
Bologna.
Seminar of Algebraic Geometry UFF.
August 20, 2025.

\medskip\noindent
{\bf Abstract.} Building from examples, we
introduce an abstract notion of a
'compactified Jacobian' of a nodal curve. We
then define a compactified Jacobian
to be 'smoothable' whenever it arises as the
limit of Jacobians of smooth
curves. We give a complete combinatorial
characterization of smoothable
compactified Jacobians in terms of some 'vine
stability conditions', which we
will also introduce. This is a joint work
with Fava and Viviani.

\medskip\noindent



Let $C$ be a smooth curve and $d \in
\mathbb{Z}$. Define
$$
J^d_C=\{L:\text{$L$ is a line bundle of
degree $d$}\}/\sim
$$
which is a smooth projective variety of
dimension $g(C)$.

If $C$ is nodal we still can consider
$J^d_C$.
\begin{enumerate}
\item One connected component. Then the
Jacobian is $\mathbb{P}^1$ minus two
points. This is not universally closed, so it
is not proper.
\item Two components intersecting at one
point.
The pullback of the normalization splits the
degree in
intintely many ways, givieng that $J^{-1}_C$
is an infinite set of points. This
is not of finite type, so it is not proper.
\item The curve has two components
intersecting at two points.
This gives $J^{-2}_C$,
 which is a mixture of the two former items.
 (Probably not proper too.)
\end{enumerate}

\medskip\noindent
Now consider
$$
\text{TF}_C^d=\{\mathcal{F}:
\text{coherent on $C$, torsion-free, rank-1
on $C$}\}/\sim
$$
This satisfies the existence oart of the valu
point of properness.

Now we consider the moduli. Now we consider
the ideal sheaf of the (singular?)
point(s?):
\begin{enumerate}
\item One component. The stack is proper!
\item Two components intersecting once. Now
we get stacky points,
$x=[\bullet/\mathbb{G}_m]$. These points have
generic stabilizer. The resulting
stack is not separated because a morphism of
a curve, say $\mathbb{P}^1$ minus a
point … there are infinitely many ways to
extend a morphism from this thing to a
line bundle. So you cannot include any of
these stacky points. Recall that a
sheaf is {\it simple} if its automorphism
group is $\mathbb{G}_m$.
\item  The ideal sheaf of both nodes
$\mathcal{I}(N_1,N_2)$ has a positive
dimensional automorphism
group. The stack is not proper.
\end{enumerate}

\begin{definition}
\label{definition-fine-compactified-Jacobian}
A {\it fined compactified Jacobian} of $C$ is
an open connected substack of
$\text{TF}^d(C)$ that is also proper.
\end{definition}

\begin{remark}
\label{remark-algebraic-space}
This thing is automatically an algebraic
space.
\end{remark}

\begin{definition}
\label{definition-compactified-Jacobian}
A {\it compactified Jacobian} is an open
connected of $\text{TF}^d(C)$ that
admits a proper, good moduli space.
\end{definition}

Consider the Artin stack $\mathfrak{X}
\xrightarrow{\Gamma}X$ […]
 is a {\it good
moduli space} if
\begin{enumerate}
\item Every moduli factors
$$
\xymatrix{
\mathfrak{X} \ar[r]\ar[rd]&\mathcal{I}\text{
(ACC. space)}\\
& X
}
$$
\item
$\pi_*\mathcal{O}_{\mathfrak{X}}=\mathcal{O}
_X$.
\end{enumerate}

\medskip\noindent
We expect to find a notion of stability
condition to produce these things
[…] $[\bullet/ \mathbb{G}_m]$ would be the
polystable representative.

\begin{definition}
\label{definition-smoot-compa-jacob}
A compactified Jacobian $\overline{J_C}$ is
{\it smoothable} if all smoothings
$\mathcal{C}\to \Delta=\{0,\eta\}$ (with
$\mathcal{C}_0=C$),
$$
J^d_{\mathcal{C}_\eta}\cup C \to
\overline{J_C}
$$
is proper.
\end{definition}

\medskip\noindent

\begin{definition}
\label{definition-connected}
Let $X$ be a curve.
$$
\text{BCON}(X)=\{Y \subseteq X\text{s.t.
}Y,Y^c \text{ are connected}\}
$$
\end{definition}

\begin{definition}
\label{definition-v-satbility}
A $v$-curve is a generalization of items (2)
and (3) in the lists above
[it looks like two long snakes $\sim$ that
intersect several times, and $t$ is
the number of nodes].
A {\it $v$-condition} is a pair $n=(n_1,n_2)$
such that
$$
n_1+n_2=\begin{cases}
d+1-t\qquad &\text{we say the s.c. is
nondegenerate} \\
d-t\qquad &\text{degenerate}
\end{cases}
$$
$\mathcal{F}$ on $X$ is  {\it
$n$-(semi)stable} if
$\text{deg}\mathcal{F}_{X_i}>n_i$
($\text{deg}\mathcal{F}_{X_i}\geq n_i$)
for $i=1,2$.
\end{definition}
$\mathcal{F}_{X_i}=\mathcal{F}|_{X_i}$
torsion.
$$
\text{deg}(\mathcal{F}_{X_i})+\text{deg}
(\mathcal{F}_{X_2})
=d-|\text{sing}(F)|.
$$
Then
$$
\overline{J_C}(n)=\{\mathcal{F}\text{ is
semistable}\},
$$
a smooth compact Jacobian.

\begin{definition}
\label{definition-n-degeneracy}
A {\it degeneratation} of $v$-stab. on $X$ is
$n:\text{BCON}(X) \to \mathbb{Z}$
such that
\begin{enumerate}
\item
$$
n_Y+n_{Y^c}+|Y \cap Y^c|=
\begin{cases}
d+1\qquad &\text{ we say $Y$ is
$n$-nondegenerate} \\
d\qquad &\text{$Y$ is $n$-degenerate}
\end{cases}
$$
\item $Y_i$ no pa. common component
$n_{Y_1}+n_{Y_2}+\ldots$
$$
\xymatrix{
&Y_1\ar[dr]\ar[dl]\\
Y_2\ar[ur]\ar[rr]&&Y_3\ar[ll]\ar[ul]
}
$$
\end{enumerate}
\end{definition}

\begin{theorem}[-, et al]
\label{theorem-bijec-stabi-condi-nodal}
(bijection between stability conditions and
nodal curves) The map
$$
\left\{ \substack{\text{sm. comp.} \\
\text{Jac of $X$}} \right\} \to
\left\{ \substack{\text{$v$-stab.} \\
\text{cond. of $X$}} \right\}
$$
$$
n\mapsto
\overline{J_X}(n)=\{n\text{-semistable
sheaves}\}
$$
is a bijection. (The arrow should be from
right to left!)
\end{theorem}

F. Viviani had proved it for fine compact
Jacobians.

\section{Equivariant spaces of matrices of
constant rank}
\label{section-equiv-space-matri-const}

\noindent
Ada Boralevi, France.
Algebraic Geometry Seminar, IMPA.
August 27, 2025.

\medskip\noindent
{\bf Abstract.} A space of matrices of
constant rank is a vector subspace V, say
of dimension n+1, ofthe set of matrices of
size axb over a field k, such that
any nonzero element of V has fixed rank r. It
is a classical problem to look for
different ways to construct such spaces of
matrices. In this talk I will give an
introduction up to the state of the art of
the topic, and report on my latest
joint project with D. Faenzi and D. Fratila,
where we give a classification of
all spaces of matrices of constant corank one
associated to irreducible
representation of a reductive group.

\medskip\noindent

We are interested in vector spaces $U \subset
\text{Mat}_{m,n}(\mathbb{C})$,
with $m \leq n$, of {\it constant rank}, i.e.
such that for all $f \in U$,
$r:=\text{rank}f$ is the same.

Let $\ell(r,m,n):=\text{max}\dim U:U$ is of
rank $r$.

\medskip\noindent
{\bf Questions.}
\begin{enumerate}
\item $\ell(r,m,n)=$? In general not known.
\item Find relations among $\ell, r,m$ and
$n$.
\item Construction of examples and
classification.
\end{enumerate}

\begin{example}
\label{example-what-we-know}
\begin{enumerate}
\item $\ell(1,m,n)=n$,
$$
\begin{pmatrix}
x_1&x_2&\cdots &x_n\\
\\
\\
\\
\end{pmatrix}
$$
\item $\text{rank}=2$? There are two cases
([Atkinson '83], [Eisenbud-Haus '88])
\begin{itemize}
\item Compression spaces,
$$
\begin{pmatrix}
*&*&\cdots &*\\
*&0&\cdots & 0\\
\vdots &\\
*&0&\cdots&0
\end{pmatrix}
$$
\item Skew-symmetric matrices of $3 \times
3$.
\end{itemize}
\item $\ell(r,m,n) \geq n-r+1$. Because you
can put a matrix of
$m\times n$ (with the first $r$ rows that can
have nonzero entries):
$$
\begin{pmatrix}
x_1 &  x_2 &  \cdots & & x_{n-r+1}\\
& x_1 & x_2 & \cdots &\\
\\
\\
\\
\end{pmatrix}
$$
\end{enumerate}
\end{example}

\begin{theorem}[Westwick '86]
\label{theorem-Westwick}
\begin{enumerate}
\item $n-r+1 \leq \ell(r,m,n) \leq 2n-2r+1$.
\item If $n-r+1 \not|
\frac{(m-1)!}{(r-1)!}\implies \ell$
\end{enumerate}
\end{theorem}

We can see these spaces as (subvarieties?) of
determinantal varieties
$M_n=\{f \in
\text{Mat}_{m,n}(\mathbb{C}):\text{rank}(f)
\leq  r\}$.
[Interpretation via secant varieties inside
Segre embedding.]

\medskip\noindent
Consider a map $\varphi: U \to \Hom(V,W)$.
Then $\varphi \in U^* \otimes  V^* \otimes
W$.
We get that
$\varphi$ is of constant rank if and only if
some kernel, image and cokernel are vector
bundles.

\medskip\noindent
{\bf Focus of today.} What happens when
$U,V,W$ are irreducible
representations of a complex reductive group
$G$?

\medskip\noindent
{\bf Question.} What is the natural
equivariant morphism
$$
U \to \Hom(V,W)=V^* \otimes  W
$$
of constant rank?

\medskip\noindent
Consider the case of $G=\text{SL}_2$. All the
irreducible representations (which are
self-dual) are given by
$V(m-1)\cong
\mathbb{C}[x,y]_{\text{deg}=m-1}$.

Recall the Clebsh-Gordon decomposition ($m
\leq  n$)
$$
V(m-1)\otimes
V(n-1)=\bigoplus_{i=0}^{m-1}V(n-m+2i)
$$

\begin{theorem}[B-Faenzi-Lella '22]
\label{theorem-BFL22}
$$
V(n+m-2)\hookrightarrow \Hom(V(m-1),V(n-1))
$$
is of constant rank (corank 1) if and only if
$$
n-m+2i|m-1
$$
\end{theorem}

\begin{theorem}[B. Faenzi, Fratila '25]
\label{theorem-BFF25}
Let $V(\nu)$, $V(\mu)$ and $V(\lambda)$ be
irreducible representations
of a comple reductive group $G$, with
$$
\dim(V(\mu)) \leq \dim(V(\lambda))_-
$$
If there exists a morphism of representations
$$
\varphi: V (\nu) to \Hom (V (\lambda),V(\mu))
$$
then $\varphi$ is of constant corank 1 if and
only if
there exists a simple root $\alpha_i$ such
that
\begin{enumerate}
\item $\lambda=\mu+\nu-\alpha_i$,
\item $\nu$ is a multiple of $\nu$,
\item  $\nu$ is a multple of $\omega_i$
\end{enumerate}
\end{theorem}

\section{On wrapped Floer homology barcode
entropy and hyperbolic sets}
\label{section-wrapp-floer-homol-barco}

\noindent
Rafael Fernandes, UC Santa Cruz.
Differential Geometry Seminar, IMPA.
September 4, 2025.

\medskip\noindent
{\bf Abstract.} In this talk, we will discuss
the interplay between the
wrapped Floer homology barcode and
topological entropy. The concept of barcode
entropy was introduced by Çineli, Ginzburg,
and Gürel and has been shown to be
related to the topological entropy of the
underlying dynamical system in various
settings. Specifically, we will explore how,
in the presence of a topologically
transitive, locally maximal hyperbolic set
for the Reeb flow on the boundary of
a Liouville domain, barcode entropy is
bounded below by the topological entropy
restricted to the hyperbolic set.

\medskip\noindent

Let $M^n$ be a manifold.
$\omega \in \Omega^2(M)$ is {\it symplectic}
if $d \omega=0$ and it is
nondegenerate.

\begin{example}
\label{example-Rn-is-symplectic}
$\mathbb{R}^n$ is symplectic with canonical
Darboux form.
\end{example}

Recall the definition of Hamiltonian vector
field
associated to a function $H \in C^\infty(M)$.

\begin{definition}
\label{definition-nonde-diffe}
A diffeomorphism $\varphi:M \to M$ is called
{\it non-degenerate} if
$\Phi(\varphi)\cap\Delta \subset M\times M$
(pitchfork, i.e. transversal intersection!).
\end{definition}

Let $M^{2n}$ be a closed symplectic manifold.
Arnold's conjecture says
\begin{enumerate}
\item If $\varphi=\varphi_H$ (Hamiltonian
flow) is nondegenerate, then
$$
\# \text{Fix}(\varphi_H) \geq
\sum_{i=0}^{2n}\dim H_i(M,k)=\dim H_*(M,k)
$$
\item If $\varphi=\varphi_H$ is degenerate,
then
$$
\# \text{Fix}(\varphi)\geq \text{Cl}(M)+1
$$
where $\text{Cl}(M)$ is the maximum number of
homology classes
we can add before getting to zero.
\end{enumerate}

Why do we care? Because
$$
\#
\text{Fix}(\varphi_H)
\leftrightarrow\{\text{1-periodic
orbits of $X_H$}\}
$$

Idea by Floer. Construct an invariant that
would
say something about periodic orbits.

\medskip\noindent
{\bf Question.} Can Floer theory capture
other ``dynamical information''?
(Other than the periodic orbits.)

\medskip\noindent
A {\it persistence module} is a pair
$(V,\Pi)$, where $V=\{V_t\}_{t \in
\mathbb{R}}$ is a family of
$\mathbb{F}$-vector spaces
and $\Pi=\{\Pi_{st}\}_{s \leq  t}$ is a
family of maps such that
\begin{enumerate}
\item $\Pi_{ss}=, \Pi_{ts}\circ
\Pi_{rt}=\Pi_{rs}$.
\item $\exists  s \subset \mathbb{R}$ such
that
$\Pi_{s t}$ is an isomorphism for $s,t$ in
the same connected component of
$R\setminus S$.
\item $\Pi_{s t}$ have finite rank.
\item $\exists s_0$ $V_s=\{0\}$, $s \leq
s_0$.
\item $V_t= \lim_{s \to t} V_s$ (lower
limit!!)
\end{enumerate}

\begin{theorem}
\label{theorem-persi-modul-sum-integ}
Any persistence module is a sum of integral
persistence modules,
$$
(V,\Pi) \cong \bigoplus_{I \in B(V)}F(I).
$$
\end{theorem}

\begin{example}
\label{example-hart-and-sphere}
Heart and sphere. There is a noise in the
persistence module
of the heart due to an unnecessary critical
point.
\end{example}

$(M^{2n},\omega)$ a {\it Liouville domain} is
a compact
symplectic manifold and $X \in
\mathfrak{X}(M)$
with $X \cap\partial M$ (pitchfork, i.e.
transversal intersection!)
pointing outwards
and preserved by the symplectic form, i.e.
$\mathcal{L}_X \omega=\omega$
($\omega=d \alpha$).

When we restrict $\omega$ to the boundary,
we obtain a contact form and get some
interesting dynamics.

A Lagrangian $(L,\partial L) \subset
(M,\partial M)$
is {\it asympotically conical} if
\begin{enumerate}
\item $\partial L \subset \partial M$ is
Legendrian.
\item $L \cap [1-\varepsilon,1] \times
\partial M=
[1-\varepsilon,1] \times \partial L$.
\end{enumerate}

\begin{remark}
\label{remark-Hamiltonian}
Take a Hamiltonian $H: \hat{M}\to \mathbb{R}$
such that
$$
\begin{cases}
H(r,x)=h(r)\qquad &r=1 \\
H(r,x)=rT-B\qquad &
\end{cases}
$$
then $X_H=h'(r)R_\alpha$.

For $L_0,L,A\subset$ Lagrangians, $H$ linear
at infinite,
then $A_{H}^{L_0 \to L_1}$,
$$
\begin{aligned}
A_{H}^{L_0 \to L_1}: P_{L_0\to L_1}
&\longrightarrow \mathbb{R} \\
\gamma &\longmapsto \int_0^1 \gamma^* \alpha-
\int_0^1H(x(t))dt
\end{aligned}
$$
where $P_{L_0\to L} =\{\gamma:[0,1] \to
\hat{M}:\gamma(0) \in L_0,
\gamma(1) \in L_1\}
$ is the set of chords.
\end{remark}

\begin{remark}
\label{remark-1-chords-are-crticial-points}
$\text{crit}(A_{H}^{L_0 \to
L_1})=\{\text{1-chords of $X_H$ from
 $L_0$ to $L_1$}\}$.
\end{remark}

Putting a metric on $P_{L_0 \to L_1}$ we can
consider $\varphi:\mathbb{R} \times
[0,1] \to \hat{M}$, solutions of some PDE
which is some kind
of generalization of a gradient, $- \nabla
A_{H}^{L_0 \to L_1}$.
These solutions can be put in a moduli space
$$
\tilde{\mathcal{M}}(x_-,x_+,H,J)=
\{\varphi\text{
solutions s.t. …}\}
$$
Then we define a boundary operator
$\partial$.
\begin{theorem}
\label{theorem-homology}
$\partial^2=0$
\end{theorem}

So that we have a homology, called {\it
wrapped Floer homology}
$HW^t(H,L_0,L_1,J)$

\begin{remark}
\label{remark-embedding}
We have $H \leq  K \rightsquigarrow
HW^t(H,L_0,L_1,J) \to
HW^t(K,L_0,L_1,K)$.
\end{remark}

\begin{definition}
\label{definition-direct-limit}
For $t \geq 0$
$$
HW^t(M,L_0,L_1)=\underline{\lim }_H
HW^t(H,L_0,L_1,J)
$$
\end{definition}

(Where we have taken direct limit.)

Taking direct limit of the homology,
we make sure the homology theory is
independent of the
choice of objects (I think, complex structure
and Hamiltonian)
we used to construct it.


\begin{proposition}
\label{proposition-persistence-module}
$t \to HW^t(M,L_0,L_1)$ is a peristence
module $B(M,L_0,L_1).$
\end{proposition}

Finally we can define {\it barcode entropy}.
Fix $\varepsilon>0$,
$t \geq 0$,
$$
b_\varepsilon(M,L_0,L_1,t)=
\# \{\text{ of bars in }B(M,L_0,L_1)
\text{with length $\geq \varepsilon$
and start before $t$}\}
$$
Then
$$
\bar{h}^{HW}(M_0,L_0,L_1)=
\lim_{\varepsilon \to 0} \text{lim sup}_{t
\to \infty}
\frac{\text{log}^+(b_\varepsilon(M,L_0,L_1,t)
}{t}
$$

\medskip\noindent
Consider a contact manifold
$(\Sigma,\lambda,L_0,L_1)$
and $A$ the Lagrangians. (Example raising
question of filling.)

\begin{theorem}[M '24]
\label{theorem-entro-indep-filli}
$\bar{h}^{HW}$ is independent of the filling.
\end{theorem}

\begin{theorem}[M '24]
\label{theorem-bound}
$\bar{h}^{HW}(M,L_0,L_1) \leq
h_{\text{top}}(\alpha)$.
\end{theorem}

\begin{theorem}[M '25]
\label{theorem-restriction}
Consider $(M,L_0,L_1)$.
Let $K$ be a compact topologically transitive
hyperbolic set
for the Reeb flow $\alpha$. Assume $W_\delta
^s (q) \subset \partial L_0$,
$W_\delta^s(p) \subset \partial L_1$. Then
$$
\bar{h}^{HW}(M,L_0,L_1)\geq
h_{\text{top}}(\alpha |_{K})>0.
$$
\end{theorem}
Which says it captures dynamics beyond
unconditional phenomena.
In lower dimensions these tend to coincide,
but in higher dimension we don't
know. This is related to
 ``the sup over hyperbolic sets […]''

Here's a conjecture:
$$
\text{sup}_{L_0,L_1}\bar{h}^{HW}(M,L_0,L_1)=
h_{\text{top}}(\alpha).
$$
\medskip\noindent
{\bf Extra comments.}
One of the aims is to describe topological
entropy $h_\text{top}$
using Floer theory. Theorems by
Çineli-Ginzburg-Gürel show bounds
of topological entropy and barcode entropy
(one of which is for
hyperbolic sets).

There is a notion of  {\it admissible
Hamiltonian and Reeb vector fields}
which is related to some asymptotical
behaviour  ``linear at infinity''.
I understand that admissible vector fields
give
the interesting chords for the Floer homology
construction.

\section{Revisiting cotangent bundles}
\label{section-revisiting-cotangent-bundles}

\noindent
Mieugel Cueca, KU Leuven.
Symplectic Geometry Joint Seminar, IMPA.
September 5, 2025.

\medskip\noindent
{\bf Abstract.} Cotangent bundles provide key
examples of symplectic
manifolds. On the other hand, one can think
of Lie groupoids as generalizations
of manifolds. In this context, Alan Weinstein
constructed their cotangent
bundles and proved that they are so-called
symplectic groupoids. In this talk, I
will recall this construction and explain
what happens when one replaces a Lie
groupoid with a Lie 2 (or n)-groupoid. If
time permits, I will exhibit some of
their main applications. This is joint work
with Stefano Ronchi.

\medskip\noindent

Recall the basic properties of the cotangent
bundle $T^*M$ for a symplectic manifold:
\begin{enumerate}
\item It's a vector bundle.
\item $\left\langle{}\cdot,\cdot\right\rangle{}:TM \otimes
T^*M \to \mathbb{R}_M$
is the dual pairing.
\item $\omega_{\text{can}}\in \Omega^2(T^*M)$
is symplectic.
\item $\mathcal{L}_\varepsilon
\omega_{\text{can}}=\omega_{\text{can}}$,
$\varepsilon$ Euler vector field.
\end{enumerate}

\medskip\noindent
{\bf Main goal.} Reproduce the above for Lie
$n$-groupoids.

For $n=0$ we get the above situation. For
$n=1$ [Duzud-Weinstein], [Prezun],
for $n\geq 2$ ? and we care about $n=2$.

\begin{definition}
\label{definition-groupoid}
A {\it Lie $n$-groupoid}
$\mathcal{G}:\Delta^{\text{op}}\to
\mathsf{Man}$
such that
$$
P_{e,j}:\mathcal{G}_{\ell}\to \Lambda_j^\ell
\mathcal{G}
$$
are surjective submersions $\forall  \ell,j$
are differomorphisms for $\ell > n$.
\end{definition}

\begin{remark}
\label{remark-generators}
This sort of manifolds-valued presheaf
category is generated by
$$
\begin{aligned}
d_i^\ell:\mathcal{G}_\ell \to
\mathcal{G}_{\ell-1}&\text{face maps}
0 \leq  i,j \leq  \ell\\
s_j^\ell:\mathcal{G}_\ell \to
\mathcal{G}_{\ell+1}&  \qquad
\text{degeneracies}
\end{aligned}
$$
\end{remark}

\medskip\noindent
The tangent is a functor, it satisfies
$$
T_\bullet(\mathcal{G})=T_k(\mathcal{G})=
T\mathcal{G}_k
$$
(It looks like $T$ preserves diagrams.)

\medskip\noindent
{\bf Dold-Kan.} The category $\mathsf{SVect}$
of simplicial vector spaces
has objects
$$
\xymatrix{
\mathbb{V}_\bullet&
\mathbb{V}_n\ar[r]&\cdots
\ar[r]&\mathbb{V}_2\ar[r]^{\text{3 arrows}}&
\mathbb{V}_1\ar[r]^{\text{2
arrows}}&\mathbb{V}_0
}
$$
where the $\mathbb{V}_i$ are vector spaces.

There is a functor
$$
\mathsf{SVect}\overset{N}{\to}\{\text{chain
complexes}\geq 0\}
$$
$$
\mathbb{V}_{\bullet}\to N\mathbb{V}=
\left(N_\ell \mathbb{V}\Ker
P_{\ell,\ell},\partial=d_\ell\right)
$$

\begin{theorem}[Dold-Kan]
\label{theorem-Dold-Kan}
That's an equivalence of categories. [Confirm
this!]
\end{theorem}

Those categories are monodial:
$$
(\mathsf{SVect},\otimes),\qquad
(\mathbb{V}_\bullet \otimes \mathbb{W})_\ell
=\mathbb{V}_\ell \otimes \mathbb{W}_\ell
$$
$$
(\mathsf{ch}_{\geq 0},\otimes),\qquad (V
\otimes W)_i=
\bigoplus_{\ell +k=1}V_\ell \otimes W_k
$$
And $N$ is Lax monoidal with Lax structure
given by the
Eilenberg-Zilber map, though we won't explain
the details of this.

There are duals given by internal Hom:
$$
\begin{aligned}
\mathbb{V}^{n*}&=\underline{\Hom}(\mathbb{V},
B^n\mathbb{R})
\end{aligned}
$$
Where the internal Hom is given by
$$
\underline{\Hom}(\mathbb{V},B^n\mathbb{R})
_\ell
=\Hom_{\mathsf{SVect}}(\mathbb{V}
\otimes\Delta_n[\ell],B^n\mathbb{R})
$$
for an object
$\Delta[\ell]=\mathbb{R}[\Delta[\ell]]$.


\medskip\noindent
{\bf Properties.}
\begin{enumerate}
\item $\mathbb{V}^{n*}$ isa simplicial vector
bundle.
\item $N(V^{n*})$ and $N(\mathbb{V})^*[n]$ is
a quasi isomorphism.
\item
$\left\langle{}\cdot,\cdot\right\rangle{}:\mathbb{V}\otimes
\mathbb{V}^{n*}\to
B^n\mathbb{R}$ is non-degenerate on homology.
\item $\mathbb{V}\hookrightarrow
(\mathbb{V}^{n*})^{n*}$
Mont. a eq.
\end{enumerate}

\medskip\noindent
{\bf The vector bundle case.}
$$
\text{Maps}(\Delta[i],\mathcal{G})_k=
\Hom_{\mathsf{SSet}}(\Delta[i]\times\Delta[k]
,\mathcal{G})
$$
\begin{proposition}
\label{proposition-on-an-n-Lie-groupoid}
Let $\mathcal{G}$ be a Lie $n$-groupoid.
\begin{enumerate}
\item $\text{Maps}(\Delta[i],\mathcal{G})$
Lie $n$-groupoids ME $\mathcal{G}$.
\item
$\text{Maps}(\Delta[i],\mathcal{G})_0=
\mathcal{G}_i$.
\item $\text{ev}:\Delta[i]\times
\text{Maps}(\Delta[i],\mathcal{G})\to
\mathcal{G}$.
\end{enumerate}
\end{proposition}

[Staircase looking diagram.]

\begin{definition}
\label{definition-what-is-this}
$\mathcal{G}_\bullet$.
$$
T_i^{n*}\mathcal{G}=\Hom_{\mathsf{SVect}}(1^*
\Pi_{\Delta[i]}(T\mathcal{G}),
B^n\mathbb{R}_{\mathcal{G}_i})
$$
$$
(d_j,F)_{K|d_j\mathcal{G}}(x^a)=(F_k)
|_{\mathcal{G}}(x^{\delta_ja}).
$$
\end{definition}

\begin{proposition}
\label{proposition-properties-of-that}
$\mathcal{G}$ Lie $n$-groupoid, then
$T^{n*}\mathcal{G}$ satisfy
\begin{enumerate}
\item is a vector bundle $n$-groupoid
\item dual to $T\mathcal{G}$
$$
\left\langle{}\cdot,\cdot\right\rangle{}:T\mathcal{G} \otimes
T^{n*}\mathcal{G}
\to B^n\mathbb{R}_{\mathcal{G}}
$$
non-degenerate on homology.
\item $n$-shifted symplectic
$$
T^{n*}_n\mathcal{G}\overset{p}{\to}
T^*\mathcal{G}_n
$$
and $p^*\omega_{\text{can}}$.
\end{enumerate}
\end{proposition}

[More computations I missed]

\section{Holomorphic extensions of s-proper
Lie groupoids}
\label{section-holom-exten-s-prope-lie}

\noindent
Rui L. Fernandes, Instituto Superior Técnico
(Lisboa)/University of Illinois
at Urbana-Champaign.
Symplectic Geometry Seminar, IMPA.
September 17, 2025.

\medskip\noindent
{\bf Abstract.} Every smooth manifold admits
a compatible analytic
structure, and a classical result of
Whitney–Bruhat states that any analytic
manifold has a holomorphic extension. Lie
groups also admit compatible analytic
structures, and another classical result, due
to C. Chevalley, shows that any
compact Lie group has a holomorphic extension
to a complex Lie group. D.
Martínez Torres has shown that any proper Lie
groupoid admits a compatible
analytic structure. I will discuss an
extension of the classical results of
Whitney–Bruhat and Chevalley, establishing
that any s-proper Lie groupoid has a
holomorphic extension. This talk is based on
recent joint work with Ning Jiang
(arXiv:2508.18036).

\medskip\noindent

\begin{theorem}[Whitney-Beuhat]
\label{theorem-Whitney-Beuhat}
If $M$ is analytic there exists a complex
manifold  $M_\mathbb{C}$
together with an analytic map $i: M \to
M_\mathbb{C}$
which is totally real ($T_HM_{\mathbb{C}}=TM
\oplus J(TM)$)
and
\begin{enumerate}
\item For every complex manifold $X$ and
every
analytic map $\phi:N \to X$ there exists
$\phi^*:U \to X$
holomorphic contained open $M \subset U
\subset M_\mathbb{C}$.
\item If $\psi:V \to X$ holomorphic on $M
\subset V \subset M_\mathbb{C}$ then
$\psi=\phi^*$ where $\phi=\psi \circ i$ on a
possibly
smaller open contained in $U \cap V$.
\end{enumerate}
\end{theorem}

\begin{theorem}[Chevalley]
\label{theorem-Chevalley}
If $G$ is a Lie group, there exists a complex
Lie group $G_{\mathbb{C}}$ and a morphism
$i:G \to G_\mathbb{C}$
satisfying the following universal property.
For every complex Lie group $H$ and morphism
$\phi:G \to H$ there exists a unique
holomorphic map $\phi^*:G_\mathbb{C} \to H$
such that
$\phi=\phi^* \circ i$.
\end{theorem}

Here's the construction of $G_\mathbb{C}$:

$$
\xymatrix{
N\subset \tilde{G}\ar[r]\ar[d]&G^*\ar[d]\\
\tilde{G}/N \simeq
G\ar[r]_{i}&G_\mathbb{C}=G^* /\overline{i^*N}
}
$$
where $\tilde{G}$ is the universal cover
group of $G$,
 $G^*$ is the group that integrates to
the complexification of the Lie algebra of
$G$,
i.e.
$\text{Lie}(G^*)=\mathfrak{g}_\mathbb{C}$
and $\overline{i^*(N)}$ is the smallest
closed
normal complex Lie subgroup of $G^*$
containing $i^*(N)$.

\begin{example}
\label{example-complexification-groupoid}
$$
\xymatrix{
\widehat{\text{SL}_2(\mathbb{R})}\ar[d]
\ar[dr]\\
\text{SL}_2(\mathbb{R})\ar@{^{(}->}[r]&
\text{SL}_2(\mathbb{C})
}
$$
which is a simple case,
but consider instead
$\tilde{\text{SL}_2}(\mathbb{R})\times
\mathbb{R}$
and the normal subgroup
$N=\left\langle{}a\right\rangle{}\times \left\langle{}\lambda\right\rangle{}$
for irrational $\lambda$.
Then we obtain
$$
\xymatrix{
\widehat{\text{SL}_2(\mathbb{R})}
\times\mathbb{R}\ar[r]\ar[d]
&\text{SL}_2(\mathbb{C})\times\mathbb{C}
\ar[d]\\
\widehat{\text{SL}_2(\mathbb{R})}
\times\mathbb{R}/N\simeq
G\ar[r]&
G_\mathbb{C}
}
$$
and $G_\mathbb{C}$ is 3 complex dimensions!
It is not
$\text{SL}_2(\mathbb{C})\times\mathbb{C}$.
\end{example}

\section{Everything you always wanted to know
about polygons but were too
afraid to ask}
\label{section-every-you-awlwa-wante-know}

\noindent
Alessia Mandini, UFF.
GAAG, IMPA.
September 22 and 23, 2025.

\medskip\noindent
{\bf Abstract.} Moduli spaces of polygons
form a family of Kähler
manifolds that can be constructed as a Kähler
reduction of coadjoint orbits.
These spaces have deep connections to various
areas of mathematics, including
symplectic and algebraic geometry as well as
representation theory. In this
talk, I will define these spaces and explore
their connections. I will then
discuss wall-crossing phenomena in these
spaces and demonstrate how it can be
used to determine their cohomology rings.
Finally, I will introduce the
hyperkähler analogue of these spaces, known
as hyperpolygon spaces, and describe
some of their generalizations.

\medskip\noindent



Plan of the talk
\begin{enumerate}
\item Polygons in $\mathbb{R}^3$ and
relations with
other moduli spaces.
\item Polygons in other spaces.
\end{enumerate}

Let $\alpha=(\alpha_1,\ldots,\alpha_n)\in
\mathbb{R}^n_{>0}$.
Consider $S^2$ with its usual symplectic,
Kähler form.
Also consider the product of several spheres.
There's a Hamiltonian action of
$\text{SO}(3)$ by rotations.
The moduli space of polygons is
the symplectic reduction obtained with this
Hamilatonian action,
$$
M(\alpha)=\prod
S^2_{\alpha_i}/\!/\text{SO}(3).
$$

Why is it called the space of polygons?
We think that
$$
[v_1,\ldots,v_n] \in M(\alpha) \iff
\sum_{i=1}^nv_i=0
$$
are polygons.

When smooth, $M(\alpha)$ is a
$(n-3)$-complex-dimensional
Kähler manifold. Consider
$$
\varepsilon_I(\alpha)=\sum_{i \in I}\alpha_i
-\sum_{j \in I^c}\alpha_j
$$
for some index set $I \subseteq \{1,\ldots,
n\}$.
$M(\alpha)$ is {\it smooth} if
$\varepsilon_I(\alpha)
\neq 0$ for all $I\subseteq \{1,\ldots,n\}$.
If so, we say $\alpha$ is {\it generic}.

\begin{theorem}[Kapovich-Millson]
\label{theorem-Kapovich-Millson}
$M(\alpha)$ is a complex analytic space
with (eventually) isolated singularity
(homogeneous quadratic cones).
\end{theorem}

\begin{example}
\label{examples-polygon-moduli-spaces}
\begin{enumerate}
\item $(n=3.)$ Then $M(\alpha)$ is either
empty or a point.

\item $(n=4.)$ $M(\alpha)$ is either empty or
a sphere.
\end{enumerate}
\end{example}

\begin{remark}[Hausmann-Knutson]
\label{remark-Hausmann-Knutson}
Let $M_n$ be the space of all $n$-gons
modulo rigid motions. Then $M_n$
can be equipped with a Poisson structure
for which $M(\alpha)$ are the symplectic
leaves.

\medskip\noindent
{\bf Polygons as quiver varieties.}
Consider the star-shaped quiver,
which is a distinguished point with some
points around it,
and arrows from every point to the
distinguished one.
Let the distinguished point be
$V_0=\mathbb{C}^2$ and
the rest $\mathbb{C}$. Then a representation
of this quiver
is
$$
\text{Rep}Q=\bigoplus_{i=1}^n
\Hom(\mathbb{C},\mathbb{C}^2)=\mathbb{C}^{2n}
.
$$
Now put
$$
K=(U(2) \times U(1)^n)/\Delta
$$
where $\Delta$ is the diagonal $S^1$.
This gives a Hamiltonian action on
$\mathbb{C}^{2n}$
as follows. For
$$
(A,\lambda_1,\ldots,\lambda_n)\cdot
(q_1,\ldots,q_n)
$$
we map
$q_i \mapsto  A^{-1}q_i \lambda_i$.
That is
$$
\begin{aligned}
\mu: \mathbb{C}^{2n} &\longrightarrow K^* \\
(q_1,\ldots,q_n) &\longmapsto
\left(\sum_{n=1}^n (q_iq_i^*
),\ldots,\frac{1}{2}|q_i|^2\ldots\right)
\end{aligned}
$$

\begin{theorem}[Hausmann-Knutson]
\label{theorem-Hausmann-Knutson}
$$
\mathbb{C}^{2n}/\!/_{(0,\alpha)}K=M(\alpha)
$$
\end{theorem}

\begin{proof}
Uses
$$
\begin{aligned}
\mathbb{R}^3  &\xrightarrow{\simeq}
\mathfrak{su}(2)^* \\
v_i &\longmapsto (q_iq_i^*)_0
\end{aligned}
$$
\end{proof}
\end{remark}

$$
\xymatrix{
&\mathbb{C}^{2n}\ar[dl]_{/\!/_\alpha U(1)^n}
\ar[dr]^{/\!/U(2)}\\
\prod S^2 \ar[dr]_{/\!/_\alpha \text{SO}(3)}
 &  &
 \text{Gr}(2,\mathbb{C}^n)\ar[dl]^{/\!/
 _\alpha
 U(1)^n}\\
&M(\alpha).
}
$$
So we have (I think former symplectic
reduction)
$$
\mu_{U(1)^n}:\text{Gr}(2,\mathbb{C}^n)\to
\mathbb{R}^n
$$
\medskip\noindent
{\bf Walls and wall crossing.}
The walls are
$$
W_I=\{\alpha \in
\mathbb{R}^n_{>0}:\varepsilon_I(\alpha)=0\}
$$
for $I \subseteq \{1,\ldots,n\}$.

To understand wall crossing suppose we have
a wall $W_I$, with $a^c$ in the wall,
$\alpha^+$ on one side and $\alpha^-$ on the
other.

Not that
$$
\varepsilon_I(\alpha^+)>0 \iff
\sum_{k \in I}\alpha_i^+>
\sum_{ j \in I^c}\alpha_j^+
$$
Define
$$
\begin{aligned}
M_{I^c}(\alpha^+)
&=\{(v_1,\ldots,v_n) \in M(\alpha^+):
v_i=\lambda v_j \forall i,j \in I^c,
\lambda>0\}\\
&=M(\tilde{\alpha}),\qquad \tilde{\alpha}=
\left(\alpha_{i_1},\ldots,\alpha_{i_k},
\sum_{j
\in I^c}\alpha_j\right)
\end{aligned}
$$
Then $\varepsilon_I(\alpha^+)<0$
implies $M_I(\alpha^-) \subseteq
M(\alpha^-)$.

$$
\xymatrix{
&\tilde{M} \subseteq E=M_I \times
M_{I^c}\ar[dd]^{\text{blow up}}\\
M(\alpha^+)\ar[dr]&  &  M(\alpha^-)\ar[dl]\\
&M(\alpha^c)
}
$$
\begin{remark}
\label{remark-modul-space-polyg-propo}
\begin{itemize}
\item $M(\alpha)\cong M(\sigma(\alpha))$
for $\sigma$ a permutation on the order of
sets.
\item $M(\alpha)$ is conformally
symplectomorphic
to $M(\lambda \alpha)$ for all $\lambda>0$.
\end{itemize}
\end{remark}

\begin{definition}
\label{definition-short-and-long}
Let $\alpha \in \mathbb{R}^n_{>0}$.
$I \subseteq \{1,\ldots,n\}$ is {\it short}
if $\varepsilon_I(\alpha)<0$ and {\it long}
if  $\varepsilon_I(\alpha)>0$.
\end{definition}

\begin{example}
\label{example-wall-crossing}
We did an example of wall crossing.
The relevant manifolds $M_{I^c}(\alpha)$ and
$M_I(\alpha)$
were projective spaces.
\end{example}

\medskip\noindent
Now recall our first quotient
$$
\xymatrix{
\mu_{U(1)}^{-1}(\alpha)\subseteq
\text{Gr}(2,\mathbb{C}^n)
\ar[d]^{/\!/U(1)^n}\\
M(\alpha)
}
$$
Let $c_i$ be the first Chern
classes associated to the
$n$ $S^1$-bundles above
(by taking reduction in stages).

\begin{theorem}[Haussmann-Knutson,M.]
\label{theorem-Haussmann-Knutson-M}
The $c_i$ generate the cohomology of
$M(\alpha)$.
\end{theorem}

\begin{theorem}[Guillemin-Stenberg]
\label{theorem-Guillemin-Stenberg}
In this setup, for $\alpha$ generic,
$$
H(M)=\mathbb{C}[c_1,\ldots,c_n]/\text{Ann}
(\text{Vol}(M(\alpha)))
$$
i.e., $Q(c_1,\ldots,c_n) \in
\text{Ann}(\text{Vol}M(\alpha))$
$\iff$ $Q\left(\frac{\partial }{\partial
\alpha_i},\ldots,
\frac{\partial }{\partial \alpha_1}\right)
\text{Vol}(M(\alpha))=0$
\end{theorem}

\begin{theorem}[Takakura, The Koi]
\label{theorem-Takakura-The-Koi}
$$
\text{Vol}(M(\alpha))
=-\frac{(2\pi)^{n-3}}{(n-3)!}
\sum_{I\text{
long}}(-1)^{n-|I|}\varepsilon_I(\alpha)^{n-3}
.
$$
\end{theorem}

\begin{example}
\label{example-Delta1}
According to Example
\ref{example-wall-crossing} (which I
did not copy) we find that
$\alpha_1 \in \Delta_1$ gives
$\text{Vol}(M(\alpha_1))=2\pi^3(1-2\alpha_3)
^2$
and
$$
H(M(\alpha_1))=\mathbb{C}[c_3]/(c_3^3)
$$
\end{example}

\medskip\noindent
{\bf Polygon game.}
Let $G$ be a Lie group and $\mathfrak{g}$ its
Lie algebra.
Then the co-adjoint orbits
$\mathfrak{g}^*  \supseteq
\mathcal{O}_{\xi_i}$
are symplectic manifolds with the KKS form.
Then
$$
\begin{aligned}
G\text{ acts on }\Pi \mathcal{O}_{\xi_i}
&\longrightarrow \mathfrak{g}^* \\
(A_i,\ldots\alpha_n) &\longmapsto \sum A_i
\end{aligned}
$$
Then the quotient
$$
M(\xi)=\Pi \mathcal{O}_{\xi_i}/\!/_0 G
$$
generalizes the previous construction,
which we get with $G=\text{SU}(2)$.

It could be interesting to investigate
which of the next constructions can be
generalized
to other Lie groups.

\begin{theorem}[Sotillo-Floentins-Gadihl]
\label{theorem-Sotillo}
Wall-crossing for $\text{SU}(m)$.
\end{theorem}

\medskip\noindent
{\bf Bending action.}
We consider a polygon of $n$ sides
and put some diagonals that don't intersect.
We introduce some
notion of ``bending'' that allows
to define a Hamiltonian function.
In turn, this defines a torus action on
$M(\alpha)$
and a moment map.

For $n=2$ we obtain the moment map
$$
\begin{aligned}
\mu: M(\alpha) &\longrightarrow \mathbb{R}^2
\\
p &\longmapsto (\ell_1(p),\ell(p))
\end{aligned}
$$
and we can see the moment polytope in
$\mathbb{R}^2$.

\begin{remark}
\label{remark-nonva-nonin-diago}
Any system of $(n-3)$ non-vanishing
and non-intersecting diagonals determine
a torus action of $T^{n-3}$ on $M(\alpha)$.
\end{remark}

\medskip\noindent
{\bf Relations of polygon spaces to other
moduli spaces.}
Representations of the fundamental group
of the punctured sphere in $SU(2)$, i.e.
$$
\text{Rep}(\pi_1(S^2\setminus
\{p_1,\ldots,p_n\},\text{SU}(2))/\text{SU}(2)
\cong
M(\alpha)
$$
This can be put very explicitly:
$$
\begin{aligned}
&\text{Rep}(\pi_1(S^2\setminus
\{p_1,\ldots,p_n\},\text{SU}(2))\\
&=\{(g_1,\ldots,g_n)\in
\text{SU}(2)^n:g_1\cdot\ldots\cdot
g_n=\text{Id},
t_rg_i=2\cos\pi \alpha_i\}
\end{aligned}
$$
See [Agapito-Godinho]. That's all we
will say about this example.

\medskip\noindent
Now consider the moduli space of parabolic
bundles.
Consider a holomorphic bundle of rank 2 over
$\mathbb{C}P^{2}$.
Choose some points $D=\{x_1,\ldots,x_n\}$ in
$\mathbb{C}P^{2}$
and a flag
$E_{x_i}=E^{x_i,1}\supseteq
E^{x_i,2}\supseteq \{0\}$
for every $x_i \in D$. Each of the
$E_{x_i,j}$ is isomorphic to $\mathbb{C}$
(1-dimensional).
A {\it quasi-parabolic bundle}
is an holomorphic bundle $E$
with such a choice of flag.
A {\it parabolic bundle} is a quasi-parabolic
bundle with a choice of parabolic weights
$0< \beta_1(x_i)< \beta_2(x_i)<1$.

There is also a notion of stability:
$$
\begin{aligned}
\text{pdeg}E&=\text{deg}E+ \sum_{i=1}^n
(\beta_1(x_i)+\beta_2(x_i))\\
\mu(E)&=\frac{\text{pdeg}(E)}{\text{rank}(E)}
\end{aligned}
$$
We say $E$ is {\it (semi)stable} if
$µ(E)> \mu(L)$ ($\mu(E) \geq \mu(L)$)
for any $L \subseteq E$ parabolic
subbundle.

Then we obtain that
$\mathcal{M}_{\pi,d}(\beta)$
is the moduli space of (semi)-stable
parabolic bundles on $\mathbb{P}^1$
of rank $r$ and degree $d$.
In particular $\mathcal{M}_{2,0}(\beta)$
is the moduli space of (semi)-stable
parabolic
bundles on $\mathbb{P}^1$ of rank 2 and
degree 0
holomorphically trivial.

\begin{theorem}[Jeffrey, Godinho, M.]
\label{theorem-Jeffrey-Godinho}
For generic $ \alpha$,
$M(\alpha)$ is diffeomorphic to
$\mathcal{M}_{2,0}(\beta)$
whenever
$\alpha_i=\beta_2(x_i)-\beta_1(x_i)$.
\end{theorem}

\begin{proof}[Idea of proof]
The correspondence can be made quite explicit
as a map
$$
\begin{aligned}
M(\alpha) &\longrightarrow
\mathcal{M}_{2,0}(\beta) \\
[q_1,\ldots,q_n] &\longmapsto
\substack{E=\mathbb{C}P^{1}\times\mathbb{C}^2
 \\ E_{x_i}\supset E_{x_i,1}\supset
 X_{x_i,2}\supset \{0\}}\end{aligned}
$$

\end{proof}

\medskip\noindent
Now consider the quiver variety we described
above:
$$
\text{Rep}Q=\bigoplus_{i=1}^n
\Hom(\mathbb{C},\mathbb{C}^2)=\mathbb{C}^{2n}
.
$$
Now put
$$
K=(U(2) \times U(1)^n)/\Delta
$$
where $\Delta$ is the diagonal $S^1$.

But this time put
$$
\text{Rep}\tilde{Q}=\bigoplus_{i=1}^n
\Hom(\mathbb{C},\mathbb{C}^2)
\oplus \bigoplus_{i=1}^n
\Hom(\mathbb{C}^2,\mathbb{C})
$$
We also have an action of $K$, and we get
$$
K\text{ acts on }\text{Rep}\tilde{Q} = T^*
\mathbb{C}^{2n},
$$
a so-called {\it hyperHamiltonian action}.
The quotient
$$
X(\alpha)=/\!/\!/_{\substack{(0,\alpha) \\
(0,0)}}K=
\frac{\mu_{\mathbb{R}}1(0,\alpha) \wedge
\mu^{-1}_{\mathbb{C}}(0,0)}{K}
$$
called {\it hyperpolygon space}.

Here
$$
\begin{aligned}
\mu_\mathbb{R}:T^*\mathbb{C}^{2n}&\to
\mathcal{K}^* \\
\mu_{\mathbb{C}}:T^*\mathbb{C}^{2n}&  \to
\mathcal{K}_{\mathbb{C}}^* \\
(q_1,\ldots,q_n,p_1,\ldots,p_n)&  \mapsto
(\sum_{i=1}^n(p_i,q_i)_0,\ldots)
\end{aligned}
$$
It turns out that $X(\alpha)$ is smooth
if and only if $\varepsilon_1(\alpha)=0$ for
all
$I \subseteq \{ 1,\ldots,n\}$.
When smooth,  $X(\alpha)$ is a hyperkähler
manifold
(non-compact, unfortunately),
$M(\alpha)=\{[0_{1i},0] \in X(\alpha)\}$.

\begin{theorem}[Boalch]
\label{theorem-Boalch}
$X(\alpha)$ is the moduli space of polygons
for $\text{GL}(2,\mathbb{C})$.
\end{theorem}

\medskip\noindent
{\bf Parabolic Higgs bundles.}
Let $E$ be a parabolic bundle as before, i.e.
$E \in \mathcal{M}_{0,2}(\beta)$.
A {\it Higgs field on $E$} is
$$
\phi \in H^{0}(\mathbb{P}^1,\text{SPEnd}(E)
\otimes K_{\mathbb{P}^1}(D)))
$$
where a {\it strongly parabolic endomorphism}
is  $f:E \to E$
such that $f(E_{x_i})\subseteq E_{x_{i+1}}$
for all $i$.
(Fix details!)

A {\it parabolic Higgs bundle} is (probably
the pair $(E,\phi)$).
A parabolic Higgs bundle $(E,\phi)$ is {\it
(semi)stable}
if  $\mu(E) > \mu(L)$ ($\mu(E) \geq \mu(L)$)
for all
$L$ parabolic Higgs subbundle.

Then $\mathcal{N}_{r,a}^{0,1}(\beta)$ is the
moduli
space of parabolic Higgs bundles over
$\mathbb{P}^1$
with rank $r$ and degree $d$ (with fixed
determinant,
and traceless; two notions that we will not
define here).


\begin{theorem}[Goldinho, M.; Biswar,
Florentino, Godinho, M.]
\label{theorem-GBSFG}
Let $\alpha$ be generic, let $\beta$
be such that
$\alpha_i=\beta_2(x_i)-\beta_1(x_i)$.
Then the hyperpolygon space $X(\alpha)$
is symplectomorphic to the moduli space
of parabolic HIggs bundles over
$\mathbb{P}^1$, $\beta$-stable,
holomorphically trivial (with fixed
determinant and trace-free).


\begin{proof}[Idea of proof]
We can define a correspondence
$$
\begin{aligned}
X(\alpha) &\longrightarrow \mathcal{H}(\beta)
\\
[p,q] &\longmapsto
\substack{E=\mathbb{C}P^{1}\times
\mathbb{C}^2
 \\ E_{x_i}\simeq E_{x_i,1}\supset
 E_{x_i,2\supset \{0\}}}
\end{aligned}
$$
No we can use the moment map condition
that the sum of the residues is zero,
where $\text{Res}_{x_i}\phi=(p_i,q_i)$
where $\phi$ is the unique function
that satisfies the last equality by
the residue theorem.
\end{proof}
\end{theorem}

\medskip\noindent
{\bf Wall-crossing for moduli spaces of
parabolic bundles.}
See ``Translation''from Thaddeus.
What happens is that there is an $S^1$-action
on $X(\alpha)$ given by
$$
\lambda\cdot[p,q]=[\lambda p,q].
$$
(This can be paralleled with the $S^1$-action
on $\mathcal{H}(\beta)$,
given by
$\lambda\cdot[E,\phi]=[E,\lambda\phi]$.)

This action has a moment map
$$
\mu_{S^1}([p,q])=\frac{1}{2}\sum_{i}|p_i|^2
$$
Fixed points are [Konno]
$$
X_S=\{[p,q] \in X(\alpha):
S,S^c \text{ are straight, }p_j=0 \forall j
\in S\}
$$
for all $|S|\geq 2$ short. Here $S\subseteq
\{1,\ldots,n\}$
is {\it straight} if $q_i=\lambda_j q_j$
for $\lambda_j>0$ $\forall  i,j \in S$
(which in the case of polygons means the
edges are aligned).

Let $U_S$ be the flow down from $X_S$
Then, crossing a wall  $W_I$ as defined
above,
i.e. $W_I=\{\alpha \in
\mathbb{R}^n_{>0}:\varepsilon_S(\alpha)=0\}$,
``replaces'' $U_S$ by $U_{S^c}$.

Let $\alpha$ and $\tilde{\alpha}$ be generic,
then $X(\alpha)$ is diffeomorphic to
$X(\tilde{\alpha})$.
The wall crossing for $M(\alpha)$
involves
$$
\begin{aligned}
M_S(\alpha^+)&=U_S \cap M(\alpha^+)\\
M_{S^c}(\alpha^-)&=U_{S^c}\cap M(\alpha)
\end{aligned}
$$
This is thought of as a Mukai transform.

\medskip\noindent
{\bf Generalization.}
See [Florentino, Godinho, Sotillo] for wall
crossing.
See also [Fisher] PhD thesis, [Fisher,
Rayan].
[Hausel et al], [Rayan-Schopnick].

\section{Vertex algebras and special holonomy
on quadratic Lie algebras}
\label{section-verte-algeb-speci-holon}

\noindent
Mario García Fernandez, ICMAT.
GAAG, IMPA.
September 24, 2025.

\medskip\noindent
{\bf Abstract.}
The chiral de Rham complex (CDR) is a sheaf
of vertex algebras on any smooth
manifold, introduced by Malikov, Schechtman
and Vaintrob, which provides a
formal quantization of the non-linear sigma
model in mathematical physics.
Motivated by the algebra of chiral symmetries
in two-dimensional superconformal
field theories, vertex algebra embeddings on
the CDR have been studied for
special holonomy Riemannian manifolds, thanks
mainly to the work of Heluani,
Zabzine, and collaborators, with interesting
applications to the elliptic genus.
In these lectures, we will discuss extensions
of some of these results to the
case of special holonomy manifolds with
skew-torsion. The presence of torsion
typically allows for continuous symmetries in
the geometry, with an enhanced
interplay with Lie theory and algebra, as
well as the application of techniques
from generalized geometry.


\medskip\noindent
Based on joint work with Luis Álvarez Cónsul,
Andoni De Arriba de la Hera. arXiv:2012.01851
(IMRN '24).

\medskip\noindent
{\bf Motivation.}
Let $(M^n,g)$ be a Riemannian spin manifold
with
parallel spinor $\nabla^g \varphi=0$.
Then $\text{hol}(g) \subset G_\varphi
\subseteq \text{SO}(n)$,
and $\text{Ric}(g)=0$.

There is a construction by Markov-S-V that
puts a sheaf
of vertex algebras $\mathcal{V} \to M^n$.
How to construct special embeddings of
trivial vertex algebras
in the cohomology of $\mathcal{V}$, i.e.
$\mathcal{V}_\varphi \hookrightarrow
H^*(\mathcal{V})$
by Heluani et al.

\medskip\noindent
Applications. Construction of topological
invariants.
Elliptic genus, [Borisov-L].
How to understand mirror symmetry using
vertex algebras [Borisov].
Holography [Witten].

\medskip\noindent
{\bf Geometry in algebra.}
We shall do geometry in quadratic Lie
algebras.
Recall that a {\it quadratic Lie algebra} is
 $(\mathfrak{g},[\cdot,\cdot],
 (\cdot,\cdot))$ where
$(\mathfrak{g},[\cdot,\cdot])$ is a Lie
algebra
(in this course we use $\mathbb{R}$ as base
field)
and $(\cdot,\cdot)$ is a symmetric bilinear
invariant form.

\begin{example}
\label{example-quadratic-Lie-algebra}
Let $R$ be a Lie algebra with
$\left\langle{}\cdot,\cdot\right\rangle{}_R: R \otimes R \to
\mathbb{R}$.
Take $\mathfrak{g}=R \oplus R^*$,
and pick $H \in \Lambda^{3}R^*$.
Put $H(a,b,c)=\left\langle{}[a,b],c\right\rangle{}$, and
define
$$
[v+\alpha,w+\varphi]=[v,w]-\varphi([v,-])+
\alpha([v,-])+H(v,w,-)
$$
for $v,w \in R$ and $\varphi,\alpha \in R
^*$. This turns out to be a bracket.

If you like geometry you can pick $K$
compact, $\text{Lie}K=R$,
this is ``Generalized geometry on $T K \oplus
T^* K$''
à la Hitchin.
\end{example}

\medskip\noindent
QLA:
\begin{itemize}
\item Courant algebroids /$\{*\}$.
\item  Symplectic supermanifolds.
\end{itemize}

\medskip\noindent
\begin{definition}
\label{definition-generalizaed}
\begin{enumerate}
\item A {\it (generalized) metric} on
$(\mathfrak{g},(\cdot,\cdot))$ is
$G \in \text{End}(\mathfrak{g})$ with
$G^2=\text{Id}$, $(G,G)=(\cdot,\cdot)$,

This gives $\mathfrak{g}=V_+ \oplus V_-$
where $G$ acts as identity on the
first term and as $-\text{Id}$ on the second
one.
$(G_-,-)$ is a non-degenerate pairing.
$(-,-)_{V_I}$ non-degenerate tensor.

\item A divergence $\varphi \in \mathfrak{g}
^*$.

\item A {\it connection} is
$$
D: \mathfrak{g} \to \mathfrak{g}^* \otimes
\mathfrak{g}
$$
such that $(D_ab,c)+(b,D_ac)=0$.
So that $D \in \mathfrak{g}^*  \otimes
\Lambda^{2}\mathfrak{g}$
(where we identify $\mathfrak{g}$ with its
dual using the pairing).

$D$ is  {\it compatible} with $G$ if
$[D_a,G]=0$. This condition says that $D$
splits into four operators:
$$
\xymatrix{
D^+_+&  &  D_+^-\\
&  D\ar[ul]\ar[ur]\ar[dl]\ar[dr]\\
D^+_-& & D_-^-
}
$$
\item  Let $a \in \mathfrak{g}$, $a=a_++a_-$.
A {\it generalized connection} satisfies
$D_a^+=[a_-,b_+]_+ +[a_+,b_-]_-$. (This will
happen to things
living in $\mathcal{D}^0$, see below.)
\end{enumerate}
\end{definition}

\begin{definition}
\label{definition-}
Given $D$ define
\begin{enumerate}
\item $\varphi_D(a)=-\text{tr}Da$.
\item $\text{Torsion}(AX):T_D \in
\Lambda^{2}\mathfrak{g}^*$,
$$
T_D(a,b,c)=(D_ab-D_ba-[a,b],c)+(D_ca,b).
$$
(Just copying the definition of torsion.)
\end{enumerate}
\end{definition}

\begin{lemma}
\label{lemma-}
Define
$\mathcal{D}^0(G,\varphi)=\{D:G\text{-
compatible
}, \varphi_D=\varphi\}$.
Then this is non-empty, but is not a point.
Furthermore, $\forall  D \in
\mathcal{D}^0(G,\varphi)$,
$$
D_{a_-}b_+=[a_-,b_+]_+,\qquad
D_{a_+}b_-=[a_+,b_-]_-.
$$
\end{lemma}

(Checking non-emptyness is just writing out
the equations.)

\medskip\noindent
Given $G$ consider $Cl (V_+)$,
 $a_+\cdot a_+=(a_+,a_+)$ for $a_+ \in V_+$.

Fix an irreducible representation for
$Cl(V_+)$: $S_+$.
Given  $D \in \mathcal{D}^0(G,\varphi)$:
\begin{enumerate}
\item $D_-^{\delta_+} \in V_-^* \otimes
\text{End}((S_+)\supset
V_-^*\otimes \Lambda^{2}V_+ \ni D_-^+$.
\item An operator like Dirac operator:
$$
\begin{aligned}
\underline{D}^+: S_+ &\longrightarrow S_+ \\
\zeta &\longmapsto \sum_{j=1}^{\dim
V}e_j^+\cdot D_{e_j^+}\zeta.
\end{aligned}
$$
\end{enumerate}

\begin{lemma}
\label{lemma-independendant}
$D^{S_+}$ and $\underline{D}^+$ are
independnt of
$D \in \mathcal{D}^0(G,\varphi)$.
\end{lemma}

\begin{definition}
\label{definition-Killing-spinor-equations}
$(G,\varphi,\zeta)$, $\varphi \in S_+$
satisfies
{\it Killing spinor equations} if
$$
KSE\qquad
\underbrace{D_-^{S_+}\zeta=0,}
_{\text{Gravitino
Eq.}}\qquad
\underbrace{\underline{D}^+\varphi=0
}_{\text{Dilatino Eq.}}$$
\end{definition}

Expectation.
\begin{enumerate}
\item
$\mathcal{M}=\{(G,\varphi,\zeta):KSE\}/
\mathfrak{g}$
special metric.
\item Given a solution $(G,\varphi,\zeta)$
then
$$
V_{(g,\varphi,\zeta)}\hookrightarrow
V^k(\mathfrak{g}).
$$
(Where I think $V^k(\mathfrak{g})$ is the
Kac-Moody affinization.)
\end{enumerate}

\begin{remark}
\label{remark-Lie-algebra}
In the case $T K \oplus T^*K: \mathcal{M}$
2-stack. See [Bursztyn].
\end{remark}

\begin{lemma}
\label{lemma-Ricci-tensors}
Associated to $(G,\varphi)$ there are
well-defined ``Ricci tensors''
$\text{Ric}^+ \in V_- \otimes V_+$,
$\text{Ric}^- \in V_+ \otimes V_-$.
$$
\text{Ric}^+(a_-,b_+)=\text{Tr}(C_+ \to
R_D(c_+,a_-)b_+)
$$
$$
R_D(a,b)=[D_a,D_b]_c-D_{[a,b]}c,\qquad D \in
\mathcal{D}^0(G,\varphi).
$$
\begin{remark}
\label{remark-in-geometric-setup}
In geometric setup, the vanishing of the
Ricci corresponds to
the motion equations of some physical
supersymmetry theory.
\end{remark}

\begin{proposition}
\label{proposition-KSE}
If $(G,\varphi,\zeta)$ is solution of KSE,
then
 $\text{Ric}^+_{G,\phi}=0$.
If $[\varphi,G]=0 \implies
\text{Ric}^-_{G,\varphi}=0$.
\end{proposition}
\end{lemma}

\begin{proof}
$\text{Ric}^+(a-,-)\cdot
\varphi=[\underline{D}^+,D_{a_-}^{S_+}\zeta
-D_{D_+a_-}^{S_+}\zeta=0$.
$[\varphi,G]=0 \implies
\text{Ric}^+(a_-,b_+)=\text{Ric}^-(b_+,a_-)$.
\end{proof}

\medskip\noindent
{\bf Generalized Ricci flow.}
Finding solutions to these equations.
Evaluating by
$$
G_t^{-1}\partial_tG_t=-2(\text{Ric}^+-
\text{Ric}^-)
$$
$\Hom(V_+,V_-) \oplus \Hom(V_-,V_+)$.
$G G+G - G=0$,
$G=\text{Id}_{V_+}-\text{Id}_{V_-}$.

\begin{exercise}
\label{exercise-}
\begin{enumerate}
\item Prove STE for GRF.
\item Assuming there exists a solution of
KSE, prove long
time existence and convergence. (See theorem
by Streets, Jordan, GF;
and a book by Streets, GF.)
\end{enumerate}
\end{exercise}

\begin{remark}
\label{remark-sigma-model}
For physicists the generalized Ricci flow
(GRF) is the GRF of a $2d$
$\sigma$-model with target a compact Lie
group
$K$,  $\mathfrak{g}=R \oplus R^*$
\end{remark}

Let's explain something in this situation.
consider a compact Lie group $K$ and
$\mathfrak{g}=R \oplus R^*$.
A generalized metric: $\mathfrak{g}=V_+
\oplus V_-$,
$(-,-)|_{V_\pm}$ non-degenerate,
$(-,-)|_{V_+}>0$.
$V_+=\mathcal{C}^b\{X+g(X)\}$,
$X \in R$, $g \in S^2(R^*)$, $b \in
\Lambda^{2}R^*$,
$H=H_0+\overline{\partial}$.

\medskip\noindent
{\bf Case $\dim V_+=2n$.}
Complex pure spinor $\varphi$ on $V_+$ is
equivalent $V_+ \otimes \mathbb{C}=\ell
\oplus \overline{\ell}$,
where $\ell=\{0 \in V_+ \oplus \mathbb{C}:0
\cdot\zeta=0\}$.
$(\ell,\ell)=0$,
$(\overline{\ell},\overline{\ell})=0$.

This gives a complex structure
$J_\zeta:V_+ \to V_+$, $J_\zeta=i
\text{Id}_\ell-i
\text{Id}_{\overline{\ell}}$.

\begin{lemma}
\label{lemma-solutions-KSE}
$(G,\varphi,\zeta)$ with $\zeta$ pure
satisfies KSE if and only if
\begin{enumerate}
\item $[\ell,\ell]\subset \ell$.
\item Take bars
$\{\varepsilon_j,\overline{\varepsilon}_j\}
^n_{j=0}$,
$\varepsilon_j \in \ell$,
$\overline{\varepsilon}_j\in \overline{\ell}$
$(\varepsilon_j,\overline{\varepsilon}_j) =
s_{j,k}$.
$\sum_{j=1}^n[\varepsilon_j,
\overline{\varepsilon}_j=-J_\zeta
\varphi_+$.
\end{enumerate}
\end{lemma}
(This sum is a moment map condition.)

Consider $\mathcal{L}=\{\ell \subset
\mathfrak{g} \otimes \mathbb{C}:
\dim\ell=n,\ell\cap\overline{0},(-,-)|_{\ell
\oplus \overline{\ell}}
\text{non-degenerate}\}$.
This is a complex manifold.
$T_\ell \mathcal{L}
\cong \Hom(\ell,\overline{\ell}\oplus V_-
\oplus \mathbb{C})$,
$V_- \otimes \mathbb{C}=(\ell \oplus
\overline{\ell})^\perp$
Pick a vector $\dot \ell = \dot J + \dot G
\in
\Hom(\ell,\overline{\ell}\oplus V_- \oplus
\mathbb{C})$
That is, $\dot J \in \Hom(\ell,
\overline{\ell})$,
$\dot J :V_+ \to V_+$,
$\dot J J+ J \dot J=0$. (Deformations of the
complex structure.)
$G \in \Hom(V_+,V_-) \oplus \Hom(V_-,V_+)$.

\begin{proposition}[Romero]
\label{proposition-Romero}
$\mathcal{L}$ has a pseudo-Kähler structure
preserved by the $\mathfrak{g}$-action and
there exists
a moment map
$$
\begin{aligned}
\mu: \mathcal{L} &\longrightarrow
\mathfrak{g}^* \\
\ell &\longmapsto
\frac{i}{2}\sum_{j=1}^n[\varepsilon_j,
\overline{\varepsilon}_j],-),
\end{aligned}
$$
which is the quantity we mentioned in Lemma
above.
\end{proposition}

\begin{proof}
Use the natural complex structure on the
space of complex structures
(found by Fujiki),
$$
\text{tr}_{V_+}(J \dot G_2,\dot
G_1)-\text{tr}_{V_+}(J \dot J_{\varepsilon}
\dot J_1).
$$
\end{proof}

Problem: prove that
$\mathcal{M}=\{(G,\varphi,\zeta):\zeta\text{
puse}\}/\mathfrak{g}$
is a pseudo-Kähler manifold.

\begin{enumerate}
\item Pseudo-Kähler.
\item Shifter symplectic stuff.
\end{enumerate}

\begin{example}
\label{example-SU2xSU1}
Take $K=\text{SU}(2)\times\text{U}(1)\cong
S^3\times S^1$
and $\mathfrak{g}=R \oplus R^*$.
Take generators $v_1,v_2,v_3$ of
$\text{SU}(2)$ and $v_4$ of $\text{U}(1)$.
We have $[v_2,v_3]=-v_1$, $[v_3,v_1]=-v_2$,
$[v_1,v_2]=-v_3$.
Put $H_\ell=\ell v^{123}$. $\ell \in
\mathbb{R}$. $x,a \in \mathbb{R}_{>0}$.
$$
g_{x,a}=\frac{a}{x}\left(\sum_{i=1}
^3\omega^{\otimes
2}
+x^2 (v^4)^{\otimes 2}\right)
$$
$$
V_+=\{ x+g_{xa}\}\subset \mathfrak{g}
$$
$$
I_X v_4=xv_1,\qquad U_X v_2=v_3, \qquad
\varphi=-xv_4
$$
\end{example}

\begin{exercise}
\label{exercise-solution-of-KSE}
If $\ell=\frac{a}{x} \implies $ solution of
KSE.
\end{exercise}

Furthermore $[\varphi_+,\ell] \subset \ell$
(holomorphic divergence).

The $a$ parameter is naturally complexified
by  $b=yv^{23}$.
$y+ia=z$.
$$
V=\text{log}
\left(\frac{\omega^2_{x,n}}{v^{1234}}\right)=
\text{log}a.
$$
KSV hyperbolic metric, metric on
$\mathbb{H}$.

\medskip\noindent
{\bf $\dim V_+=7$}.
A real spinor on $V_+$ is equivalent to $\phi
\in (\Lambda^{3}V_+^* )_{>0}$.
The action of $\text{GL}(\mathbb{R}^7)$ on
$\Lambda^{3}(V_+^*)$.
The space of spinors $S_+=\mathbb{R}^8
=\mathbb{R}^7\oplus\mathbb{R}
\left\langle{}\varphi\right\rangle{}$.
$\phi(x,y,z)=\left\langle{}x\cdot y\cdot z\cdot
\zeta,\zeta\right\rangle{}$.

\begin{remark}
\label{remark-7}
$v^1,\ldots,v^7 \cong \mathbb{R}^6 \oplus
\mathbb{R}$.
$\phi=(v^{12}+v^{34}+v^{56})\wedge v^7+
\text{Re}((v^1+i
v^2)\wedge(v^3+iv^4)\wedge(v^5+iv^6))$.
\end{remark}

\begin{example}
\label{example-SU2+SU2+R}
Take $R=\text{SU}(2)\oplus \text{SU}(2)\oplus
\mathbb{R}$.
$\mathfrak{g}=R \oplus R^*$.
$e,s \in \mathbb{R}$, $H=sv^{123}+\ell
v^{456}$.
$\phi=\omega \wedge \zeta+\Omega^+$,
$\zeta=\sqrt{\varepsilon/\ell}v^7$,
$\omega=\sqrt{s\ell}(v^{14}+v^{25}-v^{36})$,
$\Omega^+=\sqrt{s^3}v^{123}+\ell
\sqrt{s}v^{156}-\ell \sqrt{s}v^{345}$.
\end{example}

\section{Non-Kähler Hodge Lefschetz theory
and the Bianchi identity}
\label{section-non-kahle-hodge-lefsc-theor}

\noindent
Arpan Saha, UNICAMP.
Geometric Structures Seminar, IMPA.
September 25, 2025.

\medskip\noindent
{\bf Abstract.}

Being Kähler imposes severe constraints on
the cohomology of compact complex
manifolds such as the Hard Lefschetz
property, and the question of how far this
generalises beyond the class of Kähler
manifolds has been of great interest for
a while. In this talk, I shall report on
ongoing joint work with Mario García
Fernández and Raúl González Molina that
abstracts out the definition of a
variation of Hodge--Lefschetz structure and
provides evidence that, under
certain natural assumptions, such a structure
exists more generally on
distinguished subspaces within moduli spaces
of Bismut--Ricci-flat metrics that
are pluriclosed up to source terms. In
particular, these distinguished subspaces
may be regarded as replacements for the
Kähler cone, with affine structure
modelled on a subspace of the (1,1) Aeppli
cohomology of the compact complex
manifold.

\medskip\noindent
{\bf Broad motivation.}
As we move on the parameter space associated
to a CY manifold we encounter walls.
The Kähler cone is contained in this
parameter space, and we may cross
its walls. When we cross a wall
we obtain a birational transformation
(I think this means that the moduli
on either side are birational).
But the Kähler condition is not a
birational invariant.

\medskip\noindent
{\bf Hodge-Lefschetz theory.}
For $(X,J,\omega)$ Kähler,
the cohomology $H^{\bullet}(X)$ comes with:
\begin{itemize}
\item the Lefschetz operator $L$,
given by $[\omega]\wedge-$.

\item Hodge star operator.
\item $\Lambda=* L *$ satisfying
the $\mathfrak{sl}_2$ relations:
$$
[L,\Lambda]=H,\qquad [H,L]=2L,\qquad
[H,\Lambda]=-2\Lambda.
$$
These relations are equivalent to
the Hard Lefschetz property that
$$
L^q:H^{d-q} \to H^{d+q}.
$$
\item The Poincaré pairing
(which is a different pairing from Lefschetz)
of the form
$$
H^{d-q}\cong (H^{d+q})
$$
given by integration.

\item Hodge-Riemann bilinear relations
which is a positive definite form given by
$\int_X\cdot \wedge *\cdot$.

\end{itemize}

\noindent
Moving about on the Kähler cone gives
a variation of Hodge-Lefschetz structure.

\begin{definition}
\label{definition-varia-hodge-lefsc-struc}
A {\it variation of Hodge Lefschetz structure
(VHLS)}
of weight $d$ over a manifold $K$ consists of
\begin{itemize}
\item a graded real vector bundle
$E=\bigoplus_{q=0}^d E^q \to K$
with $\text{rk}_\mathbb{R}E^0=1$.

\item An isomorphism $E^0 \otimes TK
\xrightarrow{\simeq}E^1$
(this is reminiscent of the usual definition
of variations of Hodge structures, and
we will not use it in this talk).

\item Endomorphism fields
$L, \Lambda,H \in
\mathcal{A}^0(K,\text{End}(E))$
satisfying the $\mathfrak{sl}_2$ relations
that
$$
[L,\Lambda]=H,\qquad [H,L]=2L,\qquad
[H,\Lambda]=-2\Lambda.
$$
and that
$$
H|_{E_q}=(2q-d)|_{\text{id}_{E^q}}.
$$
\item An involution $* \in
\mathcal{A}^0(K,\text{End}(E))$
such that $* L *=\Lambda$.

\item A nondegenerate symmetric pairing
$P \in \mathcal{A}^0(\mathcal{K},E^* \otimes
E^*)$
w.r.t. $L$ and $*$ are self-adjoint.

\item (like the Gauss-Manin connection)
a flat connection $D:
\mathcal{A}^0(K,E)\to\mathcal{A}^1(K,E)$
such that $DH=0=DP$.
\end{itemize}
The VHLS is said to be {\it positive} if
$P(*-,-)$
is positive definite.
\end{definition}

\noindent
{\bf Calabi's dream beyond Kähler manifolds.}
Let $(X,\Omega)$ be a compact Calabi-Yau
manifold,
where $\Omega$ is the volume form.
In general such a manifold is not Kähler.

\begin{example}[Non-Kähler manifolds]
\label{example-non-kahler}
These manifolds are not Kähler:
Heisenberg algebra quotiented by some
lattice,
torus bundle, Hopf surface.
\end{example}

\noindent
By [Yau], the Kähler cone is the moduli
space of Kähler Ricci-flat metrics.

We look for PDEs on  $X$ such that
\begin{itemize}
\item Assumption C (Calabi). Moduli
space of alutions is an open set $K$
of an affine space modelled on
a subspace $\mathbb{V}$ of
$(1,1)$-Aeppli cohomology group,
which is given by
$$
H^{1,1}_A(X)=\frac{\Ker
\partial\overline{\partial}|_{\Omega^{1,1}}}
{\text{Im}
\partial+\text{Im}\overline{\partial}}
$$

\item Assumtion D (dilaton).
A volume function $v:K\to \mathbb{R}_{>0}$
such that $\underline{d}\text{log} V$
is nowhere zero (where $\underline{d}$
denotes
Aeppli differential), and the bilinear form
$g_K:=-Dd \text{log}V$
is nondegenerate (basically the Hessian
metric).
(Ideally this would be positive definite.)

\item Assumption V (vector field).
$Z=g_K^{-1}d \text{log}V$
is such that $DZ$ is invertible
(when thought of as an endomorphism)
$g_K(Z,Z)=d \in \mathbb{Z}_{>0}$,
 $$
(D_{(DZ)^{-1}})^{d+1}V=0,
$$
\end{itemize}

\begin{remark}
\label{remark-pairing}
$$
(D_{(DZ)^{-1}})^dV=\int_X \wedge \wedge
\wedge.
$$
\end{remark}

\begin{theorem}[García Fernández-González
Molina-AS]
\label{theorem-GGA}
If assumptions C, D and V all hold with  $d
\leq 3$,
then $K$ admits a VHLS of weight $d$.
\end{theorem}

\begin{proof}[Idea of proof for $d=3$]
$$
E=\underbrace{\mathbb{R}_K}_{E^0} \oplus
\underbrace{TK}_{E^1}
\oplus \underbrace{T^*K}_{E^2} \oplus
\underbrace{\mathbb{R}_K^*}_{E^3}.
$$
Define $L$ in each of the terms of the direct
sum as
$$
L1=Z,\qquad Lv=D^2_vV,\qquad
L_\alpha=i_Z\alpha
$$
and $*$ as
$$
*1=V 1^\vee,\qquad  *v=VD^2_0
\text{log}V=Vg_K(v,-).
$$
\end{proof}

\medskip\noindent
Recall the Hull-Strominger
system, which is given by two conditions
on the connection and the metric
$$
d\left(\|\Omega\|_\omega
\frac{\omega^{n-1}}{(n-1)!}\right)=0,\qquad
F_\theta \wedge
\frac{\omega^{n-1}}{(n-1)!}=0.
$$
This determines $d
d^c\omega=\alpha\left\langle{}F_\theta \wedge
F_\theta\right\rangle{}$.



\section{Irreducibility of the Hilbert scheme
of points and the class of 2-step
ideals}
\label{section-irred-hilbe-schem-point}

\noindent
Michele Graffeo, SISSA.
Algebraic Geometry Seminar, UFF.
September 29, 2025.

\medskip\noindent
{\bf Abstract.}
Hilbert schemes of points on a
quasi-projective variety X are classical
objects
in algebraic geometry. Roughly speaking, they
parametrise ideals of a polynomial
ring with complex coefficients having finite
colength. Although Hilbert schemes
always have a distinguished component called
the smoothable component, their
geometry is quite pathological and one of the
main open problems around them
concerns their irreducibility. In a joint
work with Giovenzana, Giovenzana,
Lella we introduce the notion of 2-step
ideals. We show that the loci
parametrising these ideals are in general not
contained in the smoothable
component of the Hilbert scheme, thus
providing new examples of
extra-components.  In my seminar I will
discuss our class of ideals and relate
them to the compressed algebras considered by
Iarrobbino in the eighties.
Finally, I will show how to extend our result
to the nested setting.


\medskip\noindent

As usual, you start with a functor and
prove representativity.

\begin{definition}
\label{definition-hilbert-scheme}
Let $X$ be a smooth
quasi-projective variety and a nonegative
integer $d$.
Define
$$
\begin{aligned}
\underline{Hilb}^d: \text{Sch}_\mathbb{C}
&\longrightarrow \mathsf{Set} \\
B &\longmapsto \left\{Z \hookrightarrow B
\times X:
\substack{B\text{-flat} \\ B\text{-finite}\\
\text{length }d}\right\}
\end{aligned}
$$
\end{definition}

\begin{theorem}[Grothendieck]
\label{theorem-grothendieck}
$\underline{Hilb}^d(X)$ is represented by a
quasi-projective
scheme $Hilb^d(X)$
\end{theorem}

Idea: $\mathbb{C}$-points of $Hilb^d(X)$
are in correspondence with $Z
\overset{\subset}{\text{finite}}X$
of length.

A fat point $Z= \text{Spec}(A)$ $(A,
\mathfrak{m})$ is a
local artinian $\mathbb{C}$-algebra of finite
type.

If $X$ smooth can consider $X=\mathbb{A}^n$,
and it's the same to consider
\begin{itemize}
\item $Z \hookrightarrow \mathbb{A}^n$ of
length $d$.
\item $I \subset \mathbb{C}[x_1,\ldots,x_n]$
of
colength $d$.
\item $A= \mathbb{C}[x_1,\ldots,x_n]/I$ with
$\dim_\mathbb{C}=d$.
\end{itemize}

\begin{definition}
\label{definition-nested-Hilbert-scheme}
$X$, $0 \leq d_1 \leq \ldots
d_r=\underline{d}$.
$$
Hilb^{\underline{d}}(X)=
\{Z_1\not\hookrightarrow
\ldots
\not\hookrightarrow Z_r  \not\hookrightarrow
X:
\text{len}(Z_i)=d_i\}
$$
\end{definition}

\begin{theorem}[Hartshorne (r=0), Fogarty,
Kalpan (r>1)]
\label{theorem-HK}
$Hilb^d(X)$ is connected.
\end{theorem}

For $r=1$, $Hilb^d(\mathbb{A}^n)$ is smooth
iff
$n \leq 2$, $d \leq 3$.
If $r>1$, $n=2$, it is smooth iff $r=2$ and
$d_1=d_2=-1$.
If $r>1$ and $n>2$, it is smooth iff $r=2$,
$(d_1,d_2)=\{(1,2),(2,3)\}$.

Iror components.
\begin{itemize}
\item $(r=1)$ $n\geq 4$,
$Hilb^d(\mathbb{A}^3)$
is irreducible iff $d\leq 7$ ($\impliedby$
[E-I], $\implies $ [Mozzola].
$n=3$ irredicuble if $d \leq  11$ (8,9 Sivic;
10 Jardim et al; 11, Jelisievv et
al.)
\item $(r=2)$, $n=2$, irreducible
[G-Rosul-Sebastian].
\item $(r \geq 3)$ $n=2$, there exist $d_1
\leq \ldots \leq d_5$
such that
$Hilb^{\underline{d}}(\mathbb{A}^i)$ is
irreducible [S-R].
\end{itemize}

There is an analogy between classical in
dimension 3
and nested in dimension 2.

Schematic structure.
\begin{itemize}
\item [I] $Hilb^{21}(\mathbb{A}^4)$ has at
least a generically
non-reduced component.
\item Problem: what about $n=3$?
\item Theorem. If $Hilb^{d_1,\ldots
d_r}(\mathbb{A}^i)$ is
irreducible $\implies $
$Hilb^{(1,d_1,\ldots,d_r)}(\mathbb{A}^n$
has a generically non reduced component.
\end{itemize}

Other results.
\begin{itemize}
\item …
\end{itemize}
Hilbert scheme function.
For a local artinian $\mathbb{C}$-algebra of
finite type there is an associated graded
$\text{Gr}_\mathfrak{m}(A)
=\bigoplus_{i \geq
0}\mathfrak{m}^1/\mathfrak{m}^{i+1}$
which allows us to define the Hilbert
function
$$
\begin{aligned}
h_A: \mathbb{Z} &\longrightarrow \mathbb{N}
\\
i &\longmapsto \dim_\mathbb{C}
\mathfrak{m}_i/\mathfrak{m}_{i+1}
\end{aligned}
$$
sistability for graded $A$-modules.
The socol is the annihilator of
the ideal which gives the algebra
e.g. $\mathbb{C}[x,y](x^2,xy,y^2)$ then
$\text{Soc}(A)=\left\langle{}x,y^2\right\rangle{}_\mathbb{C}
$.

\medskip\noindent
\begin{definition}
\label{definition-smoothable-component}
The {\it smoothable component} is
 $$
V_{sm}=\overline{\{[(z_i)_{i=1}^n \in
Hilb^{\underline{d}}(\mathbb{A}^n):
Z_r\text{ is reduced}\}}
$$
\end{definition}

\begin{definition}
\label{definition-elementary-componenet}
$V \subset
Hilb^{\underline{d}}(\mathbb{A}^n)$ is
an {\it elementary component} if
$\forall [Z_1,\ldots,Z_r] \in V$
then $Z_r$ is a fat point.
\end{definition}

Instead of looking for all components of the
Hilbert scheme
we look for elementary components. This may
not be some simple;
sometimes there's infinitely many. There
is no method to find them.

\begin{theorem}[Irrabaro]
\label{theorem-I}
Every irreducible component $V \subset
Hilb^d(\mathbb{A}^n)$
is generically étale locally product of
elementary components.
\end{theorem}

\begin{definition}
\label{definition-}
Let $h:\mathbb{Z} \to \mathbb{N}$ be a
function
with finite support, $|h|=\sum_{ i \in
\mathbb{Z}}h(i)=d$. Define
\begin{itemize}
\item $H_n=\{[A] \in Hilb^d(\mathbb{A}^n):
h_A \equiv h\}$.
This is closed by semicontinuity of the
function; but can be expressed
a by a representable functor.
\item The following also has a canonical
schematic structure
$$
\begin{aligned}
\pi_n:H_n  &\longrightarrow
\mathcal{H}_n=\{[A] \in H_n: A\text{ is
graded}\} \\
A &\longmapsto A=\text{Gr}_{\mathfrak{m}}(A)
\end{aligned}
$$
Where the first is $\not\hookrightarrow $ in
the second.
\end{itemize}
\end{definition}

Recall that
$$
T_{[I]}Hilb^d(\mathbb{A}^n)=\Hom_R(I,R/I)
$$
\begin{theorem}[B-B,J,G G G L]
\label{theorem-B-B-J-G-G-G-L}
There is a graded $[I] \in \mathcal{H}_n$.
Then
$$
\begin{aligned}
T_{[I]}\mathcal{H}_n&=\Hom_R(I,R/I)_{=0}\\
T_{[I]}\pi_n([I])&='>0\\
T_{[I]}H_n&=' \geq 0.
\end{aligned}
$$
\end{theorem}

\begin{definition}
\label{definition-compressed}
$(A,\mathfrak{m}_A)$ is {\it compressed} is
for all  $(A',\mathfrak{m}')$
with $e(A)=e(A')$ and $len(A) \geq  len(A')
$.
\end{definition}

\begin{theorem}
\label{theorem-E}
$E_h \subset H_n$ is $H$-compressed
$E_h \neq  \emptyset$
\end{theorem}

\section{Normal and Restricted Tangent
Bundles of Rational
Curves in Hypersurfaces}
\label{section-restricted-tangent-bundles}

\noindent
Lucas Mioranci, IMPA.
Algebra Seminar, IMPA.
October 8, 2025.

\medskip\noindent
{\bf Abstract.}
Let $X\subset \mathbb{P}^n$ be a degree $d$
hypersurface containing a smooth
rational curve $C$ of degree $e$. The normal
bundle $N_{C/X}$ and the restricted
tangent bundle $T_X|_C$ split as direct sums
of line bundles of the form
$\bigoplus_i
\mathcal{O}_{\mathbb{P}^1}(a_i)\qquad$ called
their splitting type.
The splitting type of $N_{C/X}$ and
${T_X}|_C\;\;$ controls the local structure
of the space of rational curves and
determines how many general points of $X$ we
can interpolate by deforming the curve $C$.
By combining explicit computations
of $T_X|_C\;\;$ and an induction argument, I
classify all triples $(e,d,n)$ such
that a general degree $d$ hypersurface
$X\subset \mathbb{P}^n$ contains a
rational curve $C$ of degree $e$ whose
restricted tangent bundle ${T_X}|_C\;\;$
is balanced. The case of quadrics is
particular, in which we show that
odd-degree rational curves do not interpolate
the expected number of points on
quadric hypersurfaces.  For the normal
bundle, I compute explicit examples of
hypersurfaces $X$ for all possible splitting
types of $N_{C/X}$ when $C$ is the
rational normal curve. Additionally, for
$d\ge 3$, we compute the dimension of
the space of hypersurfaces $X$ such that
$N_{C/X}$ has a given splitting type.



\medskip\noindent

$\mathcal{O}$ with no subindex means
$\mathcal{O}_{\mathbb{P}^1}$.
$$
\xymatrix{
&0\ar[d]&0\ar[d]\\
&T_{\mathbb{P}^1}=\mathcal{O}_(2)\ar[d]\ar[r]
^=
&T_{\mathbb{P}^1}=\mathcal{O}(2)\ar[d]\\
0\ar[r]&T_{X}|_{C}\ar[d]\ar[r]&T_{\mathbb{P}
^n}|_{C}\ar[d]\ar[r]&
N_{X/\mathbb{P}^n}|_{C}\cong
\mathcal{O}(de)\ar[d]_=\ar[r]&0\\
0\ar[r]&
N_{C/X}\ar[d]\ar[r]&N_{C/\mathbb{P}^n}\ar[d]
\ar[r]&
N_{X/\mathbb{P}^n}|_{C}\cong
\mathcal{O}(de)\ar[r]&0\\
&0&0
}
$$
Throughout this talk
$X$ is a degree $d$ hypersurface in
$\mathbb{P}^n$
and $C \subset X$ is a smooth rational curve
of degree $e$.

Since $C$ is rational we can think of it
as a map $f:\mathbb{P}^1 \to X$.
Then we can pullback any bundle over $C$
back to $\mathbb{P}^1$,
and by Birkhoff-? Theorem
we know that vector bundles over
$\mathbb{P}^1$
split as
$$
E\cong \bigoplus_{i=1}^r
\mathcal{O}_{\mathbb{P}^1}(a_i),
\qquad a_1\leq \ldots \leq a_n.
$$
This decomposition is called the
{\it splitting type} of $E$.
We say $E$ is {\it balanced} if
$|a_i-a_j|\leq 1$ for all $i,j$.

Balancedness is an open condition,
i.e. if you have a family of bundles over
$\mathbb{P}^1$
and one of them is balanced, then all of them
are,
(also you can think it's a generic
condition).

We know the normal bundle $N_{C/X}$ has rank
$n-2$
and degree  $e(n-d+1)-2$.
And the tangent bundle $T_{X/C}$
has rank $n-1$ and degree $e(n-d+1)$.

\begin{example}
\label{example-conic}
A conic $C \subset \mathbb{P}^3$.
Since a conic in $\mathbb{P}^3$ is contained
in a plane,
this obliges to have a ``distinguished normal
direction'',
$\mathcal{O}(2)$. Indeed,
$$
\begin{aligned}
N_{C/\mathbb{P}^3}&\cong(N_{Q/\mathbb{P}^3}
\oplus
N_{\mathbb{P}^2/\mathbb{P}^3}|_{C}\\
&\cong(\mathcal{O}(2) \oplus
\mathcal{O}(1))|_{C}\\
&\cong\mathcal{O}(1)\oplus \mathcal{O}(2).
\end{aligned}
$$
\end{example}

\medskip\noindent
Now we discuss interpolation of general
points.
Let $p_1,\ldots,p_n \in \mathbb{P}^1$
and $x_1,\ldots, x_n \in X$
be general points.
Let $f:\mathbb{P}^1 \to X$ be such that
$f(p_i)=x_i$.

Consider the space of morphisms mapping
the  $p_i$ to $x_i$. It's known that its
tangent space
at $f$
is given by
$$
T_{[f]}\Mor(\mathbb{P}^1,X,p_i \mapsto x_i)
\cong H^0(\mathbb{P}^1,f^*T_X(-n)).
$$
\noindent
Deformations of $f$ (mapping $p_i \mapsto
x_i$)
dominate $X$ if $a_i-n_i \geq 0$.
Also, deformations of $f$ interpolate up
to $a_1+f$ general points in $X$.

\medskip\noindent
How many points can deformations of the curve
interpolate?
This means (I think) that the deformations of
the curve
are still such that $p_i \mapsto  x_i$.
``And for the normal bundle, interpolations
means that the curves contain the points''.

If $T_X|_{C}$ is balanced, $C$ can
interpolate up to
$$
\left\lfloor \frac{e(n+1-d)}{n_1}
\right\rfloor +1
$$
general points. The case of $N_{C/X}$ is
similar.

\medskip\noindent
Next we survey some relevant results.

\begin{theorem}[Coskun-Riedl, 2018]
\label{theorem-CR18}
A general Fano $X$ of degree $d \geq 2$
contains a degree $c$ rational smooth curve
 $C$ with balanced $N_{C/X}$
for every $1 \leq e \leq  n$.
\end{theorem}

The idea is to observe that the space of
curves has
balanced normal bundle.

\begin{theorem}[Ran, 2021]
\label{theorem-R21}
Extends the above for
$1 \leq  e \leq 2n-2$
and $d \geq 4$.
\end{theorem}

Idea: general hypersurface contains a curve
of balanced normal bundle.

\begin{theorem}[Ran, 2024]
\label{theorem-R24}
$X$ general Fano hypersurface.
There exist smooth rational curve with  $C$
with balanced tangent bundle $T_{X}|_{C}$
for $e$ in some arithmetic progressions.
\end{theorem}

$e$ is in general very large.
They show there are curves of arbitrarily
large
degrees with balanced tangent bundle.

\medskip\noindent
Now we describe the new results.

\begin{theorem}
\label{theorem-Lucas}
Let $X \subseteq \mathbb{P}^n$ be Fano
hypersurface of degree $3 \leq d \leq n$.

\begin{enumerate}
\item If $C \subset X$ is a rational smooth
curve
of degree $e$ with
$e \leq  \frac{n-1}{n+1-d}$,
then $T_X|_{C}$ is not balanced.

\item A general hypersurface $X$
contains rational curves $C$
of degree $e$ with balanced
$T_X|_{C}$ for every
$e> \frac{n-1}{n+1-d}$.
\end{enumerate}
\end{theorem}

\begin{theorem}
\label{theorem-lucas2}
Let $X \subseteq \mathbb{P}^n$ be a smooth
quadric hypersurface.

\begin{enumerate}
\item If $e \leq 2$ is even,
then exists a curve $C$ of degree $e$
with $T_X|_{C}\cong \mathcal{O}(e)^{n-1}$.
(I.e. we can interpolate $e+2$ points
if $e$ is even.)

\item If $e \geq 1$ is odd,
then there is no curve $C$ of degree $e$
with balanced $T_X|_C$.
The best we can do is the most balanced
bundle
before balancedness, namely
$T_X|_C \cong
\mathcal{O}(e-1)\oplus\mathcal{O}(e)
\oplus\mathcal{O}(e+1)$.
(I.e. we can only interpolate $e$ points if
$e$ is odd.)
\end{enumerate}
\end{theorem}

Next we give the main ideas in the proof.
Let $C$ be the rational normal curve
of degree $e$ in $\mathbb{P}^1$,
$$
\begin{aligned}
C:\mathbb{P}^1  &\longrightarrow \mathbb{P}^n
\\
(s:t) &\longmapsto
(s^2:s^{e-1}t:\ldots:t^e:0:\ldots:0).
\end{aligned}
$$
$X=V(F)$.

Then find some examples. Then,

\begin{proposition}
\label{proposition-proposition}
If $\mu(T_X|_{C})=\frac{e(n+1-d)}{n-1}\leq 1$
then $T_X|_{C}$ is not balanced.
\end{proposition}

\noindent
Use the proposition to show that $T_X|_C$
is not balanced if $d \leq  n$
and $e le \frac{ n-1}{n+1-d}$.

This implies Theorem \ref{theorem-Lucas}.

\begin{lemma}
\label{lemma-split-normal-bundle}
If $N_{C/X}\cong \bigoplus_i
\mathcal{O}(a_i)$
with $a_i < 4$ for all $i$,
then $T_{X}|_C \cong N_{C/X} \oplus
\mathcal{O}(2)$.
\end{lemma}

\noindent
This implies that
$\text{Ext}^1(N_{C/X},\mathcal{O}(2))=0$.

The most important step is the following,
used as the induction step:

\begin{proposition}
\label{proposition-induction}
If $T_Y|_C$ is balanced for some degree $d$
hypersurface $Y \subseteq \mathbb{P}^{n-1}$,
then there exists a degree $d$ hypersurface
$C \subseteq \mathbb{P}^n$ with balanced
$T_X|_C$.
\end{proposition}

\noindent
Combining the examples, Lemma
\ref{lemma-split-normal-bundle}
and Proposition \ref{proposition-induction},
we establish that for $e \leq
\text{max}\{2d-2,n\}$.

The following argument is used to show that
actually what we showed so far is enough for
all degrees.

\begin{lemma}
\label{lemma-smoothing-argument}
$C=C_1 \cup  C_2$, $C_1 \cong \mathbb{P}^1$.
$E$ vector bundle on $C$ such that
$E|_{C_1}$ is balanced,
$E_{C_2}$ is perfectly balanced
(numerical condition that works when
$e=n-1$).
Then if $E$ is the specialization of a family
of vector bundles $E'$ on $\mathbb{P}^1$,
then $E'$ is balanced.
\end{lemma}

\noindent
Basically, glue and smooth low degree curves
(i.e. $e \leq  \text{max}\{2d-2,n\}$)
to get all degrees


\medskip\noindent
Now we prove the result for quadrics,
i.e. $X$ is a quadraic.
For even $e$ compute examples and conclude
 $T_X|_C \cong \mathcal{O}(e)^{n-1}$.

For odd $e$, compute examples,
$T_X|_C\cong \mathcal{O}(c-1) \oplus
\mathcal{O}(2)^{3}\oplus \mathcal{O}(-1)$.

\begin{proposition}
\label{proposition-interpolation2}
$n\geq 3$, $m\leq  e+1$.
There is a degree $e$ curve $\phi_e$
interpolating
$m$ general points if and only if there is a
degree $e-2$ curve $\psi_{e-2}$ interpolating
$m-2$ general points in $X$.
\end{proposition}

\noindent
The idea is to do induction.
Suppose you can interpolate $e$ points,
then show you can interpolate $e+2$.

Chose $y_1,\ldots,y_{m-2}$ on the quadric.
Choose a line in $\mathbb{P}^{m-1}$
intersecting the quadric,
and $x_1,\ldots,x_{m-1}$ points on the line.
Join the points on the line to the points on
the quadric,
Then you get somem points of intersection
with the quadric.
Those points and the two points of
intersection of the line and the quadric
let us construct a scroll, which intersects
the quadric.
The intersection of the scroll and the
quadric
(I guess the scroll contains the lines, and
so the points $y_i$)
gives the curve passing through the $y_i$.

\section{Regular but non smooth curves of
genus 3}
\label{section-regular-but-non-smooth}

\noindent
Cesar Hilario, .
Algebra Seminar, IMPA.
October 15, 2025.

\medskip\noindent
{\bf Abstract.} Regular but non-smooth curves
are a unique feature of
geometry in positive characteristic, that
results from the fact that over an
imperfect field the notion of regularity is
weaker than the notion of
smoothness. In the setting of algebraic
geometry over an algebraically closed
field, these curves correspond to fibrations
by singular curves, which are
fibrations of relative dimension 1 whose
fibers are singular. The most famous
examples are arguably the so-called
quasi-elliptic curves and quasi-elliptic
fibrations, which play a key role in the
Bombieri-Mumford classification of
algebraic surfaces in characteristics 2 and
3. In this talk I will discuss the
case of genus 3 in characteristic 2, with an
eye towards the classification of
regular plane projective quartic curves that
become rational after base change.

\medskip\noindent
arXiv 2409.05464.

\medskip\noindent
$C$ will denote a projective geometrically
integral curve over
a not necessaily closed field $K$.
Geometricallt integral
means that $C \otimes_K \overline{K}$ is
integra.
Suppose that $p=\text{char}(K) \geq 0$
and $g=h_1(\mathcal{O}_C)$ is of genus $C$.

We define
\begin{enumerate}
\item $C$ regular iff local rings of $C$ are
regular (DVR).
\item $C$ smooth iff $C$ regular and $K
\otimes_K \overline{K}$ regular.
\end{enumerate}

\noindent
Thus it's obvious that smoothness implies
regularity.
In this talk we explain why the converse is
not true.

Note that if $K$ is perfect (e.g. $p=0$ or
$K=\overline{K}$) implies
 smooth = regular. Recall that $K$ is
 imperfect iff $p>0$
the powers of $K$, $K^p$, is a proper
subfield of $K$
(it's always a subfield but we ask it's
proper).

\medskip\noindent
Assume that $C$ is regular and $p>0$.
The genus of rthe normalization is less than
or equal to the genus of the curve,
i.e.
$$
\tilde{g}=h^1(\mathcal{O}_{\widehat{C
\otimes_K \overline{K}}})\leq g.
$$
We know that $C$ is smooth iff
$\overline{g}=g$.
So $g-\overline{g}$ is a measure of
non-smoothness.

\begin{theorem}[Tate]
\label{theorem-tate}
$\frac{p-1}{2}$ divides $g-\overline{g}$.
\end{theorem}

As a corollary, we get that if $C$ is
non-smooth,
$g-\overline{g}>0$ and so $p \leq  2g+1$.

\begin{example}
\label{example-attain}
There are cases in which the bound is sharp,
namely $C:y^2=x^p+t$, with $t \in K\setminus
K^p$.
Then $g=\frac{p-1}{2}$ and $\bar{g}=0$.
Over $\overline{K}$ rhere exists a singular
point,
$(x,y)=(-t^{1/p},0)$ which
{\bf is not visible over $K$}, that is we get
that
$t^{1/p}\not \in K$.
\end{example}

\begin{example}
\label{example-elliptic}
$\{\text{elliptic curves}\}=\{g =
\overline{g}=1,\text{$\exists K$-rational
point }\}=\{\text{smooth cubic curves with
$K$-rational point}\}$.
$y^2=x^3+4g_i x+27$.
\end{example}

\begin{example}
\label{example-quasielliptic}
$g=1$, $\overline{g}=0$, $p=2,3$,
quasielliptic curves. (Important in the
classification of surfaces
in positive characteristic.)
\end{example}

\medskip\noindent
Now we make a digression to describe
Kodaira-Enriques (over $\mathbb{C}$)
in positive characteristic. Bombieri-Mumford:
{\bf new objects} appear,
quasielliptic surfaces.

Let $k=\overline{k}$, $S$ smooth elliptic
({\it quasielliptic}) surface over $k$.
Suppose we have a fibration
where the general fiber is an elliptic curve
an elliptic curve (resp. a plane
cuspidal cubic) and the generic fiber an
elliptic (resp. quasielliptic)
curve over $k(B)$.

$$
\xymatrix{
S\ar[d]\\
B
}
$$

\medskip\noindent
Back to our main content, let's assume that
$g=3$.
$$
\xymatrix{
\{g=1,\exists K\text{-rat. pt}\}\ar[r]&\ar[l]
\left\{\substack{\text{reg. cubic curves} \\
\text{with $K$-rat pt}}\right\}
}
$$
Recall that genus-degree formula
$$
g=\frac{(d-1)(d-2)}{2}=1,
$$
So then $d=3,4$ gives  $g=3$ (right?).

No we present our main theorem, which is a
classification result:

\begin{theorem}
\label{theorem-classification-quartics}
\begin{itemize}
\item $C$ regular non-hyperelliptic curve
over $K$.
\item  $g=3$, $\overline{g}=0$.
\item $p=2$.
\end{itemize}
(Not that non-hyperllipticity of $C$ and
$g=3$ imply that $C$ is a quartic
over $K$, and that  $\overline{g}=0$ implies
$C$ is geometrically rational.)
Then $C$ is isomorphic to one of the
following quartics.

\begin{enumerate}
\item $y^4+az^4+xz^3+bx^2z^2+cx^4=0$ where
$a,b,c \in K$, $c \not \in K^2$.
(Notice the imperfectness of $K$ is crucial.)

\item $y^4+az^4+bx^2y^2+cx^2z^2+bx^3z+dx^4=0$
for $a,b,c,d, \in K$,
$a \not \in K^2$, $b\not =0$.

\item
\item
\item
\end{enumerate}
\end{theorem}

\noindent
Now we enumerate properties of these
families.

\begin{itemize}
\item $C$ is a purely inseparable double
cover of a quasielliptic curve.
To understand this consider the induced map
from
Frobenius map $K \to K$, $a \mapsto a^p$ and
do base change:
$$
\xymatrix{
C^{(p)}\ar[r]\ar[d]\ar@{}[dr]|-{\lrcorner}&
C\ar[d]\\
\Spec K\ar[r]^{\substack{\text{induced from}
\\ \text{Frobenius}}}&\Spec K
}
$$
This gives a map $C \to C^{(p)}$ by universal
property of pullback.

$$
\xymatrix{
C \ar[rrrd]^{\text{Frobenius}}
\ar@{-->}[rrd]\ar[rrdd] & & \\
& &  C^{(p)}\ar[d] \ar[r] & C \ar[d] \\
& &  \Spec K\ar[r] & \Spec K
}
$$

Then consider the normalization of
$C^{(p)}=C_1$ which
turns out to be quasielliptic; by universal
property of normalization we can lift:
$$
\xymatrix{
& \widehat{C^{(p)}}\ar[d]\\
C\ar@{.>}[ur]\ar[r]&C^{(p)}.
}
$$
Here the lower arrow is purely inseparable of
degree 2.

\item There exists a unique non smooth point
$x \in C$.

\item Now we iterate the construction of the
first item:
$$
\xymatrix{
\underbrace{C}_{g=3}\ar[r]&\underbrace{C_1}
_{g=1}\ar[r]&
\underbrace{C_2}_{g_2-9}\ar[r]&
\underbrace{C_3}_{g_3=0}\ar[r]&
\underbrace{C_4}_{g_4=0}\ar[r]&\cdots
}
$$
Let $x$ be the non-smooth (non-rational)
point of $C$. Map it under
these maps to $x_1$( non-smooth,rational),
then  $x_2$ (smooth, rational?),
then $x_3 $ (smooth, rational.), then $x_4$
(smooth, rational), etc.

We have
$$
\begin{matrix}
&\text{$x_2$ rational}& \text{$x$ canoical
divisor} &
E=K(c_2)\\
(i)& Y & Y & Y\\
(ii) & N & Y & N\\
(iii) & Y & N & N\\
(iv) & N & N & Y\\
(v) & N & N & N
\end{matrix}
$$
Where the last column is explained next.

\item Canonical field of $C$ = subfield of
$K(C)$ generated
by the quotients of all non-zero holomorphic
differentials
=  $K(C)$ ($C$ non-hyperelliptic).

\item Pseudocanonical field of $C$ = …
non-zero \underline{exact}
holomorphic differential := $E$. So one has
$[K(c),E]=4=p^2$.

$$
\xymatrix{
&K(C)\ar[ld]_{4}\ar[rd]^{4}\\
E&&K(C_2).
}
$$
\end{itemize}

\noindent
In conclusion, we manage to characterize
these families.

\begin{example}
\label{example-chain-of-maps}
In case 1, let $a=b=0$. We get
\begin{equation}
\label{equation-curve-pencil}
C: y^4+xz^3+\underbrace{c}_{\in K\setminus
K^2}x^4=0
\end{equation}
$$
\underbrace{C_1}_{\text{quasiell.}}:xy^2+z^3+
cx^3=0.
$$
Further
$$
\xymatrix{
C\ar[r]&C_1\ar[r]&C_2 \cong\mathbb{P}^1\ar[r]
&C_3 \cong \mathbb{P}^1\ar[r]&C_4\cong
\mathbb{P}^1\ar[r]&\cdots.
}
$$
\end{example}

\medskip\noindent
With $C$ one can construct a pencil
(fibration over $\mathbb{P}^1$)
of quartics.
Let $k=\overline{k}$ (new base field from
what we fixed at the beginning!).
Idea: replace $c$ for $t$ Equation
\ref{equation-curve-pencil}.
 $K=k(\mathbb{P}^1)$,
 $S=V(t_0(y^4+xz^3)+t_1x^4)
\subseteq \mathbb{P}^2_{(x:y:z)}\times
\mathbb{P}_{(t_0:t_1)}$.
This has a singular point, so we blow up:
$$
\xymatrix{
\tilde{S}\ar[d]\ar@/^1em/[dd]^f\\
S\ar[d]\\\mathbb{P}^1
}
$$
The generic fiber is the curve
$y^4+xz^3+tx^4=0$,
so that $t_1/t_0 \in k(\mathbb{P}^1)=K$.
Fibers are plane rational quartics,
and the singular fiber $f_1(0:1)$ is a
configuration of lines
given by a Dynkin diagram
$$
\xymatrix{
E_1 \ar@{-}[r]& E_2\ar@{-}[r]
& E_3 \ar@{-}[r]& E_4\ar@{-}[d]\ar@{-}[r]
& \cdots \ar@{-}[r]& E_{15}\\
& & & E
}
$$
We do the same with $C_1$: define
$S^1=V(t_0(xy^2+z^3)+tx^3) \subseteq
\mathbb{P}^2 \times \mathbb{P}^1$.
Then $S_1$ is also singular, and we blow up
to obtain a smooth
quasielliptic fibration ($\tilde{S^1}$ is
quasielliptic).
The fibers are plane cuspidal cubics.
The singular fiber an arrangement of curves
$$
\xymatrix{
F_1 \ar@{-}[r]& F_2\ar@{-}[r]&
F_3\ar@{-}[d]\ar@{-}[r]& \cdots \ar@{-}[r]&
F_8\\
& & F
}
$$

In fact, there is a correspondence between
the
singular fibers of the two fibrations.



\section{Lagrangian blow-up and blow-down for
4-dimensional symplectic manifolds
}
\label{section-blow-up-down}

\noindent
Misha Verbitsky, IMPA.
Geoemtric Structures Seminar, IMPA.
October 23, 2025.

\medskip\noindent
{\bf Abstract.} The usual (complex geometric)
blow-up has a symplectic
version, called the symplectic cut. I will
introduce a variant of this
construction, which is valid only in
symplectic category, called ``Lagrangian
blow-up'', and its inverse, called Lagrangian
blow-down. Lagrangian blow-down
takes a Lagrangian sphere in a 4-dimensional
symplectic manifold and contracts
it to a symplectic orbifold with a double
point. Given a symplectic orbifold
with a double point, Lagrangian blow-up
produces a symplectic manifold, and this
construction is inverse to the Lagrangian
blow-down. I will explain why these
constructions are functorial, that is,
defined on the corresponding symplectic
Teichmuller spaces. This is used to prove an
orbifold counterexample to a famous
conjecture, sometimes attributed to
Donaldson, who asked whether any symplectic
form on a K3 surface is compatible with a
Kahler structure. I will prove that
this is false for orbifolds: some K3
orbifolds admit symplectic forms not
compatible with any Kahler structure. The
same argument is used to produce a
countable family of Lagrangian spheres in a
K3 surface which are not Lagrangian
isotopic to special Lagrangian spheres. This
is a joint work with Michael Entov.



\medskip\noindent


Consider the total space of the cotangent
bundle of $\mathbb{C}P^1$.
It's $\mathcal{O}(-2)$.
Let $\gamma \in H^{0}(\mathcal{O}(2))$.
$\gamma$ defines a fiberwise linear function
on $T^*\mathbb{C}P$.

The blow-up of $\mathbb{C}^2/\pm 1$, which is
the orbifold $\mathbb{C}^2$ with
 a doble point, is the cotangent bundle $T^*
 \mathbb{C}P^1$.
That is,
$T^*\mathbb{C}P^1 \to \mathbb{C}^2/\pm 1$.
This is called {\it crepant resolution}.
Consider the $dx \wedge dy$, which is
holomorphically symplectic
on $\mathbb{C}^2$ and descends to the
quotient $\mathbb{C}^2/\pm 1$
as $\Omega_0$.

\begin{proposition}
\label{proposition-}
$\pi:T^*\mathbb{C}P^1 \to \mathbb{C}^2/\pm
1$.
Then $\pi^*\Omega_0$ is holomorphically
symplectic.
\end{proposition}

\begin{proof}
Actually, it gives the standard symplectic
form on $T^*\mathbb{C}^2$.
\end{proof}

\medskip\noindent
Now we shall define the Langrangian blow-up.
It's a procedure that starts with a
symplectic manifold
with double points and outputs a symplectic
manifold
without double points.

\begin{definition}[Lagrangian blow up]
\label{definition-lagrangian-blow-up}
Let $(B,\omega)$ be the standard ball in
$\mathbb{R}^n$,
$dz_1 \wedge dz_2 +dz_3 \wedge dz_4$.
Consider the symplectic orbifold
$(B,\omega)/\pm 1$.
Consider $M$ be a symplectic orbifold
with double point singularities.
Darboux coordinates in a neighbourhood of $0$
with $(B,\omega)/\pm 1$ resolve
singularities.
\end{definition}

This definition has the caveat that there ar
e lots of choices involved.
We shall later that these choices do not
alter the result,
namely Theorem
\ref{theorem-lagrangian-blow-up-well-defined}
.

To define the Lagrangian blow-down we first
recall

\begin{theorem}[Weinstein's Lagrangian
neighbourhood]
\label{theorem-weins-lagra-neigh}
Let $X \subset M$ be a compact Lagrangian
submanifold in $(M,\omega)$.
Then there exists a neighbourhood $U$ of $X
\subset M$
which is symplectomorphic to a neighbourhood
of $X$
in $X \subset T^*M$.
\end{theorem}

\noindent
Let $S \subset M$ be a Lagrangian sphere
(diffeomorphic to sphere
and Lagrangian).

\begin{definition}
\label{definition-lagrangian-blow-down}
The {\it Lagrangian blow-down} is obtained
from $M$
by removing $W \subset M$ and gluing $W_0$ in
its place,
where $W_0$ is such that

$$
\xymatrix{
W
\ar@{^{(}->}[r]\ar[d]&T^*\mathbb{C}P^1\ar[d]
\\
W_0\ar@{^{(}->}[r]&\mathbb{C}^2/\pm 1.
}
$$
\end{definition}

\medskip\noindent
Now we recall Teichmüller spaces.

Let $\Gamma(\Lambda^2M)$ be the space of all
2-forms on a manifold $M$.
The space of symplectic forms is a subspace
of the closed 2-forms on $M$,
which in turn is a subspace of
$\Gamma(\Lambda^2M)$.
Equip this space with the $C^\infty$
topology.
The {\it Teichmüller space} is
$\text{Symp}/\text{Diff}_0$.

Two symplectic structures are called {\it
isotopic} if they lie
in the same orbit of $\text{Diff}_0$.
The {\it period map}
$\text{Per}:\text{Teich}_s \to
H^{2}(M,\mathbb{R})$
maps a symplectic structure to its cohomology
class.

\begin{theorem}[Moser, 1965]
\label{theorem-period-map-is-local-diffeo}
The period map is a local diffeomorphism.
\end{theorem}

\noindent
This implies that $\text{Teich}_s$ is smooth.

\begin{theorem}[Moser's trick]
\label{theorem-mosers-trick}
Let $\omega_t$, $t \in S$ be a smooth family
of symplectic
structures, parametrized by a connected
manifold $S$.
Assume that the cohomology class
$[\omega_t]$ is constant.
Then all $\omega_t$ are diffeomorphic.
\end{theorem}

\noindent
This means that the fibers of $\text{Symp}\to
H^{2}(M)$
are the orbits (of the action of
$\text{Diff}_0$, I suppose).

\medskip\noindent
Consider a Lagrangian submanifold $L \subset
(M,\omega)$.
Let us assume for simplicity (though the
results
can also be obtained without this) that
$b_1(L)=0$,
i.e. first cohomology group is zero.

Let $\text{Symp}_L$ be the set of symplectic
forms vanishing on $L$,
and $\text{Diff}_0^L$ the connected component
of diffeomorphisms
preserving $L$.
Define
$\text{Teich}_L=\text{Symp}_L/\text{Diff}
_0^L$.

\begin{theorem}
\label{theorem-}
Let $V \subset H^{2}(M,\mathbb{R})$ be the
space of all
cohomology classes on $M$ which vanish on
$L$,
and $\text{Per}_L:
\text{Symp}_L/\text{Diff}^L_0 \to V$
take $(M,\omega)$ to the cohomology class of
$\omega$.
Then Per is a local diffeomorphism.
\end{theorem}

\begin{proof}
We used a version of Moser's trick: for
$\omega_t$ a family
of symplectic forms with $\omega_t|_{L}=0$,
with $[\omega_t]$ constant, then $\omega_t$
are $\text{Diff}_0^L$-isotropic.
\end{proof}

\noindent
As a corollary we obtain that this
Teichmüller space
classifies Lagrangian submanifolds:

\begin{lemma}
\label{lemma-techm-class-lagra-subma}
Assuma that $L$ and $L'$ are Lagrangian
submanifolds in $(M,\omega)$,
and $\varphi \in \text{Diff}_0(M)$ a smooth
isotopy such that
$\varphi(L)=L'$. Then $(\omega,L)$ and
$(\varphi^*\omega,L)$
represent the same point in $\text{Teich}_L$
if and only if
$L$ and $L'$ are Lagrangian isotopic in
$(M,\omega)$.
\end{lemma}

\noindent
Again: points of this Teichmüller space are
isotopy classes
of Lagrangian submanifolds.

\medskip\noindent
\begin{theorem}
\label{theorem-blow-down-diffeomorphism}
Let $L$ be a Lagrangian 2-sphere in a K3
surface  $M$
and $\hat{M}$ the Lagrangian blow-down of
$M$.
The blow-up/blow-down
is a diffeomorphism $\text{Teich}_L(M) \to
\text{Teich}(\hat{M})$.
\end{theorem}

\medskip\noindent
Now we prove the theorem which asserts that
the Lagrangian
blow-up is independent of choices.

\begin{theorem}
\label{theorem-lagrangian-blow-up-well-defined}
Let $(\hat{M},\hat{\omega})$ be a orbifold.
Then its blow-up
$(M,L)$ defines a point in $\text{Teich}_L$
independent of the choices.
\end{theorem}

\begin{proof}
By taking local Darboux coordinates, and then
Moser's trick.
\end{proof}

\noindent
Actually, the same happens for blow-downs

\begin{theorem}
\label{theorem-same-for-blow-down}
Let $(M,L,\omega)$ be a 4-manifold with a
Lagrangian sphere $L$, and
$(\hat{M},\hat{\omega)}$ be its Lagrangian
blow-down.

Then the corresponding point in
$\text{Teich}_{\hat{M}}$ is
independent from the choices made.
\end{theorem}

\begin{proof}
Similar.
\end{proof}

\medskip\noindent
A conjecture by Donaldson ways that all
symplectic structures on a K3
surface are compatible with a Kähler
structure. This conjecture
(although it's commonly believed to be true)
is still open.
When trying to prove that this would hold for
orbifolds,
Misha and collaborators realised it's
actually false.

Note: Seiberg-Witten invariants may be used
to tell whether a symplectic
form is compatible with a complex structure.

\begin{theorem}
\label{theorem-conje-false-orbif}
There exists an orbifold symplectic K3
surface $M$ with a single
double point not admitting an orbifold Kähler
structures.
\end{theorem}

\begin{proof}
\noindent
Step 1. Consider
$$
\text{Teich}_{\text{Kähler-symplectic}}
(\hat{M})
=\{\eta \in H^{2}(\hat{M}):\eta^2>0\}.
$$
As follows from [Amerik-V., 2005], the
Teichmüller space of symplectic structures
compatible
with a hyperkähler structure is Hausdorff and
connected
(the same argument works for Kähler K3
orbifolds).

Essentially, if there wasn't any orbifold
as required, the Teichmüller space
classifying
Lagrangian submanifolds would be connected,
which
is impossible by [Seidel 2000].
\end{proof}

\noindent
In fact, Seidel's result is essentially
mirror symmetry.

\medskip\noindent
Now we discuss special Lagrangian
submanifolds.

\begin{definition}
\label{definition-phase}
Let $(M,I,\omega,\Omega)$ be a Calabi-Yau
manifold, where
$\Omega \in \Lambda^{n,0}_I(M)$ is
nondegenerate and $\omega$ is Kähler.
Let $L$ be a Lagrangian submanifold.
Define the {\it phase} as
$$
\begin{aligned}
\text{phase}: L &\longrightarrow \Lambda^nL
\otimes \mathbb{C} \\
x &\longmapsto \Omega|_{T_xL}.
\end{aligned}
$$
\end{definition}

\noindent
A better definition is:
$$
\begin{aligned}
\text{phase}: L &\longrightarrow
S^1=\text{U}(1) \\
x &\longmapsto \frac{\Omega|_{T_xL}}{|\Omega
T_xL|}.
\end{aligned}
$$

\begin{definition}
\label{definition-special-lagrangian}
A special Lagrangian submanifold is the phase
is constant.
\end{definition}

\noindent
Special Lagrangian submanifolds have
deformation space of
dimension 1, it's unobstructed. They are
calibrated, so minimal.

\begin{theorem}
\label{theorem-lagrangians-of-k3}
$L \subset$ K3 is Lagrangian is isotopic to a
special Lagrangian
if and only if its blow-down is Kähler type.
\end{theorem}

\medskip\noindent
One of the basic results about special
Lagrangian
submanifolds is

\begin{lemma}[Hitchin]
\label{lemma-hitchin}
If $(M,\Omega)$ is holomorphically symplectic
and $X \subset M$
is Lagrangian with respect to
$\text{Re}\Omega$ and $\text{Im}\Omega$
then $X$ is complex.
\end{lemma}

\noindent
As corollaries:

\begin{lemma}
\label{lemma-corollary-hitchin}
Let $(M,I,J,K,g)$ be a hyperkähler manifold
of real dimension $4n$,
considered as a Kähler manifold
$(M,I,\omega_I)$ and $\phi:\Omega^n \in
\Lambda^{2n,0}(M,I)$ the corresponding
holomorphic volume form.
Consider a Lagrangian submaniold $S \subset
(M,\omega_I)$
such that $a \omega_J + b\omega_K|_S=0$ for
some nonzero real numbers
$a,b \in \mathbb{R}$ such that $a^2+b^2=1$.
Then $S$ is special Lagrangian, and,
moreover,
it is holomorphic with respect to the complex
structure
$-a K+b J$.
\end{lemma}

\begin{lemma}
\label{lemma-corollary2-hitchin}
Let $\omega_I,\omega_J$ and $\omega_K$ on
$M$, with $\dim M=4$.
Let $S \subset (M,\omega_I)$ be Lagrangian.
Then $S$ is special Lagrangian if and only if
$S$ is complex analytic
on $a J + b K$ for two nonzero numbers  $a,b
\in \mathbb{R}$ such that
$a^2+b^2=1$.
\end{lemma}

\noindent
Now an auxiliary claim:

\begin{lemma}
\label{lemma-k3-type}
$\Omega$ is holomorphic symplectic
\end{lemma}

\begin{theorem}
\label{theorem-}
Let $S\subset M$ be a Lagrangian sphere in a
K3 surface.
Then  $S$ is isotropic to a special
Lagrangian if and only if
$M$ is of Kähler type.
\end{theorem}


\section{Systolic inequality for tori}
\label{section-systolic-inequality}

\noindent
Dimitri Korshunov, IMPA.
Minimal Surfaces end of course presentation.
October 24, 2025.

\medskip\noindent
This is due to L. Guth.

\begin{definition}
\label{definition-systole}
Given a Riemannian manifold,
a {\it systole} $\text{Sys}(M,g)$ is the
infimum of the lengths
of noncontractible paths.
\end{definition}

\begin{theorem}[Gromov]
\label{theorem-gromov}
For any metric on a torus $T^n$,
$$
\text{Sys}(T^n,g) \leq C_n
\text{Vol}(T^n,g)^{\frac{1}{n}},
$$
where $C_n$ is a constant depending only on
the dimension
of the torus $n$.
\end{theorem}

\noindent
This also holds for $\mathbb{R}P^n$. We will
use induction,
for which we need the following stronger
statement:
given a number bounded by half the systole,
there exists a ball
which is {\it big}:

\begin{theorem}[Guth]
\label{theorem-guth}
For $(T^n,g)$, given
$R<\frac{1}{2}\text{Sys}(T^n,g)$,
then there exists a point $p$ on $T^n$
such that the volume of ball of radius $r$
with center in $p$
satisfies
$$
\text{Vol}B(p,R) \geq C_nR^n.
$$
\end{theorem}

\noindent
It was conjectured by Gromov that $C_n=1$.

\medskip\noindent
Before stating the theorem that we will
prove, we introduce the
following definition that will allow us to
measure cohomological cycles:

\begin{definition}
\label{definition-measu-cohom-cycle}
Let $\alpha \in H^{1}(M_g,\mathbb{Z}_2)$.
$$
L(\alpha)
=\text{inf}\{|\lambda|:\lambda\text{
1-cycles, }\alpha(\lambda)\neq 0\}.
$$
\end{definition}

\noindent
The theorem that we will prove, which implies
Theorem \ref{theorem-guth}
is the following

\begin{theorem}
\label{theorem-bound-classes}
There exists a constant $\varepsilon_n>0$
such that given a manifold $(M,g)$, and
$\alpha_1,\ldots,\alpha_n \in
H^1(M,\mathbb{Z}_2)$,
\begin{enumerate}
\item $L(\alpha_i)>2R$
\label{item-positive}
\item
\label{item-cup} $\alpha_1 \cup  \alpha_2
\cup \ldots \cup \alpha_n \neq 0 \in
H^n(M,\mathbb{Z})$.
\end{enumerate}

\noindent
Then there exists $p \in M$ such that
$$
\text{Vol}B(p,R) \geq \varepsilon_nR^n.
$$
\end{theorem}

\noindent
In particular $T^n$ and $\mathbb{R}P^n$
satisfy (\ref{item-cup}).

Before we proceed to the proof we need the
following definition.
Essentially we ask that the surface minimizes
volume in its homology
class up to some constant $\delta$.

\begin{definition}
\label{definition-delta-minimal-surface}
A  {\it $\delta$-minimal surface} is an
embedded hypersurface $Z$
such that for any $Z'$ such that $[Z'-Z] = 0$
in  $H_{n-1}(M,\mathbb{Z}_2)$,
then  $\text{Vol}Z < \text{Vol}Z'+\delta$.
\end{definition}

\noindent
Moreover, we need the following

\begin{lemma}[Stability]
\label{lemma-stability}
Let $(M,g)$ de closed Riemannian manifold and
$\alpha \in H^1(M,\mathbb{Z}_2)$.
Suppose $Z$ is Poincaré dual to $\alpha$ and
that $Z$ is $\delta$-minimal.
If $R<\frac{1}{2}L(\alpha$) then the
following
holds for any $p \in M$:
$$
|Z \cap \text{Vol}B(p,\frac{R}{2}|\leq
\frac{2}{R}|\text{Vol}B(p,R)|+\delta.
$$
\end{lemma}

\begin{proof}[Proof of the stability lemma]
\medskip\noindent
Step 1. Using commas formula, namely
$$
\text{Vol}(B_R)=\int_0^R S_tdt,
$$
we see that there exists $t$ such that
$\frac{R}{2}<t<R$ and
$$
|S(p,t)|_R< \frac{2}{R}|\text{Vol}B(p,R)|.
$$

\medskip\noindent
Step 2. Consider
$[Z \cap B(p,t)] \in
H^{n-1}(B(p,t),S(p,t),\mathbb{Z})$.
This can be interpreted geometrically since
we are using relative homology.

\medskip\noindent
Step 3. We claim that $[Z \cap B(p,t)]=0$.
By contradiction,
if it weren't trivial, then there would exist
a cycle $H_1(B,\mathbb{Z})$
such that $Z' \cap \gamma \neq 0$ since
Poincaré pairing is nondegenerate ---
namely any cocycle that is not zero pairs
nontrivially with some cycle.
Then we will just split the cocycle in very
small parts,
and we will reach a contradiction with the
fact that
$R<\frac{1}{2}L(\alpha)$.

We will state this as another lemma.
Essentially it says that any
cycle in homology can be decomposed as a sum
of cycles
with small lengths.

\begin{lemma}[Curve factoring, Gromov]
\label{lemma-curve-factoring-gromov}
Let $(M,g)$ be a Riemannian manifold and
$\gamma$ a cycle in $B(p,t)$.
Then $[\gamma]=\sum \sigma_i$ is such that
 $L(\sigma_i)<2t+a$ for all $i$.
\end{lemma}

\begin{proof}
We spit a cycle (a curve) into some quantity
of intervals
such that each has length smaller than $a$,
then add paths joining the center point of
length $t$,
so we count $t$ twice: to reach the point and
to get away from it.
\end{proof}

\noindent
We apply this to $\gamma$ and reach a
contradiction
with the hypothesis that
$R<\frac{1}{2}L(\alpha)$ making $a$ smaller
and
smaller.

\medskip\noindent
Step 4. Finally we use the minimality of $Z$
and the property
derived from Commas lemma to arrive at a
contradiction.
\end{proof}

\noindent
Now we prove Theorem
\ref{theorem-bound-classes}:

\begin{proof}
By induction. The case $n=1$ is trivial.
Now take the last one of our classes,
$\alpha_n$.
Let $Z$ be Poincaré dual of $\alpha_n$ and
$\delta$-minimal.
Notice that $\alpha_1\cup \ldots \cup
\alpha_{n-1}\neq 0$
since if it were, then $\alpha_1\cup
\ldots\cup \alpha_n$ would also vanish.

By induction hypothesis we get a point $p \in
Z$
such that $\text{Vol} B_Z(p,\frac{R}{2}) >
R^{n-1}$.
This means that , $|Z \cap
B(p,\frac{R}{2})|>\varepsilon_n(R/2)^{n-1}$.

By the Stability Lemma \ref{lemma-stability},
$$
\begin{aligned}
\frac{2}{R}|B(p,R)|&>\varepsilon_n\frac{R^{n-
1}}{2^{n-1}}-\delta\\
\implies
|B(p,B)|&>\frac{\varepsilon_n}{2^n}R^n-
\frac{\delta
R}{2}
> \frac{\varepsilon_n}{2^n}R^n.
\end{aligned}
$$
\end{proof}

\section{From higher Lie groupoids to higher
Lie algebroids}
\label{section-highe-group-highe-algeb}

\noindent
Matias del Hoyo, UFRJ.
Symplectic Geometry Seminar, IMPA.
October 29, 2025.

\medskip\noindent
{\bf Abstract.}
This is a report of an ongoing collaboration
with A. Cabrera, where we develop a
higher Lie theory. Starting with simple,
concrete examples, we present the
differentiation functor in a geometric,
explicit way, using simplicial methods
and differential graded algebras. We describe
how our construction extends the
classical Lie functors, relate our project
with previous works, and show a van
Est theorem linking the cohomology at the
global and infinitesimal levels.

\medskip\noindent

Upshot: Integration-differentiation relation
between Lie groups and Lie algebras
in higher versions. Groupoids are the
counterparts of Poisson manifolds, and in
the next level there are 2-groupoids and
Courant algebroids. How to integrate a
2-groupoid to a Courant algebroid? This is an
open question which motivated this
work. See arXiv
\href{http://arxiv.org/abs/2309.14105}%
{2309.14105}.

\medskip\noindent
Let $M$ be a smooth manifold.

\begin{remark}
\label{remark-smoot-funct-ring-fully-faith}
$M \mapsto C^\infty M$ is fully faithful.
\end{remark}

\noindent
Which allows us to think of manifolds as
algebraic objects.
Notice that
$\Omega^kM=\{k\text{-forms}\}
=\Lambda^k_{C^\infty M}\text{Der}(C^\infty
M,C^\infty M)$,
which is (algebraically) locally spanned
by the exact 1-forms.

Recall that the exterior differential is
the only operator $d:\Omega^kM \to
\Omega^{k+1}M$
such that $d(\alpha \wedge \beta)=d\alpha
\wedge \beta +
(-1)^{|\alpha|}\alpha \wedge d\beta$
and that $df(X)=Xf$.

This makes $(\Omega M,\wedge,d)$ a
differential graded algebra.
Recall that $\text{dom}(X,Y)=X\omega(Y)-Y
\omega(X)-\omega([X,Y])$.
We also have de Rham cohomology.

Consider a submanifold $S \subset M$. Recall
that $\mathcal{I}_S/\mathcal{I}^2_S$ is
essentially
the sections of the normal bundle.
In particular, for $M \subset M \times M$
we can see that the embedding vanishes on the
diagonal,
so there is no linear part in the Taylor
expansion,
and the second term is in fact the de Rham
differential
$df=\pi_1^*f-\pi_2^*f \text{ mod
}\mathcal{I}^2_{\Delta(M)}$.

In fact, this construction can be extended to
a simplicial manifold
(which I think is just a presheaf valued on
category of manifolds)
using projections and diagonals,
$$
\xymatrix{
\cdots & M\times M \times M
\ar @< 3pt> [r]
\ar @< 0pt> [r]
\ar @<-3pt> [r] & M \times M \ar @< 2pt> [r]
\ar @<-2pt> [r]& M.
}
$$
To this we can apply the $C^\infty$ functor
and obtain
a cosimplicial ring (not a dga?).
The idea here
is that de Rham differential is an
alternating sum of pullbacks;
as other differentials will also be.

\medskip\noindent
Now observe that for a finite-dimensional
vector space $V$,
$$
\left\{\substack{\text{Lie bracket} \\
V \times V
\xrightarrow{[\cdot,\cdot]}V}\right\}
\leftrightarrow
\left\{\substack{\text{differential on} \\
\Lambda V^*}\right\},
$$
making $\Lambda V^*$ into a differential
graded algebra.

Recall that given a Lie algebra
$\mathfrak{g}$,
we have $(\Lambda^\bullet \mathfrak{g}^*
,d)=\text{CE}(\mathfrak{g})$,
the Chevalley Eilenberg algebra,
which gives Lie algebra cohomology.

\medskip\noindent
[Missing part on Simplicial approach,
exponential map,
Simplicial manifold nerve]

\medskip\noindent
Lie algebroid: $\Delta \to M$ vector bundle
with $[\cdot,\cdot]$ on $\Gamma A$,
$A \xrightarrow{\rho}TM$ Leibniz.

\begin{remark}
\label{remark-lie-algebroids-differentials}
$$
\left\{\substack{\text{Lie algebroids on} \\
E \to M}\right\}
\leftrightarrow
\left\{\substack{\text{differentials} \\
\text{on }\Lambda\Gamma E^*}\right\}.
$$
\end{remark}

\noindent
This gives a differential graded algebra
$$
\xymatrix{
\Delta\ar[r]&M\ar[r]&\text{CE}(A)
}
$$
and
$$
\xymatrix{
C^\infty M\ar[r]^{\tilde{\rho}}&\Gamma
E^*\ar[r]^{\widetilde{[\cdot,\cdot]}}
&\Gamma E^* \wedge \Gamma E^*\ar[r]&\cdots.
}
$$

From Lie groupoids to Lie algebroids:
Lie functor/differentiation.

$$
\Omega(G)=\Gamma \Lambda^\bullet
\underbrace{\supset}
_{\text{sub dga}}
\Gamma \Lambda^\bullet T^s
\xrightarrow{-/G}\text{CE}(A_G).
$$
There's the following ecosystem around Lie
groupoids:
$$
\xymatrix{
&\text{Lie groupoids}
\ar[dl]\ar@{->>}[d]\ar[dr]^{\text{nerve}}\\
\text{Lie algebroids}&
\substack{\text{Differentiable stacks} \\
\supset \text{orbifolds}}
&\text{simplicial manifolds.}
}
$$

\medskip\noindent
Lie groupoids are models for stacks,
which are geometric objects solving
classification
problems with objects identified up to
isomorphism.
Relaxing the identification to a weaker
notion
of equivalence we must introduce higher
versions of stacks.

Moving higher,
$$
\xymatrix{
\text{Lie groupoid}\ar[r]\ar[d]_{\text{Lie}}
&\text{Simplicial manifolds}\ar[d]^?\\
\text{Lie
algebroids}\ar[r]&\text{dg-manifolds}
}
$$
dg algebra is a  dga $(A=\bigoplus_{k=1}^n
A_k \to M,\wedge,d)$
starting at $C^\infty M$, and continuing
until you get to
the sections, namely
$$
\xymatrix{
C^\infty M \ar[r]_d & A_2 \ar[r]& \cdots &
\cong \Gamma A.
}
$$
[dH-Cabrera] provide the solution to this
problem,
namely what to put in the right vertical
arrow:

\begin{theorem}[dH-Cabrera]
\label{theorem-dh-cabrera}
$$
\frac{(C^\infty_N(G),\cup,\sum(-1)A^*_i)}{J+
\delta
J}
=\text{CE}(A_G)
$$
where $A_G$ is a higher Lie algebroid
and $J$ is defined by: $f$ is in $J$ if it
has higher vanishing.
\end{theorem}


Groupoids are the counterparts of Poisson
manifolds,
and in the next level there are 2-groupoids
and Courant algebroids.
How to integrate a 2-groupoid to a Courant
algebroids?
This is an open question which motivated this
work.



\section{Formality of Dolbeault cohomology}
\label{section-forma-dolbe-cohom}

\noindent
Misha Verbitsky, IMPA.
Geometric Structures Seminar, IMPA.
October 29, 2025.

\medskip\noindent
{\bf Abstract.}
A quasi-isomorphism of differential graded
algebras (DGA) is a multiplicative
map inducing an isomorphism on cohomology. A
DGA is called  formal if it can be
connected by a chain of  quasi-isomorphisms
to its cohomology algebra. For the
de Rham algebra of a compact manifold,
formality is an important  topological
property which has many implications (such as
vanishing of Massey products). For
the Dolbeault DGA of a compact complex
manifold, the situation is not well
understood, unless it is Kahler. The
Dolbeault DGA of a torus, symmetric spaces
and hyperkahler manifolds are known to be
formal.  We prove that the Dolbeault
DGA of a complex nilmanifold is formal only
if it is a torus, and the Dolbeault
algebra of (0,p)-forms is formal if and only
if the complex structure is
abelian. This is a joint work with Tommaso
Sferruzza.

\medskip\noindent
Upshot: the Dolbeault dga of a nilmanifold
$M$ is never formal.

\medskip\noindent
If a space is formal, Massey products vanish.
Unfortunately, Massey products are not
well-defined.
Let $a,b,c \in \Lambda^*(M)$ be closed forms
such that their cohomology classes satisfy
$[a][b]=[b][c]=0$.
Let  $\alpha,\gamma \in \Lambda^*(M)$
be such that $d(\alpha)=a \wedge b$ and
$d(\gamma)b \wedge c$.
Then the {\it Massey product} of $a,b,c$
is the class  $[\alpha \wedge
c-(-1)^{\tilde{a}}a \wedge \gamma]$
which is only defined up to
 $\left\langle{}\text{img}L_a + \text{img}
 L_c\right\rangle{}$
where $L_a$ and $L_c$ are multiplication by
$a$ and $c$ operators.

The {\it Heisenberg nilmanifold} is the
quotient of
the Heisenberg group $G$ of upper triangular
matrices
with $1$ on the diagonal over
$G_{\mathbb{Z}}$
of the same group with integer entries.
The Heisenberg nilmanifold fibers over a
torus
with fiber a circle:
$$
\xymatrix{
0\ar[r]&S^1\ar[r]&G/G_{\mathbb{Z}}\ar[r]&
\mathbb{R}^2/\mathbb{Z}^2\ar[r]&0
}
$$
which in fact comes from the central
extension
$$
\xymatrix{
0\ar[r]&\mathbb{R}\ar[r]&G\ar[r]&\mathbb{R}
^2\ar[r]&0.
}
$$


\begin{lemma}
\label{lemma-heise-nilma-masse-produ}
Massey products on  $G/G_{\mathbb{Z}}$ are
non-zero.
\end{lemma}

\begin{proof}
Step 1. There is an action of $G$ in the
cohomology of $G$.
It's not hard to see that the cohomology
classes on $G/G_{\mathbb{Z}}$
can be represented by right $G$-invariant
forms,
and the cohomology of $G/G_{\mathbb{Z}}$ is
equal to the cohomology of the complex
of right-$G$-invariant forms on $G$.

\medskip\noindent
Step 2. In fact, this complex is the same as
the
Chevalley-Eilenberg complex of the Lie
algebra $\mathfrak{g}$ of $G$,
and the cohomology is the Lie algebra
cohomology.

\medskip\noindent
Step 3. From a basis $a,b,t$ of
$\mathfrak{g}$ with the only nontrivial
commutator $[a,b]=t$ and $\alpha,\beta,\tau$
dual basis in $\mathfrak{g}^*$
with the only nontrivial differential
$d\tau=\alpha\wedge\tau$,
we obtain a basis $\alpha \wedge \beta,
\alpha \wedge \tau, \beta \wedge \tau$
in $\Lambda^2(\mathfrak{g}^*)$ with
$d|_{\Lambda^2\mathfrak{g}^*}=0$.
This gives $\text{rk}
H^1(G/G_{\mathbb{Z}})=2$
and $\text{rk}H^2(G/G_{\mathbb{Z}})=2$.

\medskip\noindent
Step 4. Compute the Massey product of
$\alpha,\beta,\tau$:
it does not vanish.
\end{proof}

\medskip\noindent
Digression: this is Sullivan's motivation for
working on dgas.
Consider the class of spaces whose homotopy
and homology
groups are vectors spaces over $\mathbb{Q}$
(and it's better to also assume that $\pi_1$
is nilpotent).
There is a equivalence of categories between
old category
and homotopy category of rational homotopy
spaces.
And in fact it's also equivalent to the
(homotopy?)
category of dga algebras.

\medskip\noindent
Now we move to differential graded algebras
(dga).
A morphism of dgas (which I think only means
that commutes with
the differentials) is called {\it
quasi-isomorphism} if it
induces isomorphisms on all $H^i(A^\bullet)$.

A dga is {\it formal} if it is connected to
$H^\bullet(A^\bullet)$
by a chain of quasi-isomorphisms.

Recall the $dd^c$-lemma. It states that for
the twisted
differential $d^c=Jd J^{-1}$, if a form
$\alpha$ is $d$-closed and $d^c$-exact,
(or $\partial$-exact and
$\overline{\partial}$-exact),
then $\alpha\in \text{img}dd^c$.

\begin{lemma}
\label{lemma-quasi-isomorphisms}
The map $(\Lambda^\bullet(M)_{d^c},d)\to
(\Lambda^\bullet(M),d)$
is a quasi-isomorphism.
\end{lemma}

\begin{proof}
Uses Hodge decomposition of a form to prove
induced maps
on cohomology are injective. Then prove
surjectivity.
\end{proof}

\begin{lemma}
\label{lemma-kahler-formal}
A compact Kähler manifold is formal (i.e. the
usual cohomology algebra
is formal).
\end{lemma}

\begin{theorem}
\label{theorem-kahler-dolbeault-formal}
$M$ Kähler then
$(\Lambda^\bullet(M),\overline{\partial})$ is
formal.
\end{theorem}

\medskip\noindent
A manifold $M$ is {\it nilmanifold} if it has
an action of aa nilpotent Lie
group. All nilmanifolds are obtained as
qutient spaces $M=G/\Gamma$
for a cocompact lattice in a nilpotent Lie
group.

I gather there are two things:
one is an {\it integrable complex structure
on a Lie algebra}, which
is a subalgebra
$\mathfrak{g}^{1,0}\subset \mathfrak{g}
\otimes_\mathbb{R} \mathbb{C}$
such that $\mathfrak{g}^{1,0} \oplus
\overline{g^{1,0}}
=\mathfrak{g} \otimes \mathbb{C}$.
Integrability is
$[\mathfrak{g}^{1,0},\mathfrak{g}^{1,0}]
\subset \mathfrak{g}^{1,0}$.

And then there's a {\it complex nilmanifold},
which is $M=G/\Gamma$
with a complex structure $I$ such that $G$
has
a right-invariant complex structure and $G
\to M$ is holomorphic.
In this case we say $I$ is {\it invariant}.
Not all complex structures on a nilmanifold
are invariant:
a compact torus admits a non-Kähler complex
structure that is not invariant.
A Kodaira surface, a locally trivial
holomorphic fibration over $T$ with fiber
$T'$ with non-trivial Chern class, for $T$
and $T'$ elliptic curves,
is a nilmanifold.

Recall the Chevalley-Eilenberg cohomology
on $\Lambda^\bullet\mathfrak{g}^*$ where the
Lie bracket becomes
a differential by Cartan formula.
More generally,

\begin{proposition}
\label{proposition-cheva-eilen-diffe}
Let $w \in \Hom(\Lambda^2V,V)$.
Consider the dual map $d_w:V^* \to
\Lambda^2V^*$.
Extend this map to $d_w:\Lambda^kV^*  \to
\Lambda^{k+1}V^*$
using Leinbiz rule. Then $d^2_w=0$ if and
only if $w$
defines the Lie algebra structure on $V$.
\end{proposition}

\noindent
We may identify the Chevalley-Eilenberg
complex of the Lie algebra  $\mathfrak{g}$
with the
complex of left-invariant differential forms
in its
Lie group $G$. This defines a natural
homomorphism of dgas
$(\Lambda^\bullet\mathfrak{g}^*,d) \to
(\Lambda^\bullet(G/\Gamma),d)$
for any discrete group $\Gamma \subset G$.

\begin{theorem}[Nomizu]
\label{theorem-nomizu}
Let $M=G/\Gamma$ be a nilmanifold.
Then the map above is a quasi-isomorphism.
\end{theorem}

As a corollary, we see that nilmanifolds,
with exception of tori, are never formal.

\begin{proposition}
\label{proposition-}
Let $(\mathfrak{g},I$) be a complex structure
on a Lie algebra. Then we may consider the
Dolbeault
cohomology of
$\partial_w:\Lambda^\bullet(\mathfrak{g}^*)
\to\Lambda^2(\mathfrak{g}^*)$.
Then $[\cdot,\cdot]_{\partial_w}$ is defined
as follows: let
$x,y \in \mathfrak{g}^{1,0}$, $x',y' \in
\mathfrak{g}^{0,1}$.
Then
$$
\begin{aligned}
[x,y]_{\partial_w}&=[x,y]\\
[x,x']_{\partial_w}&=[x,x']^{1,0}\\
[x',y']_{\partial_w}&=0.
\end{aligned}
$$
\end{proposition}

\noindent
This implies that $(CE,\overline{\partial})$
is formal if and only if
$[\cdot,\cdot]=0$.

\medskip\noindent
We intend to apply the construction made
before for the usual
Lie algebra cohomology but for Dolbeault
cohomology.
Unfortunately there is no analogue of
Nomizu's theorem Theorem
\ref{theorem-nomizu}.

\noindent
[Missing slides]

\medskip\noindent
We say a dga has {\it Poincaré duality} if
the product
map $H^i(A^* ) \times H^{n-i}(A^*)\to
H^n(A^*)=\mathbb{C}$
is a non-degenerate pairing in cohomology.
A morphism of dgas $(B^*,d)\to(A^*,d)$, where
$(B^*,d)$ admits
Poincaré duality, that induces injective maps
on cohomology
is called a {\it domination}.

\begin{theorem}[Milivojevic,Stelzig,Zoller
arXiv
\href{http://arxiv.org/abs/2306.12364}%
{2306.12364}]
\label{theorem-msz}
Let $\psi:(B^*,d)\to (A^*,d)$ be a domination
of dga.
If $(B^*,d)$ is not formal, then $(A^*,d)$ is
not formal.
\end{theorem}

\begin{theorem}[Barberis-Dotti-V., 2007]
\label{theorem-bdv7}
A complex nilmanifold has trivial canonical
bundle,
trivialized by an invariant holomorphic top
form.
\end{theorem}

As corollaries we get that the natural dga
morphism
$(\Lambda^\bullet\mathfrak{g}^*,
\overline{\partial}_w)\to
(\Lambda^\bullet
G/\Gamma,\overline{\partial})$ is a
domination,
and that the Dolbeault dga of a nilmanifold
$M$ is never formal.






\section{Hyperbolic surfaces and Mirzakhani's
curve counting}
\label{section-hyper-surfa-mizak-curve}

\noindent
Viveka Erlandson, University of Bristol, UK.
8th Brazilian School of Dynamical Systems,
IMPA.
November 4, 2025.

\medskip\noindent
{\bf Abstract.}
It’s a classical problem to study the growth
of the number of periodic orbits of
bounded length $L$ under a geodesic flow on a
manifold. When it comes to
(closed) hyperbolic surfaces Huber proved
that the number is asymptotic to $e^L
L$ and this results has since been
generalized in many directions, maybe most
famously by Margulis to negatively curved
manifolds as well as more general
flows. However, a fundamentally different
result is due to Mirzakhani who
counted the subset of those closed geodesics
on hyperbolic surfaces which have
no self-intersections (or more generally,
those inside a fixed mapping class
group orbit) and showed that in this setting
the growth is asymptotic to a
polynomial in L of degree only depending on
the topology of the surface. In this
mini-course we will give the necessary
background on hyperbolic surfaces to
understand these results, give an idea of the
proof of Mirzakhani’s theorem and
introduce some tools used such as measured
laminations and train tracks. Time
allowing we present some generalizations (for
example, to non-simple closed
geodesics or to various metrics) and
applications to the distribution of closed
geodesics.

\medskip\noindent
Our goal is to count certain closed geodesics
on hyperbolic surfaces.
By ``certain'' we main, for example, simple
or fixed topological type;
we will define these properly. Also, we may
treat more general
types of curves than only hyperbolic, e.g.
combinatorial, algebraic.
As ``closed geodesics'', we mean free
homotopy classes of curves.

As historical note,
results on coarse polynomial growth for
simple curves
were are due to Birman-Serres, Rees in the
80's.
In 2001, McShane-Runn studied asymptotic
growth
for simple curves on punctured torus.
Asymptoric growth on any hyperbolic
surface was proved by Mirzakhani for simple
curves in 2008
and then in 2016 for any topological type of
curve.

Here's an outline of this minicourse:
\begin{itemize}
\item Introduction, motivation.
\item Measured laminations, train tracks,
Thurston measure.
\item  (Geodesic) currents.
\item Curve counting theorem.
\end{itemize}

\medskip\noindent
Throughout $S$ will be a topological
(connected, orientable
of finite type, i.e. finite genus). Let
$g<\infty$ be the genus
and $r< \infty$ the number of punctures.
We assume that $\chi(S)=2-2g-r<0$,
which is equivalent to saying that $S$
admits hyperbolic metrics.
We will not consider the pair of pants,
since there's only three simple closed
geodesics on it,
making it uninteresting for our means.

Let $X$ be a hyperbolic surface homeomorphic
to $S$.
Consider the universal cover
$\tilde{X}=\mathbb{H}^2$ of $X$,
and
$\pi_1(X)=\Gamma\underbrace{<}
_{\substack{\text{discrete}
\\ \text{free}}}
\text{PSL}(2,\mathbb{R})=\text{Isom}^+
(\mathbb{H}^2)$.
Recall that a $4g$-gon yields a genus $g$
surface $S_g$ by
adequate identification of its sides.

For us, a {\it curve} $\gamma$ is a free
homotopy class
of a closed immersion $\mathbb{S}^1
\hookrightarrow S$,
which are closed $X$-geodesics.
In turn, a {\it multi-curve}
is a formal sum $\sum_{i=1}^na_i\gamma_i$
for $\gamma_i$ distinct curves and $a_i>0$,
$a_i \in \mathbb{Q}$.
The {\it length} of a curve $\gamma$ is
$\ell_X(\gamma)$ is the length of
$X$-geodesic representing $\gamma$,
and the length of a multi-curve is the sum of
the lengths
of the curves in the formal sum.

We won't use this: we can think of $X$ as
an element of Teichmüller space
$\mathcal{T}(S)\cong \mathbb{R}^{6g-6+2r}$
parametrizing complex structures on $S$.

\medskip\noindent
Notice that given $X \in\mathcal{T}(S)$,
the number of curves with length bounded by a
given number $L$,
$\#\{\gamma|\ell_X(\gamma)\leq L\}$,
is finite. This can be checked going to the
universal cover.
However, as $L \to \infty$, the number also
tends to infinity.

Huber, Hekhal (59' or later) showed that if
$X$ is hyperbolic
on $S_{g,r}$, then
$\#\{\gamma|\ell_X(\gamma)\leq L\}\sim
\frac{e^L}{L}$
asymptotically, that is,
$$
\lim_{L\to \infty}\frac{\#
\{\gamma|\ell_X(\gamma)\leq L\}}{e^L/L}=1.
$$

\medskip\noindent
What is we restrict to a fixed topological
type?

Recall that the {\it mapping class group} of
$S$ is
$\text{Map}(S):=\text{Homeo}^+(S)/
\text{homotopy}$.
Then $\gamma$ and $\gamma_0$ are of the same
(topological)
type if $\gamma \in \text{Map}(S)\gamma_0$.

We define the {\it (geometric) intersection}
of $\alpha,\beta$
distinct curves by
$$
i(\alpha,\beta)=\text{min}\{\#(\alpha' \cap
\beta'|
\substack{\alpha'\text{ homotopic to }\alpha
\\
\beta'\text{ homotopic to }\beta}\}.
$$
(In fact, this minimum is always realized by
geodesics,
but it's more a topological definition.)
The {\it self-intersection} of a curve is
$$
i(\alpha,\alpha)=\frac{1}{2}\text{min}
\{\#(\alpha' \cap
\alpha''|\alpha',\alpha''\text{
homotopic to }\alpha\}.
$$

We say $\gamma$ is {\it simple} if
$i(\gamma,\gamma)=0$.

It is known that $i(\cdot,\cdot)$ is
$\text{Map}(S)$-invariant,
and that $\{\gamma:i(\gamma,\gamma)=k\}$ is a
finite union
of topological types.

\begin{theorem}[Mirzakhani, '08, '16]
\label{theorem-mirzakhani}
Let $X$ be a hyperbolic surface, $X \cong
S_{g,r}$
and $\gamma_0$ any (multi-)curve.
Then
$$
\#\{\gamma\text{ type
}\gamma_0|\ell_X(\gamma)\leq L\}
\sim CL^{6g-6+2r},
$$
where $C>0$,
$C=\frac{c(\gamma_0)}{b_\gamma}B(X)$.
\end{theorem}

\noindent
As a corollary, we obtain
$$
\begin{aligned}
\#\{\gamma\text{ simple}|\ell_X(\gamma)\leq
L\}&
\sim\text{const.}L^{6g-6+2r}\\
\#\{\gamma|i(\gamma,\gamma)\leq
k,\ell_X(\gamma)\leq L\}
&\sim (\text{const. depending on
$k$})L^{6g-6+2r}.
\end{aligned}
$$

\medskip\noindent
As an example, consider the torus with simple
(multi-)curves.
If we put a puncture it becomes a hyperbolic
curve.
Forgetting the puncture, we have an Euclidean
flat torus,
of which we know the geodesics: they are
straight lines
with rational slope in the plane (before
quotienting).
We see that the amount of geodesics of length
bounded by $L$
is $\frac{1}{2}\mathbb{Z}^2 \cap B((0,0),L)$,
which is in
fact $\sim \pi L^2$.
A result by Mc-Shane-Rivn, 2001 shows that
 $$
\{\text{(primitive) simple curves}\}\sim
\frac{\pi}{6}L^2.
$$
Where primitive (most likely) excludes
counting
a curve twice as a multicurve.


\medskip\noindent
Let us introduce a measure on $\mathbb{R}^2$.
Let $L>0$. Then
$$
m^L=\frac{1}{L^2}\sum_{p \in
\mathbb{Z}^2}\delta_{\frac{1}{L}p}
\xrightarrow{L \to
\infty}\text{Leb}_{\mathbb{R}^2}.
$$
For $p=(a,b)$,
$\frac{1}{L}p=(\frac{1}{L}a,\frac{1}{L}b)$.
$$
\begin{aligned}
U_{\mathbb{E}}&=\{x \in \mathbb{R}^2:|x|\leq
1\}\\
m^L(U_{\mathbb{E})}&=\frac{1}{L^2}\#\{p \in
\mathbb{Z}^2
\left|\frac{1}{L}p\right|\leq 1\}\\
&=\frac{1}{L^2}\#\{p \in \mathbb{Z}^2:|p|\leq
L\}\\
&=\frac{1}{L^2}\#\mathbb{Z}^2 \cap
B_{\mathbb{R}^2}((0,0),L)\\
&\to\text{Leb}_{\mathbb{R}^2}(U_{\mathbb{E})}
=\pi.
\end{aligned}
$$

\medskip\noindent
Let's go back to study the geodesics on a
torus.
First we recall a fact, which is that
on any hyperbolic surface $S$ there exists a
compact $K \subset S$ such that all
simple closed geodesics stay in $K$.

Now consider a geodesic given by a curve of
irrational slope,
which is dense in the torus. However,
it is not dense on a hyperbolic torus
$S_{1,1}$.
We may approximate this geodesic by $q_n \in
\mathbb{Q}$ with
$q_n \to \alpha$. Each $q_n \rightsquigarrow
\gamma_n$
is a closed geodesic on $S_{1,1}$.
The lamination in this case is the Hausdorff
limit
of the $\gamma_n$.

\begin{definition}
\label{definition-lamination}
A {\it lamination} $L$ on a (hyperbolic)
surface $S$
is a compact subset of $S$ consisting of
simple, pairwise
disjoint complete geodesics.

A {\it measured lamination} $\lambda$ is
$\lambda=(L,\lambda)$ is a lamination $L$
together with a {\it transverse measure}
$\lambda$:
a measure is assigned to each arc $I$
transverse to $L$.

We write $ML(S)$ for the space of all
measured laminations on $S$.
\end{definition}

\noindent
While the definition of lamination depends on
a choice
of hyperbolic metric on the surface, the
different $ML(S)$
that we obtain with different metrics are
actually homeomorphic,
and we thus do not distinguish between them.

As an example, every simple multi-curve is a
measured lamination.

Morally we think of measured laminations as
limits of simple multicurves:
simple multi-curves are dense in $ML(S)$.

\medskip\noindent
Now we will discuss train tracks, which will
give us a combinatorial model for measured
laminations.

\begin{definition}
\label{definition-train-track}
A {\it train track} $\tau$ on $S$ is an
embedded (trivalent)
graph such that at every vertex the three
half edges
meet with the same tangent line and such that
the complementary region (cc. $S \setminus
\tau$)
are not disk
(with a certain caveat: it can be a
contractible
region as long as the vertices cannot
preserve
it's smooth incidence structure
along the homotopy),
a cyllinder (ring) nor a punctured disk.

A {\it set of weights} on $\tau$ is an
assignment
of a non-negative real number $w_e$
to each edge $e$. They satisfy the {\it
switch equations}
if at every vertex we have $w_a=w_b+w_c$
(where $b$ and $c$ ``stem'' from $a$).

We denote $w(\tau)$ the solutions to the
switch equations
and $w_{\mathbb{Z}}(\tau)$ the integral
solutions.

We say a measured lamination $\lambda \in
ML(S)$
is {\it carried} by $\tau$ ($\lambda < \tau$)
if its support can be homotopied into $\tau$.
We denote
$$
\begin{aligned}
ML(\tau)&=\{\lambda \in ML(S):\lambda <
\tau\}\\
ML_{\mathbb{Z}}(\tau)&=\{\lambda \in
ML_{\mathbb{Z}}(S):\lambda <\tau\}\\
ML_{\mathbb{Z}}(S)&=\{\lambda \in
ML(S):\substack{\text{
simple multi-curve} \\ \text{with integral
coefficients}}\}.
\end{aligned}
$$
\end{definition}

\medskip\noindent
Now we will explain the fact that there are
one-to-one correspondences
$$
\begin{aligned}
ML(\tau)&\xymatrix{\ar@{<->}[r]&}w(\tau)\\
ML_{\mathbb{Z}}(\tau)&\xymatrix{\ar@{<->}[r]&
}w_{\mathbb{Z}}(\tau).
\end{aligned}
$$

\noindent
In a simple case where we pick a multicurve
$\lambda \in ML_{\mathbb{Z}}(\tau)$, the
associated measure
is just counting the intersection number
of the curve with a perpendicular segment
near
the train track edge.

The following argument appears to be a
converse construction.
Given a solution to the switch equations,
i.e. an element in $\omega(\tau)$, we can
produce
a multicurve $\lambda \in
ML_{\mathbb{Z}}(\tau)$.
(The method is essentially putting as many
strands
as the weights say near the vertex;
this yields formal sums of curves.)
This also works for formal sums with rational
coefficients.
And for non-negative real coefficients, we
take a thickening of $\tau$.

\begin{theorem}[Thurston]
\label{theorem-thurston-train-tracks}
The weights on $\tau$ uniquely determine a
transverse measure on $L$.
\end{theorem}

\medskip\noindent
Facts:
\begin{enumerate}
\item Every $\lambda \in MS(S)$
is carried by some train track $\tau$ and
determines
a unique point in $\omega(\tau)$.
\item There are finitely many train tracks
$\tau_i$ such that
every $\lambda \in ML(S)$ is carried by at
least $\tau_i$.
From this it follows that we can write
$ML(S)=ML(\tau_1)\cup  ML(\tau_2)\cup
\ldots\cup ML(\tau_n)$
and $ML(\tau)\approx \omega(\tau)\approx
\mathbb{R}^{E-V}$
(under very mild conditions on the train
tracks).

Here is an exercise: suppose $S$ is closed.
Show that $E-V \leq  6g-6$, and
$E-V=6g-6$ if $\tau$ is maximal (which means
that all complementary
regions are triangles).

\item $ML(\tau_i)\cap ML(\tau_j)$ for $i \neq
j$ has dimension
$<6g-6$.
\end{enumerate}

Given those facts,
$$
\begin{aligned}
ML(S)&=ML(\tau_1)\dot\cup \ldots \dot\cup
ML(\tau_n)\\
&=\omega(\tau_1)\dot\cup \ldots\dot\cup
\omega(\tau_n).
\end{aligned}
$$

\noindent
In particular, $ML(S)\approx
\mathcal{E}^{6g+6L+2r}$
($\omega(\tau_1)\approx\mathbb{R}^{6g+6L-2r}
$).
We may now introduce a measure on the space
of laminations:

\begin{definition}
\label{definition-thurston-measure}
The {\it Thurston measure} on $ML(S)$ is
$$
\lim_{L\to\infty}\frac{1}{L^{6g-6+2r}}
\sum_{\lambda
\in ML_{\mathbb{Z}}(S)}
\delta_{\frac{1}{L}\lambda}
$$
\end{definition}

\noindent
One checks that this is a Lebesgue measure
(in the limit) as follows:
$$
\begin{aligned}
ML_{\mathbb{Z}}(S)&=\bigcup_{i=1}
^nML_{\mathbb{Z}}(\tau_i)\\
&\frac{1}{L^{6g-6+2r}}\sum_{\lambda \in
ML_{\mathbb{Z}}\tau_1}
\delta_{\frac{1}{L}\lambda}\\
&\frac{1}{L^{6g-6+2r}}\sum_{(\omega)\in
\omega_{\mathbb{Z}}(\tau_i)}
\delta_{\frac{1}{L}\omega}\\
&\frac{1}{L^{6g-6+2r}}\sum_{p \in
\mathbb{Z}^{6g-6+2r}}
\delta_{\frac{1}{L}p}
\end{aligned}
$$

\noindent
on $\mathbb{R}^{6g-6+2r}$.

Notice that there is an action of the mapping
class group
on the space of laminations: a multicurve is
taken to another multicurve.

The following theorem was used for proving
the speaker's results:

\begin{theorem}[Masur '85]
\label{theorem-masur-85}
$m_{Thu}$ is $\text{Map}(S)$-invariant
and is ergodic with respect to
$\text{Map}(S)$.
\end{theorem}

\medskip\noindent
Now we discuss (geodesic) currents.
Consider de space $\mathcal{G}(\tilde{S})$
of unoriented geodesics on
$\tilde{S}=\mathbb{H}^2$
as $(S^1\times S^1\setminus\Delta)/(a,b)\sim
(b,a)$.
This  is well defined (even topologically,
i.e.
independently of the choice of metric)
since two geodesics on Poincaré plane are
determined by their endpoints.

\begin{definition}
\label{definition-geodesic-current}
A {\it (geodesic) current} is a
$\pi_1(S)$-invariant
Radon measure on $\mathcal{G}(S)$.
\end{definition}

\begin{example}
\label{example-multicurves-are-currents}
Any (multi-)curve is a current.
\end{example}

\noindent
We shall see that not only curves, but also
measured laminations
and in fact the Teichmüller space are
contained in currents:
$$
\begin{aligned}
\{\text{curves}\}&\subset C(S)\\
ML(S)&\subset C(S)\\
\mathcal{T}(S)&\subset C(S)
\end{aligned}
$$

Let $X$ be a hyperbolic metric on $S$.
Let $U=\{\lambda \in
ML(S):\ell_X(\lambda)\leq 1\}$.
If we let $m^L(U_X)=\frac{1}{L^{6g-6}}
\#\{\gamma \in
ML_{\mathbb{Z}}(S)|\ell_X(\gamma)\leq  L\}$,
then $m^L(U_X) \to m_{\text{Thu}}(U_X):=B(X)$
 as $L \to \infty$.

\medskip\noindent
In the following, we think os simple
curves as measured laminations and currents.
Recall that a measured lamination is
a union of disjoint curves.
$$
\xymatrix@C=4em{
ML_{\mathbb{Z}}(S)\ar@{^{(}->}[r]&\overline{
\{\text{simple
multicurves}\}}\ar@{^{(}->}[d]\ar@{^{(}->}[r]
^-{\substack{\text{equality for} \\
\text{closure only!}}}
&ML(S)\ar@{^{(}->}[d]\\
&\overline{\{\text{multi-curves}\}}\ar@{^{(}->}[r]^-{
\substack{\text{equality for} \\
\text{closure only!}}}
&CS=\{\text{geodesic currents}\}.
}
$$

\medskip\noindent
Facts:
\begin{enumerate}
\item (Banahon.) The geometric intersection
number of curves
extends to
$$
i:C(S) \times C(S) \to \mathbb{R}_{\geq 0}
$$
continuously, bi-linear,
$\text{Map}(S)$-invariant.

\item $ML(S)=\{\mu \in
C(S):i(\mu,\mu)=0\}\rightsquigarrow
m_{\text{Thu}}$ measure on $C(S)$
(supported on $ML(S)$).

\item $C_{\text{fin}}(S)=\{\mu \in C(S):
(\mu,\lambda)>0 \forall \lambda \in C(S)\},$
the {\it filling currents}, is an open dense
subset of $C(S)$.

\item (Banahan.) For all $X$ hyperbolic
metric on $S$,
there exists a unique $\lambda_X \in C(S)$
such that
$i(\lambda_X,\gamma)=\ell_X(\gamma)$ for all
curves $\lambda$
(Liouville current) $\rightsquigarrow $
$\mathcal{T}(S) \subset C(S)$.
This way we have a notion of hyperbolic
length of currents $\mu \in C(S)$:
$\ell_X(\mu)=i(\lambda_X,\mu)$.
\end{enumerate}

\begin{definition}
\label{definition-nice-length-function}
A length function on curves on $S$
is {\it nice} if it extends to a function
$F:C(S) \to \mathbb{R}_{\geq 0}$ wich is
continuous, homogeneous
($F(c\cdot\mu)=cF(\mu)\;\forall c>0$,
positive ($F(\mu)>0\;\forall \mu\neq 0$.
\end{definition}

\begin{example}
\label{example-nice-length-functions}
Following are examples of nice length
functions:
\begin{itemize}
\item Hyperbolic length.
\item (Otal.) Negatively curved Riemannian
metric.
\item (E-Parlier-Souto.) (Singular)
Riemannian metric.
\item (Duchen-Leinhein-Rafi.) Euclidean cone
surface.
\item (Banachan.) $\sigma \in C_{fill}(S)$,
$\ell_{\sigma}=(\sigma,M)$.
\item Word length w.r.t. finite generating
set $\pi_1(S)$.
\end{itemize}
\end{example}

\medskip\noindent
Let us go back to our original problem
of counting curves of a specific type.
Fix $\gamma_0$ a (multi-)curve (not
necessarily simple).
$\gamma$ is of type $\gamma_0$ if
$\gamma \in \text{Map}(S)\cdot \gamma_0$.

We want asymptotics of
$\#\{\gamma\text{ type
}\gamma_0:\ell_X(\gamma)\leq L\}$.
Fix $\gamma_0$ and $L>0$. Consider
$m^L_{\gamma_0}=\frac{1}{L^{6g-6}}
\sum_{\gamma\text{ type
}\gamma_0}\delta_{\frac{1}{L}\gamma}$
measure on $C(\gamma)$.

Note that $m^L_{\gamma_0}(U_X)=
\frac{1}{L^{6g-6}}\#
\{\gamma\text{ type
}\gamma_0:\ell_X(\gamma)\leq L\}$.

We are ready to state the main theorem:

\begin{theorem}[Mirzakhani '08, E-Souto '22]
\label{theorem-mirzakhani-e-souto}
For any $\gamma_0$ on $S=S_g$
$$
m^L_{\gamma_0}\to\frac{c(\gamma_0)}{b_g}
m_{\text{Thu}},\qquad
L\to \infty.
$$
\end{theorem}

\noindent
Mirzakhani proved it for simple curves,
while we now presented a general version
using currents.
We write the constants as
$$
\begin{aligned}
c(\gamma_0)&=m_{\text{Thu}}(\{\lambda \in
ML(S)/\text{Stab}(\gamma_0)
|i(\lambda,\gamma_0)\leq 1\},\\
b_g&=\sum_{\alpha_0 \in
ML_\mathbb{Z}/\text{Map}(S)}c(\alpha_0).
\end{aligned}
$$

\noindent
Mirzakhani used another expression.
This one has the virtue of being purely
combinatorial.

Mirzakhani showed the following example.
Consider the genus 2 surface $S_2$.
Then
$$
\lim_{L \to \infty}
\frac{\#\{\text{non-separating
}\gamma|\ell_X(\gamma)\leq L\}}{
\#\{\text{separating
}\gamma|\ell_X(\gamma)\leq L\}}
=48.
$$
This may be interpreted by saying that when
taking a closed simple geodesic,
the probability of it being separating is
$\frac{1}{48}$.






\section{Variations on a theorem of Baum and
Bott}
\label{section-baum-and-bott}

\noindent
Stéphane Druel, Université Claude Bernard
Lyon.
Algebra Seminar, IMPA.
November 5, 2025.

\medskip\noindent
{\bf Abstract.}
In this talk, I will discuss an analogue of
the Baum-Bott vansihing theorem for
the tangent bundle of the space of leaves of
regular foliations together with
applications to transversely projectively
flat foliations of dimension one. This
is work in progress.

\medskip\noindent



Baum-Bott 70'.

Let $E$ be a vector bundle on $X$
complex manifold is a sheaf of 1-jets of $E$:
$$
J'_X(E)= E \oplus \Omega^1_X \otimes E
$$
where equality is the category of sheavs of
abelian groups.
Where the $\mathcal{O}_X$ structure is given
by
$$
\underbrace{f}_{\substack{\text{local
function} \\ \text{on }U}}
(\underbrace{e}_{\substack{\text{local
section} \\ \text{of $E$ on  $U$}}}
, \underbrace{\alpha}_{\substack{\text{local
section} \\
\text{of $\Omega^1_X \otimes E$}}})
=(fe,f\alpha-df \otimes e)
$$
The {\it Atiyah class} of $E$ is the element
$$
\text{at}_X(E) \in H^1(X,\Omega_X^1 \otimes
\text{End}(E))
$$
corresponding to the exact sequence of
$\mathcal{O}_X$-modules
$$
\xymatrix{
0\ar[r]&\Omega_X^1\otimes
E\ar[r]&J_X^1(E)\ar[r]&E\ar[r]&0.
}
$$

\begin{remark}
\label{remark-holomorphic-connection}
 $\text{at}(E)=0$ if and only if there exists
$\nabla:E \to \Omega^1_X \otimes E$ map of
abelian sheaves
such that $e \mapsto  (e-\nabla e)$ is
$\mathcal{O}_X$-linear
$$
-\nabla fe=-f\nabla e - df \otimes e
$$
if and only if $E$ admits a holomorphic
connextion.
\end{remark}

\begin{example}
\label{example-line-bundl-holom-conne}
If $E$ is a line bundle, at
$(E)=(d\text{log}f_{ij})_{ij}
\in H^1(X,\Omega_X^1)$, $E=(f_{ij})$.
\end{example}

\begin{theorem}[Atiyah, '57]
\label{theorem-atiyah-57}
Let $r =\text{rank}E$, $k$ an integer, $X$ be
compact Kähler
and $\Sigma_k$ a polynomial of degree $k$ on
$M_r$
such that  $\Sigma_k(g)$ is the $k$-th
elementary
symmetric function on the characteristic
zeros of
$g \in M_r$. $S_k=$ polynomial given by
polarization.
$$
\begin{aligned}
M_r^{\otimes k}&\to \mathbb{C}\\
\text{End}E^{\otimes k}& \to \mathcal{O}_X\\
\text{Ex., $k=1$}\qquad
\text{End}E&\xrightarrow{\text{Tr}}
\mathcal{O}_X.
\end{aligned}
$$

$$
\xymatrix{
H^1(X,\Omega_X^1 \otimes
\text{End}(E))^{\otimes k}\ar[r]^{\cup }
&H^k((\Omega^1_X)^{\otimes k}\otimes
\text{End}E^{\otimes k}\ar[r]
&H^k(X,\Omega_X^k)\ar[r]
&H^{2k}(X,\mathbb{C})\\
\text{at}(E)^{\otimes k}\ar[rrr]
&&&(2\pi i)^kc_{2k}(E).
}
$$
(Double check the term $(\Omega^1_X)^{\otimes
k}$…
I think there's a typo.)

Partial connection. $\mathcal{F}$ regular
distribution.
$\mathcal{F} \subseteq T_X$ subbundle.
An {\it $\mathcal{F}$-connection} on $E$ is
$\nabla:E \to \mathcal{F}^\vee \otimes E$
such that
$$
\nabla(fe)=f\nabla e+f_{\mathcal{F}}f \otimes
e
$$
where
$d_{\mathcal{F}}:\mathcal{O}_X\xrightarrow{d}
\Omega_X^1
\to\mathcal{F}^\vee$.
\end{theorem}

\begin{proposition}[Bott-Baum, '70]
\label{proposition-bott-baum}
If $E$ admits an $\mathcal{F}$-connection
then all characteristic
classes of degree  $>\text{codim}\mathcal{F}$
vanish.
\end{proposition}

\begin{proof}
There exists an $\mathcal{F}$-connection if
and only if
$$
\xymatrix{
0\ar[r]&\Omega_X^1 \otimes
X\ar[r]\ar[d]&J_X^1\ar[r]&E\ar@{=}[d]\ar[r]&
0\\
0\ar[r]&\mathcal{F}^\vee\otimes E\ar[r]&
\mathcal{F}^\vee \otimes E \oplus E\ar[r]&
E\ar[r]&0
}
$$
if and only if $\text{at}(E) \in
H^1(X,\Omega_X^1\otimes\text{End}E)$
maps to 0 in
$H^1(X,\mathcal{F}^\vee\otimes\text{End}E)$
is and only if $\text{at}(E)$ comes from
$$
H^1(X,N^\vee\otimes\text{End}E)\to
H^1(X,\Omega_X^1\otimes\text{End}E)
$$
where $N=T_{X/\mathcal{F}}$.
\end{proof}

\begin{example}[Bott's partial connection]
\label{example-bott-partial-connection}
Suppose $\mathcal{F}$ is a regular foliation
and $N=T_{X/\mathcal{F}}$ has a
$\mathcal{F}$-connection.
Take $n$ a local section of $N$,  $v$ a local
section of $\mathcal{F}$
(over the same open) and suppose exists a
section $t$ of $T_X$
that is a lift of $t$, i.e. $p(t)=n$ for
$p:T_X \to N$.
Then $\nabla_v^Bn=p([v,t])$ is a section of
$N$.
This defines an $\mathcal{F}$-connection on
$N$.
\end{example}

\begin{lemma}[Beauville]
\label{lemma-beauville}
Let $X$ be a complex manifold and $E$ is a
direct submanifold of $T_X$.
Then $\text{at}(E) \in H^1(X,\Omega_X^1
\otimes \text{End}E)$
comes from $H^1(X,E^\vee \otimes
\text{End}(E))$.
\end{lemma}

\begin{proof}
Consider
$$
T_X= E \oplus F \xrightarrow{p}E.
$$
We claim that there exists an
$\mathcal{F}$-connection on $E$.
Let $v$ be a local section of $F$,
$e$ a local section of $E$.
Define $\nabla_ve=p([v,e])$.
One checks that this satisfies the required
properties.
\end{proof}

\noindent
Notice that no involutivity is imposed on $E$
- it's just a distribution.

\medskip\noindent
Now we will discuss some variations of
Baum-Bott's theorem
\ref{theorem-baum-bott}.

Let $\mathcal{F}$ be  regular foliation, $E$
a vector bundle on $X$ equipped with a flat
$\mathcal{F}$-connection.
For example, $\nabla^B$ flat.

Recall that  $(E,\nabla)$ can be thought of
as
$\text{Mon}(\mathcal{F})$-equivariant vector
bundles
where $\text{Mon}(\mathcal{F})$ is the
monodromy groupoid.

Let $X/\mathcal{F}$ be the space of leaves,
and
$\mathcal{O}_{X/\mathcal{F}}=\mathcal{O}
_X^{d_{\mathcal{F}}}=
\{\text{functions $f$ such that
$d_{\mathcal{F}}f=0$}\}$.

It is a fact that $E^{\nabla}$ is a locally
free $\mathcal{O}_{X/\mathcal{F}}$
sheaf of rank the rank of $E$ and such that
$E=E^{\nabla}\otimes_{F_{X/\mathcal{F}}}
\mathcal{O}_X$.

Now one defines bundles of 1-forms and
tangent bundle as follows.
 $\Omega_{X/\mathcal{F}}^1=(N_{\mathcal{F}}^*
 )^{\nabla B}$
and
$T_{X/\mathcal{F}}=N_{\mathcal{F}}^{\nabla
B}$,
which are locally free sheaves of
$\mathcal{O}_{X/\mathcal{F}}$-modules
of rank $q=\text{codim}\mathcal{F}$.

The Atiyah class of $(E,\nabla)$ is
$$
J^1_{X/\mathcal{F}}(E,\nabla)=E^{\nabla}
\oplus
\Omega_{X/\mathcal{F}}^1
\otimes_{\mathcal{O}_{X/\mathcal{F}}}
E^{\nabla}
$$
with $\mathcal{O}_{X/\mathcal{F}}$-module
structure given by
 $$
(e,\alpha)=(f\alpha,f\alpha-\underbrace{df}
_{\in
\Omega^1_{X/\mathcal{F}}}
\otimes e.
$$
Then
$$
\text{at}_{X/\mathcal{F}}(E,\nabla)
\in
H^1(X,\Omega_{X/\mathcal{F}}
^1\otimes_{\mathcal{O}_{X/\mathcal{F}}}
\text{End}_{\mathcal{O}_{X/\mathcal{F}}}
E^{\nabla})
$$
corresponding to
$$
\xymatrix{
0\ar[r]&\Omega_{X/\mathcal{F}}
^1\otimes_{\mathcal{O}_{X/\mathcal{F}}}\ar[r]
&J_{X/\mathcal{F}}^1(E)\ar[r]&E^{\nabla}
\ar[r]&0.
}
$$

\begin{lemma}
\label{lemma-atiya-class-maps-atiya-class}
$\text{at}_{X/\mathcal{F}}(E,\nabla)$ maps to
$\text{at}_X(E)$ under
$$
H^1(X,\Omega_{X/\mathcal{F}}
^1\otimes_{\mathcal{O}_{X/\mathcal{F}}}
\text{End}_{\mathcal{O}_{X/\mathcal{F}}}
E^{\nabla}
\to H^1(X,\Omega_X \otimes \text{End}E).
$$
\end{lemma}

\noindent
Partial connections. Let $G \subseteq
T_{X/\mathcal{F}}$
be a subbundle a {\it $G$-connection}
$(E,\nabla)$ is
$$
\begin{aligned}
D: E^{\nabla} &\longrightarrow G^\vee\otimes
E^{\nabla} \\
\text{s.t.} D(fe)&=fDe+d_Gf\otimes e
\end{aligned}
$$
where
$d_G:\mathcal{O}_{X/\mathcal{F}}\to\Omega_{X/
\mathcal{F}}^1\to
G^\vee$.

It's a fact that a $G$-connection
$(E,\nabla)$ exists if and only if
$\text{at}_{X/\mathcal{F}}(E,\nabla)$ maps to
0 in
$H^1(X,G^\vee
\otimes_{G_{X/\mathcal{F}}}\text{End}
_{\mathcal{O}_{X/\mathcal{F}}}
E^{\nabla}$.

The analogue of Beauville's lemma
\ref{lemma-beauville} is

\begin{lemma}[Beauville]
\label{lemma-analogue-beauville}
Suppose $N_{\mathcal{F}}=\mathcal{A} \oplus
B$ such that
$\nabla^B\mathcal{A} \subseteq
\mathcal{F}^\vee
\otimes_{\mathcal{O}_X}\mathcal{A}$
and $\nabla^B \mathcal{B} \subseteq
\mathcal{F}^\vee\otimes_{\mathcal{O}_X}
\mathcal{B}$.
Let $\nabla_{\mathcal{A}}$ and
$\nabla_{\mathcal{B}}$ be a flat
 $\mathcal{F}$-connection on $\mathcal{A}$
 and $\mathcal{B}$.

Then
$\text{at}_{X/\mathcal{F}}(\mathcal{A},
\nabla_{\mathcal{A}})
\in
H^1(X,\Omega_{X/\mathcal{F}}
^1\otimes_{\mathcal{O}_{X/\mathcal{F}}}
\text{End}_{\mathcal{O}_{X/\mathcal{F}}}
\mathcal{A}^{\nabla_{\mathcal{A}}}$
comes from
$H^1(X,(\mathcal{A}^{\nabla_{\mathcal{A}}})
^\vee
\otimes_{X/\mathcal{F}}\text{End}
_{\mathcal{O}_{X/\mathcal{F}}}\mathcal{A}
^{\nabla_{\mathcal{A}}}$.
\end{lemma}

\begin{proof}
Same proof but we need a Lie bracket on
$T_{X/\mathcal{F}}=N_{\mathcal{F}}\subseteq
N_{\mathcal{F}}\overset{p}{\leftarrow}T_X$.
Let $v_1,v_2$ be local sections on
$T_{X/\mathcal{F}}$
with $v_1=p(t_1)$ and $v_2=p(t_2)$.
Define
$$
[v_1,v_2]=p([t_1,t_2]).
$$
One can check that this a Lie bracket on
$T_{X/\mathcal{F}}$.
Suppose that $t_1 \in \mathcal{F}$. Then
$$
p([t_1,t_2])\overset{\text{def}}{=}
\nabla_{t_1}^{\mathcal{B}}v_2=0.
$$
For $v_3 \in \mathcal{F}$,
$$
\nabla_{v_3}^{\mathcal{B}}[v_1,v_2]
=p([v_3,[t_1,t_2]])=\pm\nabla_{v_3}
^{\mathcal{B}}v_1\pm
\nabla_{v_3}^{\mathcal{B}}v_2=0.
$$
\end{proof}

\begin{lemma}
\label{lemma-vanishing-chern-class}
Suppose that $\mathcal{F}$ is a regular
foliation on $X$.
Let $N_{\mathcal{F}}=\mathcal{L}_1 \oplus
\ldots \oplus \mathcal{L}_q$
be line bundles on $X$,
$\nabla^{\mathcal{B}}\mathcal{L}_i\subseteq
\mathcal{F}^\vee\otimes\mathcal{L}
_i\rightsquigarrow
\nabla_i$.
Suppose
$(\mathcal{L}_i,\nabla_i)\cong(\mathcal{L}_j,
\nabla_j)$
for all $i,j$.
Then $c_1(\mathcal{L}_i)=0$.
\end{lemma}

\begin{proof}
By the Comparison Lemma,
$$
\begin{aligned}
H^1(X,\Omega_{X/\mathcal{F}}^1
&\longrightarrow H^1(X,\Omega_X^1) \\
\text{at}_{X/\mathcal{F}}(\mathcal{L},
\nabla_i)
&\longmapsto
\text{at}(\mathcal{L}_i),
\end{aligned}
$$

\noindent
it suffices to prove that
$\text{at}_{X/\mathcal{F}}(\mathcal{L}_i,
\nabla_i)$
vanishes.
$$
\text{at}_{X/\mathcal{F}}(\mathcal{L}_i,
\nabla_i)
\in H^1(X,\Omega_{X/\mathcal{F}}^1)
=\bigoplus_{j=1}^9
H^1(X,(\mathcal{L}_j^{\nabla_j})^\vee)
$$
comes from
$H^1(X,(\mathcal{L}_i^{\nabla_i})^\vee)$.
\end{proof}

\begin{example}
\label{example-degree-2-line-bundle-on-e}
$A=E^1 \times B^{n-1}$, an abelian variety.
Let $ L$ be the pullback of a degree 2
divisor on $E$.
Then there exists  $T_X
\to\!\!\!\!\!\to L^{\oplus n-1}$.
\end{example}

\begin{proposition}
\label{proposition-}
$\mathcal{F}$ regular foliation of dimension
1 on $X$ complex projective
of dimension $\geq 3$.
Suppose $\mathcal{F}$ is transversally
projective flat
and that $X$ contains a rational curve.
Then $\mathcal{F}$ is a $\mathbb{P}^1$-bundle
and it is tangent
to $\mathcal{F}$.
\end{proposition}

\medskip\noindent
Now we explain what it means to be
projectively flat.
It is when the projectivization of the normal
bundle $\mathbb{P}(N)$
is induced by a representation of $\pi_1(X)
\to \text{PGL}_g$
it comes with an Ehresmann connection on
$$
\xymatrix{
\mathbb{P}(N)=\mathbb{P}\ar[d]^{f}\\
X
}
$$
$T_{\mathbb{P}}=T_f \oplus g$.

Now, it is also called transversally
projectively flat
if $\nabla^B$ is a partial connection on $N$
and there exists a foliation $\mathcal{H}$ on
$$
\xymatrix{
\mathbb{P}(N)\ar[d]\\ X
}
$$
such that $\mathcal{H}$ projects onto
$\mathcal{F}$.
Transversality means that $\mathcal{H}
\subseteq \mathcal{G}$.

Finally, transversally projectively flat
means that
it is induced on a
$\text{Mon}(\mathcal{F})$-module
by a representation of
$\pi_1(\text{Mon}(\mathcal{F}) \to
\text{PGL}_g$.


\section{Dolbeault cohomology of
nilmanifolds}
\label{section-dolbe-cohom-nilma}

\noindent
Misha Verbitsky, IMPA.
Geometric Structures Seminar, IMPA.
November 6, 2025.

\medskip\noindent
{\bf Abstract.}
A nilmanifold is a (left) quotient of a
nilpotent Lie group by a cocompact
lattice. If this Lie group is equipped with a
left-invariant complex structure,
the quotient is called a complex nilmanifold.
Examples of complex nilmanifolds
include Kodaira-Thurston surface and many
other non-Kahler complex manifolds.
The differential graded algebra of invariant
differential forms on a nilmanifold
has the same cohomology as a nilmanifold, as
shown by Nomizu. I am going to
explain a Dolbeault cohomology version of
this theorem. A Lie algebra of a
nilmanifold admits a natural rational
structure, induced by the cocompact
lattice. If the complex structure operator,
acting on the Lie algebra, has
rational coefficients, Console and Fino
proved that the Dolbeault cohomology of
a nilmanifold is equal to the Dolbeault
cohomology of its invariant forms. I
will describe an attempt to prove this result
for a general nilpotent Lie
algebra, based on number-theoretic arguments.

\medskip\noindent
Let $G$ be a nilpotent Lie group.
[For me this would mean that the Lie algebra
of $G$ is
nilpotent as in Lie Algebras Definition
\ref{lie-algebras-definition-nilpo-%
solva-lie-algeb},
but according to Misha this would be the case
for $G$
connected only; so there's must some other
notion
of nilpotent Lie groups.]
A {\it nilmanifold} $M$ is a manifold with a
transitive action
of a nilpotent Lie group. This means that
$M=G/\Gamma$.

\begin{theorem}[]
\label{theorem-}

\end{theorem}

\noindent
Let $G$ be a real Lie group
and $I \in \text{End}(TG)$ with
$I^2=-\text{Id}$.
Suppose $I$ is right-invariant.
Then $I$ is integrable iff
$[T^{1,0}G,T^{1,0}G]\subset T^{1,0}G$
iff
$\mathfrak{g}^{1,0}\subset\mathfrak{g}
\otimes\mathbb{C}$
where $\mathfrak{g}^{1,0}$ is the
$\sqrt{-1}$-eigenspace of $I$.

A {\it complex nilmanifold} is
$(G,I)/\Gamma$.

\begin{example}
\label{example-iwasawa-manifold}
An {\it Iwasawa manifold} is a quotient of
the set
of upper triangular matrices with 1 on the
diagonal
(this group is $U_3(\mathbb{C})$) by a
lattice $\Gamma$.
For example, we can take
$\Gamma=U_3(\mathbb{Z}[\sqrt{-1}])$.
\end{example}

\noindent
The complex structure $I$ is bi-invariant iff
it is paralelizable.

\medskip\noindent
Let $E$ be an elliptic curve with an ample
line bundle $L$.
We put an action on $\text{Tot}L$ given by
multiplication by
$\lambda \in \mathbb{C}$ for $|\lambda|>0$.

\begin{theorem}
\label{theorem-de-rham-cheva-eilen}
De Rham cohomology of a nilmanifold coincides
with the
CE cohomology of $\mathfrak{g}$.
\end{theorem}

\begin{theorem}[Console-Fino]
\label{theorem-console-fino}
For complex nilmanifolds,
$H^{\bullet,\bullet}(\mathfrak{g}^\bullet,
\overline{\partial}$
and
$H^{\bullet,\bullet}_{\overline{\partial}}(M)
$
are isomorphic.
\end{theorem}

\noindent
There's several results addressing a converse
situation
(when the algebras are isomorphic) and other
situations
involving fibrations.

\medskip\noindent
Now we make a sudden change to deformation
theory.

\begin{lemma}
\label{lemma-}
Complex spaces on $V=\mathbb{R}^{2n}$ are
uniquely determined
by $W \in V \otimes_{\mathbb{R}}\mathbb{C}$
such that $W \oplus \overline{ W}=V \otimes
\mathbb{C}$
which is an element in $\text{Gr}(n, V
\otimes \mathbb{C})$.
\end{lemma}

\begin{definition}
\label{definition-}
Let $\mathfrak{g}$ be a Lie algebra and
$\mathfrak{g}_\mathbb{C}$
its complexification.
The set of complex structures on
$\mathfrak{g}$ is $\text{Comp}(\mathfrak{g})$
which is $\{ W \in
\text{Gr}_0(\mathfrak{g}_\mathbb{C}:
W \subset W \subset
\mathfrak{g}_{\mathbb{C}}\text{ is a
subalgebra}\}$
since being a subalgebra means
$[\mathfrak{g}^{1,0},\mathfrak{g}^{1,0}]
\subset \mathfrak{g}^{1,0}$.
\end{definition}

We can define a variety out of this:
$$
\overline{\text{Comp}(\mathfrak{g}
_{\mathbb{C}}}
=\{W \in
\text{Gr}(n,\mathfrak{g}_{\mathbb{C}}:
\text{$W$ subalgebra}\}
$$
is a complex subvareity of
$\text{Gr}(n,\mathfrak{g}_{\mathbb{C}})$.


\begin{definition}
\label{definition-field-of-defition}
Let $X \subset \mathbb{C}P^n$ be a projective
subvariety.
Consider the smallest field $k \supset
\mathbb{Q}$ such that
any ideal $I$ of $X$ [ideal in the coordinate
algebra of $X$?]
can be generated by polynomials with
coefficients
in $k[k^n]$ is called a {\it field of
definition}.
\end{definition}

\noindent
Apparently there's some discussion about this
notion
not being easily defined.

\begin{lemma}
\label{lemma-}
$\text{Gal}_{\mathbb{Q}}(k)$ acts on
$\mathbb{P}_k^n$.
Let $X \subset \mathbb{P}^n$ be a projective
subvariety
and $\Gamma \subset
\text{Gal}_{\mathbb{Q}}(k)$
the stabilizer of $X$.
Then the field of definition of $X=\{\lambda
\in k:\Gamma(\lambda)=\lambda\}$.
\end{lemma}

\noindent
We have a simpler definition of field of
definition for a point
(which is apparently the former definition
when the variety is a point)

\begin{definition}
\label{definition-field-of-definition-point}
For a point $(z_1,\ldots,z_n) \in
\mathbb{C}^n$, its
{\it field of definition} is the field
generated by $z_1,\ldots,z_n$.
\end{definition}

\noindent
So, for instance, a rational point has
transcendence degree 0.

\begin{remark}
\label{remark-}
Let $X$ be an $n$-dimensional algebraic
variety
defined over $\mathbb{Q}$ or any finite
extension of $\mathbb{Q}$,
and let $\pi:X\to \mathbb{C}^n$ be finite
dominant over $\mathbb{Q}$.
\end{remark}


\noindent
We have a fun theorem:

\begin{theorem}
\label{theorem-fun-theorem}
Let $X$ be an irreducible complex algebraic
variety defined over $\mathbb{Q}$.
Then the group
$\text{Aut}_\mathbb{Q}(\mathbb{C})$ acts
transitively
on points of maximal transcendence degree in
$X$.
\end{theorem}

\begin{definition}
\label{definition-algebraic-point}
We say that a point $x \in X$ is {\it
algebraic} if
its transcendence degree over $k_X$ is zero.
\end{definition}

\noindent
It's not very clear whether the variety $X$
is projective or affine.

\begin{lemma}
\label{lemma-algebraic-points-are-dense}
The set of algebraic points in
$X\subset\mathbb{C}P^n$ is dense
in the analytic topology.
\end{lemma}

\begin{theorem}
\label{theorem-aut-orbit-of-point}
Let $X$ be a complex (affine) algebraic
variety defined over $\mathbb{Q}$.
Consider an
$\text{Aut}_{\mathbb{Q}}\mathbb{C}$-orbit of
a point
$p \in X$. Then the closure of its orbit is a
minimal subvariety $Z \subset X$ defined
over $\mathbb{Q}$ and containing $p$.
\end{theorem}

\noindent
Now we introduce a Galois action on Dolbeault
cohomology.

\begin{theorem}
\label{theorem-}
$\mathfrak{g}$ nilpotent Lie algebra, $M$
corresponding nilmanifold,
$\sigma \in
\text{Aut}_{\mathbb{Q}}(\mathbb{C})$,
and let $J=\sigma(I) \in \text{Comp}(G)$.
Then $\dim H^{p,q}_{\overline{\partial}}(M,I)
=\dim H^{p,q}_{\overline{\partial}}(M,J)$ for
all $p,q$.
\end{theorem}

From this it follows that

\begin{lemma}
\label{lemma-console-fino-complex-structure}
For a nilmanifold with complex structure $I$,
if the field of definition of $I$ is
$\mathbb{Q}$,
then the Console-Fino map
$H_{\overline{\partial}}^{p,q}(\mathfrak{g})
\to
H_{\overline{\partial}}^{p,q}(M)$ is an
isomorphism.
\end{lemma}

\noindent
The idea is that if we know the result for
algebraic
points then we'll have it for transcendental
points as well.
Namely,

\begin{lemma}
\label{lemma-if-conso-fino-algeb-then-hold}
If Console-Fino is an isomorphism for all
algebraic complex structures,
then it holds for all transcendental complex
structures too.
\end{lemma}

\begin{proof}
Let $I \in \text{Comp}(\mathfrak{g})$ be
trascendental.
Consider the variety
$\overline{\text{Aut}_{\mathbb{Q}}\mathbb{C}}
$,
which is algebraic and thus contains
algebraic points.
Then Console-Fino is an isomorphism on such
points.
Then it is so in a neighbourhood of this
point.
Then it is so in the orbit of this
neighbourhood.
\end{proof}

\noindent
Consider $(\mathfrak{g},I)$ with $I$ with
number field $k$.
Consider
$\tilde{\mathfrak{g}}=\bigoplus_{\sigma \in
\text{Gal}(k:\mathbb{Q})}
(\mathfrak{g},\sigma(I))$…
this constructions leads to a Lie algebra
$(\tilde{\mathfrak{g}},\tilde{I})$
equipped with a complex structure defined
over $\mathbb{Q}$.
This satisfies Console-Fino. We can lift
$(\tilde{M},\tilde{I})$
since it remains Dolbeault harmonic --- which
is the main tool used by
Colsole-Fino in their theorem.
This is how to proof Console-Fino for complex
structures over a closed field.





\section{Holomorphic symmetric differentials,
semiample bundles,
fundamental groups}
\label{section-holom-symme-diffe-semia}

\noindent
Ernesto Mistretta, Università di Padova.
Algebra Seminar, IMPA.
November 12, 2025.

\medskip\noindent
{\bf Abstract.}
We will present some recent results and some
open conjectures relating
holomorphic symmetric differentials and the
geometry or topology of a compact
complex manifold, we will illustrate the
difference in the projective, or
Kaehler, or compact complex cases. Then we
will apply some elementary
constructions on semiample vector bundles to
partially answer some of these
questions.



\medskip\noindent



\noindent
As a motivation,
$$
\xymatrix{
&\text{differential forms}\ar[ld]\ar[rd]\\
\text{topology}&&\text{geometry}
}
$$
$X$ compact complex manifold, $E$
(holomorphic) vector bundle.
$L$ line bundle, the  {\it Kodaira-Itaka
dimension} is
$\text{KI}(X,L)=k$ if $h^0(X,L^{\otimes
m})=\mathcal{O}(m^k)$.

{\it Holomorphic $p$-differentials} are
elements of
$H^0(X,\text{Sym}^m\Omega^p_X)$ for some
$m>0$.

A conjecture by Mumford: suppose that $X$ is
a compact
Kähler manifold, then $X$ is rationally
connected if and only if
for every  $m>0$ and for every $p>0$,
$H^0(X,\text{Sym}^m\Omega^p_X)=0$.

Some relevant results are the following.

\begin{theorem}[Campana-Demailly-Peterwell,
2015]
\label{theorem-cdp}
Suppose $X$ is a smooth and projective
variety over $\mathbb{C}$.
$X$ is rationally connected if and only if
for every
ample line bundle there exists a positive
constant $C_A>0$,
such that $\forall p>0$, $\forall K>0$ and
$\forall m>C_A K$,
$$
H^0(X,\text{Sym}^m\Omega^p_X \otimes A^K)=0.
$$
\end{theorem}

\noindent
Recall that if $X$ is a Kähler manifold and
rationally connected,
then $X$ is projective and $\pi_1(X)=0$.

\begin{theorem}[Brunarbe-Campana, 2016]
\label{theorem-bc16}
Let $X$ be compact Kähler.
Suppose that $\forall m>0$, $\forall p>0$,
$H^0(X,\text{Sym}^m\Omega^p_X)=0$.
Then $X$ is projective and $\pi_1(X)=0$.
\end{theorem}

\noindent
What if we restrict to $p=1$?
$H^0(X,\text{Sym}^m\Omega^1_X)$?

A conjecture by Esnault says: let $X$ be a
compact Kähler manifold.
Suppose that $\forall m>0$,
$H^0(X,\text{Sym}^m\Omega^1_X)=0$.
Then $|\pi_1(X)|<\infty$.

\begin{theorem}[Burnabbe-Klinger-Totsaro]
\label{theorem-bkt}
Suppose $X$ is a compact Kähler manifold.
Suppose that there exists a
finite-dimensional
representation of $\rho:\pi_1(X)\to
\text{GL}_r(k)$
over some field $k$ such that
$|\pi(\pi_1(X))|=\infty$.
Then there exists $m>0$ such that
 $H^0(X,\text{Sym}^m\Omega_X^1) \neq 0$.
\end{theorem}

\noindent
Notice the relationship between Theorem
\ref{theorem-bkt}
and the conjecture by Esnault.

\medskip\noindent
Conjecturally, we have:
\begin{enumerate}
\item No holomorphic symmetric
$p$-differentials for all $p$
if and only if $X$ is rationally connected.
\item No holomorphic symmetric
1-differentials, them $|\pi_1(X)|<\infty$.
\end{enumerate}

\noindent
Question. Suppose there are some (many)
nonzero holomorphic symmetric
1-differentials, does this imply that
$|\pi_1|=\infty$?
What if $X$ is compact Kähler? And if $X$ is
projective?

Let's first answer the easiest of these
questions.
Let $X$ be compact Kähler with
$H^0(X,\Omega^1_X)\neq 0$.
By Hodge theorem,
$b_1(X)=\dim_\mathbb{Q}H^1(X,\mathbb{Q})
=2\dim_\mathbb{C}H^0(X,\Omega_X^1)$.
Also,
$H_1(X,\mathbb{Z})=\text{Ab}(\pi_1(X))$.
In particular, if $b_1 \neq 0$,
then
$H_1(X,\mathbb{Z})=\mathbb{Z}^{b_1}\oplus$
torsion.

\begin{theorem}[Brotbeck-Drondeau,2018]
\label{theorem-brotbeck-drondeau18}
$X \subseteq \mathbb{P}^N$ very gen. complete
intersection
of high degree such that $\dim X\leq N/2$.
Then $\Omega_X^1$ is ample.
(In particular, $\text{Sym}^m\Omega_X^1$
is globally generated for $\gg m$.)
Now if $\dim X\geq 2$, then
$\pi_1(X)=\pi_1(\mathbb{P}^N)=0$.
\end{theorem}

\noindent
In this case $\omega_X=\det(\Omega_X^1)$ is
an ample line bundle.
So the Kodaira dimension of $X$ is
$\text{kod}(X)=\text{kod}(X,\omega_X)=\dim(X)
$.

\medskip\noindent
Now we explain joint work with Franccesco
Exposito from Padova.

First, a conjecture (by E-M).
Suppose $X$ is a compact complex manifold
with $\text{kod}(X)<\dim(X)$. If $\Omega_X^1$
is (weakly) semiample, then
$|\pi_1(X)|=\infty$.

The evidence is:
\begin{enumerate}
\item If $\dim X=1,2$, holds because of
classification.
\item If $X$ is projective, then the
conjecture follows from
a result of Hörng, which states that
$\Omega_X^1$ is semiample,
then there exists $X'\to X$ finite étale
cover
such that $X'\sim_{\text{bitat}}Y\times A$
where $A$
is an abelian variety and  $\dim
Y=\text{kod}(X)$.
\item If $X$ is Kähler, the conjecture (E-M)
follows
from a conjecture of Wu-Zheng from 2002,
which says that
if  $\Omega^1_X$ is (NEF) semiample (which
implies NEF), then
there exists $X' \to X$ finite étale that
admits a smooth fibration
$X' \to Y$ in complex tori with $\dim
Y=\text{kod}X$.
In this case the fundamental groups of the
tori are injected
in the fundamental group of $X$, i.e.,
$\pi_1(F) \hookrightarrow \pi_1(X)$.
(It appears that this is proved via the Serre
homotopy
long exact sequence, though the Kähler
hypothesis may needed, and
it's not clear how…)
This makes $|\pi_1(X)|=\infty$.

\item If $X$ is a compact complex manifold
and $\text{kod}(X)=0$,
and $\Omega_X^1$ is (weakly) semiample, then
$|\pi_1(X)|=\infty$.
\end{enumerate}

\noindent

\begin{definition}
\label{definition-}
Let $E$ be a vector bundle on $X$, a compact
complex manifold.
Recall the construction
$$
\xymatrix{
\mathbb{P}(E)\ar[d]^{\pi}\ar@{-}[r]&\pi^*
E\ar@{->>}[r]
&\mathcal{O}_{\mathbb{P}(E)}(1)\\
X\ar@{-}[r]&E
}
$$
$$
\pi_*(\mathcal{O}_{\mathbb{P}(E)}(m))=
\text{Sym}^m(E).
$$
\begin{enumerate}
\item $E$ is {\it strongly semiample} if
there exists $m>0$
such that $\text{Sym}^mE$ is globally
generated.
\item $E$ is {\it weakly semiample} if
$\mathcal{O}_{\mathbb{P}(E)}(1)$ is semiample
line bundle
(i.e. $\mathcal{O}_{\mathbb{P}(E)}(m)$ is
globally generated
for some $m\gg 0$).
\end{enumerate}
\end{definition}

\noindent
Notice that these two definitions are not
equivalent.
Strongly implies weakly, but not conversely.
To see that strongly implies weakly, consider
the pullback
$$
\xymatrix{
\pi^*\text{Sym}^nE=\text{Sym}^n\pi^*E\ar[r]
\ar[d]&\mathbb{P}(E)\ar[d]^{\pi}\\
\text{Sym}^mE\ar[r]&X
}
$$
then
 $$
\text{Sym}^m\pi^* E
\to\!\!\!\!\!\to\mathcal{O}
_{\mathbb{P}(E)}(m).
$$
To see that weakly does not imply strongly,
suppose $X$ is a smooth projective
surface with $\text{kod}(X)=0$.
Consider $E=\Omega_X^1$ semiample. We have
that $\Omega_X^1$
is strongly semiample if and only if $X$ is
an abelian surface.
Also, if $X$ is a hyperelliptic surface, then
$\Omega_X^1$
is weakly semiample.
This means that $X$ hyperelleptic surface is
the counterexample.

\begin{theorem}[E-M]
\label{theorem-em}
Suppose that $X$ is compact complex and that
$\text{kod}(X)=0$. Then
\begin{enumerate}
\item if $\Omega_X^1$ is strongly semiample
then
$X$ is a parallelalizable complex manifold.
\item If $\Omega_X^1$ is weakly semiample,
then $X$ is a finite étale Galois quotient
$X\cong P/G$
of a parallelizable manifold $P$.
\end{enumerate}
\end{theorem}

\begin{definition}
\label{definition-paralellizable}
$X$ is {\it parallelizable} is
$\Omega_X^1\cong\mathcal{O}_X$.
\end{definition}

\noindent
A $X$ is a compact complex manifold is
parallelizable if and only if
$X \cong H/\Gamma$ for $H$ a complex Lie
group and $\Gamma\subset H$
a discrete subgroup.
Further, it is Kähler if and only if $H$ is a
complex torus.

\section{Um passeio nos trabalhos de
Gerd Faltings}
\label{section-faltings}

\noindent
Hossein Movasati, IMPA.
Pelestra especial,
May 19, 2026.

\medskip\noindent
{\bf Abstract.}
Nesta palestra vamos dar um passeio nos
trabalhos de G. Faltings, começando por sua
demonstração da conjetura de Mordell em
1983, que lhe trouxe a medalha Fields. Uma
boa parte de sua carreira foi elaborar o
lado aritmético de um dicionário heurístico
proposto por P. Deligne em 1970, que tem
origem na teoria conjetural de motivos de
Grothendieck. Isso inclui versão p-ádica de
teoria de Hodge, domínios de períodos,
correspondência de Simpson, etc. Várias
propriedades aritméticas de variedades
abelianas são outro foco da carreira dele
que vamos abordar nesta palestra. Em 2026
ele ganhou o prêmio Abel por suas
contribuições fundamentais em geometria
aritmética.

\medskip\noindent
{\bf Upshot.} 
Elliptic curves are diophantine equations.
Some are simple and fun to solve,
some are extremely complicated ---
like Fermat's last theorem.
Finding solutions to a diophantine
equation is like finding $k$-points
on an algebraic curve for suitable
choices of fields.
In this sense,
Mordell's conjecture, now Faltings'
theorem is a statement of finiteness
of solutions of diophantine equations.
Other results are adaptations
of results from complex geometry
to finite fields.

\medskip\noindent
{\bf Notes.}

Finding roots of a polynomial $f \in
\mathbb{Q}[x,y]$ becomes the geometric
problem of finding the $\mathbb{Q}$-rational
points of $\operatorname{Spec}(\mathbb{Q}[x,y]/(f))$.
Or, one can pass to the projective version
$\operatorname{Proj}(\mathbb{Q}[x,y,z]/(F))$,
where $F$ is the homogenization of $f$.
We write
$C(\mathbb{Q})=\{[x:y:z] \in
\mathbb{P}^2(\mathbb{Q}) \mid
F(x,y,z)=0\}$
and similarly $C(\mathbb{C})$ for the
analogous set.
 
It turns out that the topology
of the curve, via its genus, strongly
constrains the $\mathbb{Q}$-rational points.
The easiest cases are:
\begin{itemize}
\item ($g=0$.) Here $d=1,2$;
then the curve is projectively equivalent to
a conic, and over $\mathbb{Q}$ it has either
no rational points or infinitely many.

\item ($g=1$.) Here $d=3$; this is more difficult
but still well-studied; see the
Birch--Swinnerton-Dyer conjecture. 

\item ($g\geq 2$.) Here $d\geq 4$.
This is much more complicated!
Faltings'
theorem implies that $C(\mathbb{Q})$ is
finite.
\end{itemize}

\noindent
This is the former Mordell conjecture
from 1962.

\begin{theorem}[Faltings, 1983]
For any smooth curve of genus $\geq 2$ over
$\mathbb{Q}$, we have $\#C(\mathbb{Q})<
\infty$.
\end{theorem}

The same theorem remains true after
replacing $\mathbb{Q}$ by a number field,
that is, a finite field extension of
$\mathbb{Q}$.

\begin{theorem}[Manin 1963, Grauer 1965]
\label{theorem-manin-grauer}
For any smooth non isotrivial curve of
genus $\geq 2$ over $k=\mathbb{C}(t)$
we have
$$
\# C(k) < \infty.
$$
\end{theorem}

\begin{proof}
Using Gauss-Manin connection ---
the name of which was invented by 
Grothendieck.
\end{proof}

\noindent
The main difference between
$\mathbb{C}(t)$
and $\mathbb{Q}$
is that in $\mathbb{C}(t)$
we have a derivation
$$
\frac{\partial }{\partial t}:
\mathbb{C}(t) \to \mathbb{C}(t)
$$
satisfying the Leibniz identity.

\medskip\noindent
The proof is somehow related
to Shafarevich and Tate conjectures.
The former states there's only
finitely many isomorphism
classes of abelian varieties of a
given dimension with
fixed polarization degree, defined
over a field with
good reduction outside a finite set
of places.

Tate conjecture was proved for finite
fields by Tate in 1966,
and by Faltings for number fields
in 1983.
It relates maps between the tangent
spaces of two abelian varieties
and maps between the varieties
themselves:
$$
\Hom_K(A,B) \otimes \mathbb{Z}_p
\cong
\Hom_{Gal(\tilde{K}/K)}
(T_pA,T_pB)
$$
where $K$ is a field of characteristic
$\neq p$, and $A,B$ are abelian over $K$.

Further, the proof uses Arakelov's
theory of arithmetic genus
and Faltings height.

\medskip\noindent
{\bf Abelian varieties.} 

Any curve (over any field)
can be embedded in some variety
called the Jacobian variety.
Over $\mathbb{C}$,
abelian varieties are just tori:
$$
\mathbb{C}^n/\Lambda
\cong_{C^\infty}
S^1\times… S^1,\qquad  2n \text{ times}.
$$

An abelian variety of dimension one
over a field $k$ is called
an {\it elliptic curve}.
They are given by diophantine
equations of degree 3 in two variables:
$$
y^2+a_1xy+a_3y
=
x^3+a_2x^2+a_4x,
\qquad a_1,…,a_6 \in k.
$$

The name was invented by A. Weil.
Before him, many mathematicians like E.
Picard,
G. Humbert, were using ``hyperelliptic
surface'' for Abelian surfaces.

\medskip\noindent
{\bf Shavarevich's conjecture.} 
What is the number of non-isomorphic elliptic
curves over $\mathbb{Q}$ and with bad
reduction only for primes $2,3$?

\medskip\noindent
A consequence of Tate conjecture
and Chebotarev density theorem is:;

\begin{theorem}
\label{theorem-tate-chebot-consequence}
LEt $E_i$, $i=1,2$ be two
elliptic curves over $\mathbb{Q}$.
TFAE:
\begin{enumerate}
\item $E_1$ is isogeneous to $E_2$
over $\mathbb{Q}$.
\item For all except a finite number
of primes,
$$
\# E_1(\mathbb{F}_p)
=\# E_2(\mathbb{F}_p).
$$
\end{enumerate}
\end{theorem}

\begin{definition}
\label{definition-isogeny}
{\it AI-generated.}
An isogeny between elliptic curves over
$\mathbb{Q}$ is a nonzero morphism of
elliptic curves defined over $\mathbb{Q}$
that is also a group homomorphism and has
finite kernel.
\end{definition}

Here isogenous means that there exists an
isogeny $f:E_1 \to E_2$ over $\mathbb{Q}$.

\medskip\noindent
One of the key ingredients
in Falting's proof
is the notion of {\it Faltings height},
which are integrals used to
control the numerator and denominator
or rational solutions of a (diophantine?)
equation.

For example, for $r = n/m$,
$h(r)=n^2+m^2$. For all $M \in \mathbb{R}$,
$$
\{r \in \mathbb{Q}| h(r)<M\}
$$
is finite.

\begin{proposition}[Main property
of height functions]
\label{proposition-faltings-height}
For any positive integer $M$,
the number of elliptic curves $E/\mathbb{Q}$
with $h_{\text{Falt}}<M$ is finite.
\end{proposition}

\noindent
On the curve given by Equation
\ref{equation-?}
consider the 1-form
$$
\omega=\frac{dx}{2y+a_1x+a_3}.
$$
The Falting's height is basically…
[missing! an integral involving $\omega$…]

\medskip\noindent
\begin{itemize}
\item Elliptic integrals $\leq 1900$.
\item Poincaré, Pircard, $\sim 1900$.
\item Singular, de Rham 1900-1950.
\item Weil: cohomology over the integers
(enter arithmetic).
\item Grothendieck: étale, crystalline, etc.
\end{itemize}

\noindent
Then is Deligne's dictionary,
in a 1970 paper, where $\ell$-adic
cohomology is introduced
(I think --- abstract in french, check).

\medskip\noindent
From complex to arithmetic geometry:

\begin{itemize}
\item Endilchkeitssätze
fur abelsche Varietäten über Zahlkörpern.
\item 1984: Calculus on arithmetic surfaces.
\item 1988: p-adic Hodge theory.
\item 1990: Degeneration of abelian
varieties (with Chai Ching-li).
\item 1992: Lectures on the
arithmetic Riemann-Roch theorem.
\item 1999: Does there exist an arithmetic
Kodaira-Spencer class? (!!!)
\item 2005: A p-adic Simpson corresponding.
\item 2020: covering of p-adic period
domains.
\end{itemize}

\medskip\noindent
{\bf p-adic Hodge theory.}

Grothendieck and Deligne
noticed that we can obtain
the Hodge decomposition of cohomology
using filtrations, which sit
in the realm of polynomials ---
no harmonic analysis needed.

The idea of Deligne (no, Faltings?) 
was to use
this proof for the $p$-adic case.

\section{Hyperkahler manifolds with 
finite group of birational automorphisms}
\label{section-hk-misha-26}

\noindent
Misha Verbitsky, 20 May, 2026. IMPA.

\medskip\noindent {\bf Abstract.} Seja $M$ uma
coletânea hyperkahler, $b_2(M)\ge 6$, e $D$ seu
espaço de deformação, $I\in D$ uma estrutura
complexa e $H_I$ a rede Hodge de $(M,I)$, ou seja,
o conjunto de todas as classes de cohomologia
inteiras do tipo $(1,1)$. Usando o teorema de
Vinberg sobre redes geradas por reflexões,
provamos que há apenas um número finito de
classes de isomorfismo de redes de Hodge $H_I$
de grau $\ge 3$, de modo que $(M,I)$ seja projetivo
e tenha um grupo finito de automorfismos
birracionais. Usamos essa observação para
mostrar que qualquer família não trivial de
deformações projetivas de $M$ tem um conjunto
denso de estruturas complexas com um grupo
infinito de automorfismos birracionais. Este
é um trabalho conjunto com Ekaterina Amerik e
Andrey Soldatenkov.  


\medskip\noindent
{\bf Upshot.}
The main result we shall prove is the
following.
Let $M$ be a K3 surface and let $\mathcal{H}$
be its polarized moduli space.
There exist finitely many points
$z_1,\dots,z_n \in \mathcal{H}$
such that for any
$m \in \mathcal{H}\setminus \bigcup_{i=1}^n \{z_i\}$
with $\rho(M)\leq 3$, one has
$\lvert \operatorname{Aut}(M)\rvert=\infty$.

In plain language: for most members of the
family, once the Picard rank is small enough,
the surface has infinitely many
automorphisms in ``almost every case'';
only finitely many special
parameter values have infinite automorphism
group.

My favorite idea from the talk is the
following, which I think is called
Noether-Lefschetz result:

\begin{quotation}
  ``Torelli tells us that there are roughly as many
K3 surfaces as Hodge decompositions. Fix a
polarization $p \in H^2(M,\mathbb{Z})$ with
$p^2>0$. There is a group acting on
$\mathbb{P}\operatorname{er}$, namely
$\operatorname{SO}(H^2(M,\mathbb{Z}),p)$, and the
quotient is quasiprojective. This is precisely the
moduli space of polarized K3 surfaces.

\noindent
Now suppose you have a sublattice
$\Lambda \subset H^2(M,\mathbb{Z})$. The
Noether-Lefschetz component is
\[
H_{\Lambda}
\;=\;
\{\,M \in \mathbb{P}\operatorname{er}/\operatorname{SO}
\mid \Lambda \subset H^{1,1}(M)\,\}.
\]
The classes in such quotient correspond
with the K3 that have finite automorphism
group!"
\end{quotation}


\medskip\noindent
The talk starts reviewing several notions
and results like:

\begin{theorem}
A K3 surface $M$ is projective if and only if
there exists a class $\lambda \in H^{1,1}(M,\mathbb{Z})$
such that $\int_M \lambda \wedge \lambda > 0$.
\end{theorem}

\begin{remark}
From the Riemann-Roch formula it follows that
any $(-2)$-class $\lambda$ is either effective or
its negative $-\lambda$ is effective, that is,
equal to the fundamental class of a curve.
\end{remark}

\begin{remark}
This implies that a class in the orthogonal
complement of a $(-2)$-class cannot be Kähler.
\end{remark}

\noindent
The follwing theorem uses the Demailly-Paun
characterization of Kählerness:

\begin{theorem}
Let $M$ be a K3 surface, and $S \subset
H^{1,1}(M,\mathbb{Z})$ the set of all
$(-2)$-classes. Consider the decomposition of
$\operatorname{Pos}(M) \setminus
\bigcup_{s\in S} s^\perp$ into connected
components. Then the Kähler cone of $M$ is one
of the connected components.
\end{theorem}

\begin{theorem}
Let $M$ be a K3 surface. The group
$\operatorname{Aut}(M)$ of holomorphic
isometries of $M$ is naturally identified with
the subgroup of $O(H^2(M,\mathbb{Z}))$ which
preserves the Kähler cone.
\end{theorem}

\begin{definition}
A group is called virtually free abelian if it
contains a free abelian subgroup of finite
index.
\end{definition}

\begin{proposition}[Oguiso]
Let $M$ be a K3 surface, and $H^{1,1}(M,\mathbb{Z})$
its Neron-Severi lattice. Then the following
alternatives occur.
\begin{enumerate}
\item[(a)] If the intersection form on
$H^{1,1}(M,\mathbb{Z})$ is negative definite, then
the kernel of $\Psi: \operatorname{Aut}(M)\to
\operatorname{Aut}(H^{2,0}(M))=\mathbb{C}^*$ is finite,
and the image is virtually free abelian.
\item[(b)] If the intersection form on
$H^{1,1}(M,\mathbb{Z})$ is degenerate, the image and
the kernel of the map $\Psi$ are virtually free
abelian.
\item[(c)] If $M$ is projective, the image of $\Psi$
is finite.
\end{enumerate}
\end{proposition}

\begin{lemma}
Let $M$ be a projective K3 surface with Picard
rank $1$. Then $\operatorname{Aut}(M)$ is finite.
\end{lemma}

\begin{proof}
Indeed, for a projective K3, the map
$\operatorname{Aut}(M)\to O(H^{1,1}(M,\mathbb{Z}))$
has finite kernel (Proposition 1).
\end{proof}

\begin{definition}
Let $(\Lambda,q)$ be a quadratic lattice, that is,
$\mathbb{Z}^n$ equipped with an integer-valued
quadratic form. We say that $(\Lambda,q)$
represents $n\in \mathbb{Z}$ if there exists
$x\in \Lambda$ such that $q(x,x)=n$.
\end{definition}

\begin{lemma}
Let $M$ be a projective K3 surface with Picard
rank $2$. Then $\operatorname{Aut}(M)$ is finite if
its Neron-Severi lattice represents zero, and
contains an infinite cyclic subgroup of finite
index otherwise.
\end{lemma}

\begin{proof}
Since $M$ is projective, $\operatorname{Aut}(M)$ is
finite if and only if its group $\operatorname{Aut}_{\Omega}(M)\subset
\operatorname{SO}(H^{1,1}(M,\mathbb{Z}))$ of symplectic
automorphisms is finite (Proposition 1). The group
of isometries of a quadratic lattice $(\Lambda,q)$
of signature $(1,1)$ is finite if $(\Lambda,q)$
represents zero. It contains an infinite cyclic
subgroup of finite index otherwise, by Pell's
classification of solutions of an equation
$x^2-ay^2=1$.
\end{proof}

\noindent
It looks like a ``divisorial
reflection'' is just a reflection
w.r.t. a  $(-2)$-curve.

\begin{definition}
Let $v \in H^{1,1}(M,\mathbb{Z})$ be a $(-2)$-class.
Consider the map
\[
x \mapsto x-\frac{2(x,v)}{(v,v)}v,
\]
where $(\cdot,\cdot)$ is the intersection form.
This map is called the divisorial reflection.
\end{definition}

\begin{lemma}
The group generated by divisorial reflections acts
transitively on the set of connected components of
\[
\operatorname{Pos}(M)\setminus \bigcup_{s\in S}s^\perp,
\]
where $S$ is the set of all $(-2)$-classes.
\end{lemma}

\noindent
We denote by $\operatorname{Ref}(M)$ the subgroup
of $\operatorname{SO}(H^2(M,\mathbb{Z}))$
generated by divisorial reflections.

\noindent
Clearly, it is a normal subgroup in the group of
all Hodge isometries of $H^2(M,\mathbb{Z})$.

\begin{lemma}
The group $\operatorname{Aut}_{\Omega}(M)$ of
symplectic automorphisms of a K3 surface is
determined, up to an isomorphism, by its
Neron-Severi quadratic lattice.
\end{lemma}

\begin{proof}
It is identified with the subgroup
$\operatorname{St}(\operatorname{Kah}(M))$ of
$\operatorname{SO}(H^{1,1}(M,\mathbb{Z}))$
preserving the Kähler cone. However,
$\operatorname{Ref}(M)$ acts on
$\operatorname{SO}(H^2(M,\mathbb{Z}))$ by Hodge
isometries, transitively on all connected
components of
$\operatorname{Pos}(M)\setminus
\bigcup_{s\in S}s^\perp$.
Therefore, the action of $\operatorname{Ref}(M)$ is
transitive on all groups of form $\operatorname{St}(A)$,
where $A$ is a connected component in
$\operatorname{Pos}(M)\setminus
\bigcup_{s\in S}s^\perp$, and $\operatorname{St}(A)$
its stabilizer.
\end{proof}

\begin{remark}
The group $O(m,n)$, $m,n > 0$, has $4$ connected
components. We denote the connected component of $1$
by $SO^+(m,n)$. We call a vector $v$ positive if its
square is positive.
\end{remark}

\begin{definition}
Let $V$ be a vector space with quadratic form $q$
of signature $(1,n)$, $\operatorname{Pos}(V)=\{x\in
V \mid q(x,x)>0\}$ its positive cone, and
$\mathbb{P}V$ projectivization of $\operatorname{Pos}(V)$.
Denote by $g$ any $SO(V)$-invariant Riemannian
structure on $\mathbb{P}V$. Then $(\mathbb{P}V,g)$ is
called hyperbolic space, and the group $SO^+(V)$ the
group of oriented hyperbolic isometries.
\end{definition}

\begin{example}
For any K3 surface, $\mathbb{P}\operatorname{Pos}(M)$
is a hyperbolic space, and $\operatorname{Aut}(M)$
acts on this space by isometries.
\end{example}

\begin{theorem}[Sterk, 1985]
Let $M$ be a projective K3 surface. Then
$\operatorname{Aut}(M)$ acts with finitely many orbits
on the set of the walls of the Kähler cone.
\end{theorem}

\noindent
The proof may be found in a book by Kneser.

\medskip\noindent
When $M$ is a Calabi-Yau manifold, Sterk's
theorem is known as the Morrison-Kawamata
conjecture, formulated by Morrison in 1993,
and proven by Kawamata for elliptic Calabi-Yau
manifolds in 1997. It was proven by B. Totaro
for log-Calabi-Yau surfaces (2010), and by
Amerik-V. for hyperkähler manifolds (2015).

\begin{remark}
This is an AI-generated explanatory remark.
The ample cone is usually described by ample line
bundles, so its natural generators have integral
first Chern classes. Rational classes appear here
because one often passes to the rational span of
the cone when intersecting with the Kähler cone:
on a projective K3 surface, ample classes are
exactly the rational Kähler classes.
\end{remark}

\begin{lemma}
Let $\operatorname{Amp}(M):=\operatorname{Kah}(M)\cap H^{1,1}(M,\mathbb{Q})$
denote the ample cone of a projective K3
surface, and $\operatorname{Amp}_{\mathbb{R}}(M):=
\operatorname{Kah}(M)\cap H^{1,1}(M,\mathbb{Q})\otimes_{\mathbb{Q}}\mathbb{R}$
denote the corresponding cone in
$H^{1,1}(M,\mathbb{Q})\otimes_{\mathbb{Q}}\mathbb{R}$.
Consider its projectivization
$\mathbb{P}\operatorname{Amp}_{\mathbb{R}}(M)$ as a
subset in the hyperbolic space $\mathbb{P}\operatorname{Pos}(M)$.
Then $\mathbb{P}\operatorname{Amp}_{\mathbb{R}}(M)$
has finite volume if and only if
$\operatorname{Aut}(M)$ is finite.
\end{lemma}

\begin{proof}
By construction,
$\mathbb{P}\operatorname{Amp}_{\mathbb{R}}(M)$ is a
hyperbolic polyhedron (with some vertices on
absolute). If $\operatorname{Aut}(M)$ is finite, this
polyhedron has finite number of faces, hence its
volume is finite. Conversely, if $\operatorname{Aut}(M)$
is infinite, $\mathbb{P}\operatorname{Amp}_{\mathbb{R}}(M)$
contains infinite number of copies of its fundamental
domain, hence its volume is infinite.
\end{proof}

\begin{theorem}[Vinberg]
Let $M$ be a projective K3 surface of
Picard rank $\geq 3$. Then $\operatorname{Aut}(M)$
is finite if and only if the reflection group
$\operatorname{Ref}(M)$ has finite index in
$\operatorname{SO}(H^{1,1}(M,\mathbb{Z}))$.
\end{theorem}

\begin{proof}
Step 1: If $\operatorname{Aut}(M)$ is finite, the
ample cone has finite volume, hence the action of
$\operatorname{Ref}(M)$ on the hyperbolic space
$\mathbb{P}\operatorname{Amp}_{\mathbb{R}}(M)$ has
finite volume. This implies that $\operatorname{Ref}(M)$
is a lattice subgroup in
$\operatorname{SO}(H^{1,1}(M,\mathbb{R}))$, hence has
finite index in $\operatorname{SO}(H^{1,1}(M,\mathbb{Z}))$.

Step 2: Conversely, assume that $\operatorname{Ref}(M)$
has finite index in $\operatorname{SO}(H^{1,1}(M,\mathbb{Z}))$.
Then its Coxeter polyhedron is finite (Vinberg).
However, $\mathbb{P}\operatorname{Amp}_{\mathbb{R}}(M)$
is its Coxeter polyhedron.
\end{proof}

\begin{definition}
\label{definition-reflective}
We call a quadratic lattice with
$\operatorname{Ref}(\Lambda)$ of finite index in
$\operatorname{SO}(\Lambda)$ reflective.
\end{definition}

\noindent
Torelli tells us that there are roughly as many
K3 surfaces as Hodge decompositions. Fix a
polarization $p \in H^2(M,\mathbb{Z})$ with
$p^2>0$. There is a group acting on
$\mathbb{P}\operatorname{Per}$, namely
$\operatorname{SO}(H^2(M,\mathbb{Z}),p)$, and the
quotient is quasiprojective. This is precisely the
moduli space of polarized K3 surfaces.

\noindent
Now suppose you have a sublattice
$\Lambda \subset H^2(M,\mathbb{Z})$. The
Noether-Lefschetz component is
\[
H_{\Lambda}
\;=\;
\{\,M \in \mathbb{P}\operatorname{Per}/\operatorname{SO}
\mid \Lambda \subset H^{1,1}(M)\,\}.
\]
The classes in such quotient correspond
with the K3 that have finite automorphism
group!

\begin{lemma}
\label{lemma-noether-lefschetz-components}
Let $H$ be the moduli space of polarized K3
surfaces. Then there is a finite set of
quasiprojective subvarieties
$Z_1,\dots,Z_n \subset H$ such that for any K3
surface $M$ outside of $\bigcup_i Z_i$ with Picard
rank $\geq 3$, the group $\operatorname{Aut}(M)$ is
infinite.
\end{lemma}

\begin{proof}
Step 1: Fix a polarization on $M$, that is, an
integer vector $p \in H^2(M,\mathbb{Z})$ such that
its square is positive. By Torelli theorem,
$H=\mathbb{P}\operatorname{Per}/\operatorname{SO}(H^2(M,\mathbb{Z}),p)$,
where $\operatorname{Per}$ is the space of all Hodge
structures on $H^2(M)$ such that $p$ has type
(1,1), and $\operatorname{SO}(H^2(M,\mathbb{Z}),p)$
the group of isometries preserving $p$. For any
sublattice $\Lambda \subset H^2(M,\mathbb{Z})$
containing $p$, the corresponding
Noether-Lefschetz component $H_{\Lambda}\subset H$
is the image of all $I \in \operatorname{Per}$ such
that $H^{1,1}(M,I)\subset \Lambda$. By Baily-Borel
theorem, $H_{\Lambda}$ is a quasiprojective
subvariety in $H$.

Step 2: Let $Z_i$ be the set of Noether-Lefschetz
components associated with reflective lattices
$\Lambda \subset H^2(M,\mathbb{Z})$. Clearly, any
$M$ which has Picard rank $\geq 3$ and is not in
$Z_i$ has infinite automorphism group.
\end{proof}

\begin{theorem}[K. Oguiso, 2003]
\label{theorem-oguiso-2003}
Let $M_t$ be a non-trivial quasiprojective family
of K3 surfaces over a curve $C$, with general
member of Picard rank $\geq 3$. Then
$|\operatorname{Aut}(M_t)|=\infty$ for a dense set
of $t \in C$.
\end{theorem}

\begin{proof}
Remove from $C$ all proper subsets which belong to
the Noether-Lefschetz components $Z_i$. If
$C \not\subset Z_j$, we are done. Otherwise, there
is a dense subset of points in $C$ where the Picard
rank jumps, and these points are not in $Z_i$.
\end{proof}

\begin{remark}
\label{remark-oguiso-lower-picard}
This theorem is also valid for families with
general member of Picard rank $\geq 1$, but the
proof for lower Picard values is more elementary.
\end{remark}

\begin{remark}
\label{remark-oguiso-vinberg}
Oguiso's proof does not use Vinberg's theorem, he
uses explicit lattice-theoretic arguments.
\end{remark}

\begin{theorem}[Denisi, Onorati, Rizzo, Viktorova]
\label{theorem-denisi-onorati-rizzo-viktorova}
Let $M_t$ be a non-trivial quasiprojective family
of hyperkähler manifolds of known type over a
curve $C$. Then $|\operatorname{Bir}(M_t)|=\infty$
for a dense set of $t \in C$, where
$\operatorname{Bir}(M_t)$ denotes the group of
birational automorphisms.
\end{theorem}

\begin{remark}
\label{remark-dorv-lattices}
Their proof relies on known properties of the
cohomology lattices of hyperkähler manifolds.
\end{remark}

\begin{theorem}[Amerik, Soldatenkov, V.]
\label{theorem-amerik-soldatenkov-v}
This theorem holds for any family of hyperkähler
manifolds.
\end{theorem}

\begin{remark}
\label{remark-asv-vinberg}
Our proof is based on Vinberg's theorem, in fact,
it is pretty much identical to the proof for K3
surfaces that I gave.
\end{remark}


\section{The unstable flow of a partially 
hyperbolic diffeomorphism with an 
application to measure rigidity}
\label{section-unstable-flow-hyp-dyn}

\noindent
Bruno Santiago (UFF?)

Qui 21 mai 2026, 15:30

\medskip\noindent
{\bf Abstract.}
Let $f$ be a partially hyperbolic
diffeomorphism preserving a splitting
$
TM = E^s \oplus E^c \oplus E^u,
$
where the unstable bundle $E^u$ is
one-dimensional.
When $f$ is homogeneous, like the time-one
map of the geodesic flow on a surface of
constant negative curvature or a linear
automorphism of the torus, the unstable
foliation admits a natural parametrization
that makes invariant measures with
absolutely continuous conditionals
(the so-called $u$-Gibbs measures)
correspond exactly to invariant measures
of this flow.
How can this picture be generalized when
$f$ is not homogeneous?
In this talk, I will explain a construction
of a suitable suspension flow of $f$ for
which the unstable foliation becomes the
orbit foliation of a new flow, called the
unstable flow.
Ergodic $u$-Gibbs measures are in
one-to-one correspondence with measures
simultaneously invariant under both flows
and ergodic for the suspension.
As an application, I will deduce the
following measure rigidity statement:
if $E^c$ is expanding and the bundle
$E^s \oplus E^u$ is integrable, then any
$u$-Gibbs measure whose entropy exceeds its
unstable entropy is necessarily an SRB
measure.
This is joint work with Sébastien
Alvarez, Sylvain Crovisier, Martin
Leguil, and Davi Obata.

\begin{remark}[AI-generated overview]
\label{remark-ai-overview-seminar-unstable-flow}
The talk was about translating a geometric
object into a dynamical one. We start with a
partially hyperbolic diffeomorphism
$$
f:M\to M.
$$
This means that there is a $Df$-invariant
splitting
$$
TM=E^s\oplus E^c\oplus E^u.
$$
Here $Df$-invariant means that
$$
Df_x(E_x^s)=E_{f(x)}^s,\qquad
Df_x(E_x^c)=E_{f(x)}^c,\qquad
Df_x(E_x^u)=E_{f(x)}^u.
$$
Thus $f$ moves points, while $Df$ moves the
corresponding tangent directions in a way
compatible with the splitting.

\medskip\noindent
The key simplifying hypothesis is that
$E^u$ is one-dimensional. Hence the unstable
distribution integrates to curves
$$
W^u(x),
$$
called unstable leaves. The first geometric
idea is that these leaves should not merely
be thought of as a foliation, but as the
orbits of a flow. Thus one wants to construct
an unstable flow $\varphi_t^u$ whose orbit
through $x$ is exactly $W^u(x)$.

\medskip\noindent
The nontrivial point is the choice of
parametrization
[because the passage from the orbits
to a vector field is not integration,
but differentiation, which in principle
is easy --- once you have
the right parametrization!]. 
A one-dimensional leaf can
be parametrized in many different ways, and
the dynamics of $f$ selects a preferred one.
If $x$ and $y$ lie on the same unstable leaf,
the function
$$
\rho_x^u(y)
$$
compares the expansion of $Df^n$ in the
unstable direction at $x$ and at $y$. Thus
$\rho_x^u$ is a dynamical density along
$W^u(x)$.

\medskip\noindent
This density defines a coordinate on the
unstable leaf:
$$
H_x^u(y)
=
\int_x^y
\rho_x^u(z)\,d\operatorname{Leb}_x^u(z).
$$
The inverse map
$$
\Phi_x^u=(H_x^u)^{-1}
$$
then gives the natural parametrization of the
unstable leaf. This is the geometric heart of
the construction.

\medskip\noindent
The measure-theoretic part enters because in
dynamics one studies not only where an orbit
goes, but also its long-time statistical
distribution. For a point $x$, the empirical
measures
$$
\mu_{n,x}
=
\frac{1}{n}
\sum_{j=0}^{n-1}
\delta_{f^j(x)}
$$
record the proportion of time the orbit of
$x$ spends in different regions of $M$.
Equivalently, for a measurable subset
$A\subset M$, one has
$$
\mu_{n,x}(A)
=
\frac{1}{n}
\#\{0\leq j<n:f^j(x)\in A\}.
$$
Thus $\mu_{n,x}(A)$ is the proportion of
time, up to time $n$, that the orbit of $x$
spends inside $A$. This is where time enters:
we are not measuring the size of $A$ itself,
but rather how frequently the forward orbit
$$
x,\ f(x),\ f^2(x),\ldots
$$
visits $A$.

\medskip\noindent
A $u$-Gibbs measure is an invariant measure
which is smooth in the unstable direction.
More precisely, if
$$
\{\mu_x^u\}
$$
is the disintegration of $\mu$ along unstable
leaves, then
$$
\mu_x^u\ll \operatorname{Leb}_x^u.
$$
So the measure may be globally complicated,
but leaf-by-leaf along $W^u$ it has a density
with respect to ordinary leafwise length.

\medskip\noindent
The main theorem of the talk says that such
$u$-Gibbs measures can be characterized after
passing to a suspension space. The original
map $f$ is discrete-time, while the suspension
turns it into a continuous-time flow
$\psi_t$. The theorem identifies $u$-Gibbs
measures with lifted measures which are
invariant both under the suspension flow
$\psi_t$ and under the unstable flow
$\varphi_t^u$.

\medskip\noindent
Finally, the application concerns SRB
measures. An SRB measure is stronger than a
$u$-Gibbs measure: it is meant to describe
the statistics seen by a positive-Lebesgue
set of initial conditions. In the theorem
from the end of the talk, if the center
direction is expanding and a $u$-Gibbs measure
has more entropy than what is accounted for
only by the unstable direction, namely
$$
h_\mu(f)>
\int_M \log \lambda_x^u\,d\mu(x),
$$
then the measure must be SRB.

\medskip\noindent
In the language of the Heap Project, this is
not like constructing a scheme from equations;
it is more like extracting an invariant
object from asymptotic data. Just as a Hilbert
polynomial records asymptotic information
about graded pieces, an invariant measure
records asymptotic information about the
distribution of an orbit. The slogan of the
talk is therefore:
$$
\text{unstable expansion}
\quad\leadsto\quad
\text{dynamical density}
\quad\leadsto\quad
\text{unstable flow}
\quad\leadsto\quad
\text{measure rigidity}.
$$
\end{remark}

\noindent
Let $S:M \to M$ be a $C^r(M)$ map.
Let $V$ be a vector field
inducing a foliation
whose leaves form a $C^r$ foliation
on $M$.
The idea is
that this dimension-1 
foliation is dilated by the
diffeomorphism.

\begin{example}
\label{example-sl3-spectrum}
\begin{itemize}
\item $A \in \text{SL}_3(\mathbb{Z})$,
with spectrum
$\text{spec}(A)
= \{0<\lambda_A^s<1<
\lambda^c_A
<\lambda^u_A\}$.
$W_A^u
= $ folitation with orbits
of irrational flux $\varphi_t$
of $f_A:\mathbb{T}^3 \to \mathbb{T}^3$.
\end{itemize}
\end{example}

\begin{definition}[AI]
\label{definition-u-gibbs-measure}
Let $f:M \to M$ be a partially hyperbolic
diffeomorphism, and let
$\mathcal F^u$ be its unstable foliation.
A measure
$$
\mu \in \mathcal P_f^{\operatorname{reg}}(u)
$$
is called a $u$-Gibbs measure if its
conditional measures along unstable leaves
are absolutely continuous with respect to
Lebesgue measure on the leaves.

More precisely, if
$$
\{\mu_x^u\}_{x \in M}
$$
denotes the disintegration of $\mu$ along
the unstable foliation, then
$$
\mu_x^u \ll \operatorname{Leb}_x^u
$$
for $\mu$-almost every $x \in M$.
\end{definition}

\begin{remark}
\label{remark-stable-directions-dyn}
$$
TM = E^s \oplus E^c \oplus E^u.
$$
\end{remark}

\begin{remark}
\label{remark-leb-ae-ai}
For Lebesgue almost everywhere,
$x \in M$, $n_k \to \infty$,
$$
\mu_{n_k,x}
:=
1 \sum_{n_k=0}^{n_k-1}
\delta_{f_k} \to \mu
$$
which implies that $\mu \in 
\text{Leb}^u$.
\end{remark}

\noindent
I'm very lost, here's AI-corrected
above remark:

\begin{remark}
\label{remark-lebesgue-ae}
The phrase ``Lebesgue almost every $x$''
means almost every point with respect to
the Riemannian volume measure on $M$.

Given $x \in M$, its empirical measures are
$$
\mu_{n,x}
:=
\frac{1}{n}
\sum_{j=0}^{n-1}
\delta_{f^j(x)}.
$$
A measure $\mu$ is called physical if its
basin
$$
B(\mu)
:=
\left\{
x \in M :
\mu_{n,x} \to \mu
\right\}
$$
has positive Lebesgue measure, where the
convergence is weak-$*$ convergence of
probability measures.
\end{remark}

\begin{lemma}[Ledrappier density]
\label{lemma-ledrappier}
$$
\frac{d\mu^u_x}{
d\text{Leb}_x^u}
=c(X)\rho_x^u
$$
where $\rho_x^u:w^u(x)\to \mathbb{R}$,
$$
\rho_x^u(y)
=
\lim_{n \to \infty} 
\frac{\|D\bar{f}^n(y)V^u\|}{
D\bar{f}^n(x) V^u}
$$
is the dynamic density function.
\end{lemma}
\begin{lemma}[Ledrappier density, AI]
\label{lemma-ledrappier-density}
Let $\mu$ be a $u$-Gibbs measure. Then the
conditional measure $\mu_x^u$ is absolutely
continuous with respect to $\operatorname{Leb}_x^u$,
and its density is of the form
$$
\frac{d\mu_x^u}{d\operatorname{Leb}_x^u}(y)
=
c(x)\rho_x^u(y),
$$
where $c(x)>0$ is a normalizing constant and
$$
\rho_x^u:W^u(x)\to \mathbb R
$$
is given by
$$
\rho_x^u(y)
=
\lim_{n\to\infty}
\frac{
\|D\bar f^{-n}(y)|_{E_y^u}\|
}{
\|D\bar f^{-n}(x)|_{E_x^u}\|
}.
$$
Equivalently, in the one-dimensional unstable
case,
$$
\rho_x^u(y)
=
\lim_{n\to\infty}
\frac{
\|D\bar f^n(x)v_x^u\|
}{
\|D\bar f^n(y)v_y^u\|
},
$$
where $v_x^u\in E_x^u$ and $v_y^u\in E_y^u$
are unit vectors.
\end{lemma}

\medskip\noindent
{\bf Properties of EDD.}
\begin{itemize}
\item 
$$
\rho_y^(u)
=\frac{j}{\rho^u_x(y)}
$$
\item 
$$
z,y \in W^u(x)
\implies 
\rho^u_x(z)
=\rho_x^u(y)\rho_y^u(z)
$$
\item 
$$
\rho^u_{f(x)}(f(y))
=\frac{\lambda_x^u}{\lambda_y^u}
\rho_x^u(y).
$$
\end{itemize}

\medskip\noindent
{\bf Affine structure.} 
$$
H_x^u:W^u(x) \to \mathbb{R}
$$
$$
H_x^u(y)
=\int_x^y \rho_x^u(z)
d\text{Leb}_x^u(z)
$$
$$
\Phi_x^\mu
:=(H_x^u)^{-1}
:\mathbb{R} \to W^u(x).
$$

Change of coordinates:
$$
H_y^u \circ\Phi_x^u(t)
=\rho_y^u(x)H
+H_y^u(x).
$$
\begin{definition}[Unstable flux]
\label{definition-unstable-flux}
$$
\begin{aligned}
\tilde{\varphi}_t:
M \times \mathbb{R} &\longrightarrow 
M \times \mathbb{R} \\
(x,\alpha) &\longmapsto 
(\Phi^u_x(e^{-x}t),
\alpha
+\text{log}\rho_x^u\Phi_t^u(e^{-\alpha}t))
\end{aligned}
$$
\end{definition}

\noindent
Let's check why this, in fact,
defines a flux:

\begin{lemma}
\label{lemma-defines-a-flux}
$$
\tilde{\varphi}_t \circ \tilde{\varphi}_s
(x,\alpha)=\tilde{\varphi}_{t+s}(x,\alpha)
$$
\end{lemma}

\begin{proof}
Provided
\end{proof}

\begin{theorem}
\label{theorem-gibbs-measure}
$\mu \in \text{Gibbs}_{\text{reg}}^u(f)$,
then there exists a unique
$$
\hat{\mu} 
\in \mathcal{P}_{\psi}^{\text{reg}}
(\hat{\mu}) \cap
\mathcal{P}_\varphi(\hat{\mu})
$$
\end{theorem}

$$
D(f)
=\{(x,\alpha):0\leq \alpha<\lambda_x^u\}
\subset\hat{\mu}
$$
$$
\hat{\mu}
=\frac{1}{
\int_M \text{log} \lambda_x^u
d\mu(x)}\mu\otimes\text{Leb}_{\mathbb{R}}
(\mathbb{T}^{-x}(A)).
$$

\medskip\noindent
{\bf Applications of Theorem
\ref{theorem-gibbs-measure}}

\noindent
Back to the totally partially hyperbolic
context, suppose there is a decomposition
$$
TM^3 = E^s \oplus E^c \oplus E^u
$$
such that
$$
0 \leq \|Df(x)|_{E^s}\|
<I<\|Df(x)|_{E^c}\|
<\|Df(x)|_{E^u}\|.
$$
Then $f$ is Anosov partially hyperbolic.
Thus, there exist foliations
$W^s,W^c,W^u$ such that
$E^s,E^c$ and $W^u$ and
$W^{cs}, W^{du}$
$$
F^s \oplus E^c \oplus E^c \oplus E^u.
$$

\begin{definition}
\label{definition-}
$\mu \in \mathcal{P}_f(u)$
is SRB if all the conditionals
$$
\{\mu_x^{c\mu}\}_{x \in M}
$$
along $W^{c\mu}(\alpha)$
satisfy
$$
\mu_x^{c\mu}\ll \text{Leb}_x^{c\mu}.
$$
\end{definition}

\begin{theorem}[Sinoi-Rulh-Bouso]
\label{theorem-srb}
There exists a unique SRB measure
$\mu$.
\end{theorem}

\noindent
Conjecture (ALOS, Katz): 
if $E^s \oplus E^u$
are not integrable then any
$$
\mu \in \text{Gibbs}_{\text{reg}}^{\mu}(f)
$$
is SRB.

\begin{theorem}[B]
\label{theorem-gibbs-b}
If $E^s \oplus E^u$ is integrable
and
$$
\mu \in \text{Gibbs}_{\text{reg}}^u(u)
$$
satisfies
$$
h_\mu(f)
> \int \text{log}\lambda_x^u d\mu(x),
$$
then $\mu$ is SRB.
\end{theorem}

\medskip\noindent
{\bf Some ideas.}

 \noindent
$x \in M \mapsto \text{log}\lambda_x^u$
is $C^{1+\alpha}$.

Descição das $W^*$
$$
\widetilde{W}^s(x,\alpha)
=\{(x,\alpha+\theta_x^{s}(y)
:y \in W^s(x)\}
$$
where
$\theta(\alpha)=\text{log}\lambda_x^u$
and
$\theta_x^s(y)
=\sum_{\ell=0}^{\infty}
\theta(f^\ell y)
-\theta(f^\ell x)$.



\section{Weil--Petersson homeomorphisms 
and maximal surfaces in anti-de Sitter space}
\label{section-weil-petersson}

\noindent
Graham Smith

May 21, 2026.

\medskip\noindent
{\bf Abstract.}
A classe de Weil--Petersson é uma classe
notável de homeomorfismos do círculo que se
encontra na fronteira de vários campos da
matemática contemporânea, incluindo a teoria
de Teichm\"uller, a teoria de
Schramm--Loewner, a geometria hiperbólica e
assim por diante.
Neste trabalho, apresentamos uma nova
perspectiva, estudando essa classe do ponto
de vista da geometria anti-de Sitter.
Mostramos que um homeomorfismo do círculo é
Weil--Petersson se e somente se seu gráfico
limitar uma superfície máxima completa,
semelhante ao espaço, em
$\mathrm{AdS}^{2,1}$ de área finita
renormalizada.

Este é um trabalho conjunto com F. Diaf,
A. Moriani, R. Smai e E. Trebeschi.
2604.17804.

\medskip\noindent
We are interested in Weil-Petersson
Teichmüller space $\mathcal{T}_0$
This is a very interesting object
that intersects several branches of 
mathematics:
\begin{itemize}
\item Renormalized area of minimal surfaces
in $\mathbb{H}^3$.
\item Stochastic Loewner evolution.
\item Geometric computer vision.
\item Berezin Quantization and Bott-Virasoro
group.
\item Planar Coulomb gasses.
\end{itemize}

\medskip\noindent
{\bf Loewner transport.}
Let  $\Omega \subset \hat{\mathbb{C}}$
be a Jordan domain (interior of Jordan
curve).
Let $p,q \in \Omega$.
For example, $\Omega=\mathbb{H}^2$,
$p=0$ and $q=\infty$.
Let $\gamma:[0,1]\hookrightarrow
\overline{\mathbb{H}^+}$
with $\gamma(0)=0$ and $\gamma(t) \in
\mathbb{H}^+$ for all $t>0$.

For all $t$ define
$$
\Omega_t:\mathbb{H}^+\setminus
\gamma([0,t])
$$
$$
g_t:\mathbb{H}^t \to \Omega_t
$$
be the unique uniform map
$$
g_t(z)
=z + \frac{2\varphi(t)}{z}
+ o \left(\frac{1}{z}\right)
$$
where $\varphi$ is real valued and
increasing. Parametrize $\gamma$
by $\varphi$:
$$
g_t(z)=z+ \frac{2t}{z}
+o \left(\frac{1}{z}\right).
$$
Caratheory:
$$
W(t)=g^{-1}_t(\gamma(t))
$$
$$
W:\underbrace{[0,L]}_{\text{domain of }
\gamma} \to \mathbb{R}.
$$

The {\it Lowener transform} (1923)
is the map
$$
\gamma \mapsto W
$$
and $W$ is the {\it driving function},
invertible.

\begin{remark}
\label{remark-brownian-loewner}
If $W$ is Brownian we obtain a 
stochastic loewner evolution.
Yilin Wane.
\end{remark}

\medskip\noindent
{\it Loewner energy:}
$$
I(\gamma)
:= \frac{1}{2}
\int_0^L W'(t)^2 dt
$$
\begin{remark}
\label{remark-loewner-energy-independent}
Jordan curves have a well-defined
Loewner enery
[some remark I didn't understand concerning
$\Gamma$ Jordan curve, $p,q \in \Gamma$…]
\end{remark}

\noindent
In this talk we shall take this
as a definition:

\begin{theorem}[Wang]
\label{theorem-wang}
If $\Gamma \subseteq \hat{ \mathbb{C}}$
is a Jordan curve then
$I(\Gamma) < \infty$ if and only if
$\Gamma$ is in $\mathcal{T}_0$,
the {\it Weil-Petersson Teichmüller space}.
\end{theorem}

\medskip\noindent
{\bf Universal Teichmüller space.}
Let
$B_1^{\infty}$ be the open unit ball in
$L^\infty(\mathbb{H}^+)$.
Let $\mu \in B_1^\infty$,
$$
\mu_0(z)=
\begin{cases}
\mu(z)\qquad &\text{ in }\mathbb{H}^+ \\
0\qquad &\text{ in }\mathbb{H}^-.
\end{cases}
$$

The upshot of the following theorem
is that these functions behave
like holomorphic functions:

\begin{theorem}[Measurable Riemann Mapping
Theorem]
\label{theorem-measurable-rmt}
There exists a unique $f_\mu:\hat{\mathbb{C}}
\to \hat{\mathbb{C}}$ such that
\begin{enumerate}
\item $f_\mu(0)$, $f_\mu(1)=1$,
 $f_\mu(\infty)=\infty$.
\item 
$$
\frac{f_{\mu,\bar{z}}}{
f_{\mu z}}=\mu_0.
$$
\end{enumerate}
\end{theorem}

\noindent
Recall:

\begin{definition}
\label{definition-quasiconformal-map}
$f:\Omega(\subseteq\mathbb{C})
\to \mathbb{C}$
is {\it quasiconformal} if
\begin{enumerate}
\item $f$ is a local homeomorphism,
\item $f$ has locally $L^2$ 
dist. derivatives.
\item There exists an $L^\infty$
function $\mu$ such that
$\|\mu\|_{L^\infty}<1$ such that
$$
f_{\bar{z}}=f_z\cdot \mu
$$
and $\mu$ is the {\it Beltrami differential} 
of $f$.
\end{enumerate}
\end{definition}

\noindent
$\mu_0$ is the ``extension by $0$''.
There's also:
$$
\mu_s(z)
=\begin{cases}
\mu(z)\qquad &z \in \mathbb{H}^+ \\
\overline{\mu(\bar{z})}\qquad &
z \in \mathbb{H}^-
\end{cases}
$$
and by the Measurable Riemann Mapping 
Theorem, there exists a unique
$f^u:\hat{\mathbb{C}} \to \hat{\mathbb{C}}$
such that $\sim$
$$
f_{\bar{z}}^\mu
=\mu_s f^\mu.
$$

[Picture which I think
says $f_\mu$ maps the upper half plane
to a Jordan region]

\begin{definition}
\label{definition-quasicircle}
A Jordan curve $\Gamma$ is a
{\it quasicircle} whenever
$$
\Gamma=f_\mu(\hat{\mathbb{R}})
$$
\end{definition}

\noindent
Symmetry:
[Picture showing $f^\mu$
putting the upper half plane
to the whole plane.]

$f^\mu$ restricts to a homeomorphism
of $\hat{\mathbb{R}}$,
$$
\varphi:\hat{\mathbb{R}}
\to \hat{\mathbb{R}}
$$
is quasisymmetic whenever
$$
\varphi=f^\mu|_{\hat{\mathbb{R}}}
$$
for some $\mu$.

$$
\xymatrix{
& B_1\subseteq L^\infty
(\mathbb{H}^+)\ar[dl]\ar[dr]\\
\text{Quasicircles}
\ar[rr]_{\text{Homeo}}
&&
\text{Quasisymmetries}
}
$$
where the labelled homeomorpihsm
is the {\it conformal Welding}.
If $\mu_1, \mu_2 \in B_1$, then
$$
\text{Im}(f_{\mu_1})
= \text{Im}(f_{\mu_2})
\iff
f^{\mu_1}|_{\hat{\mathbb{R}}}
=f^{\mu_2}|_{\hat{\mathbb{R}}}
$$

Bers was very interested in this triangle,
and introduced what he called
``universal Teichmüller space''.

$$
\begin{aligned}
\mathcal{T}&=\text{QCirc}/
\text{PSL}(2,\mathbb{C})\\
&=\text{QSymm}/\text{PSL}(2,\mathbb{R})
\times \text{PSL}(2,\mathbb{R})\\
&=B^\infty_1/\sim.
\end{aligned}
$$

Bers showed that
\begin{itemize}
\item $\mathcal{T}$ is a Banach
manifold.
\item Every $TM$ space embeds smoothly
in $\mathcal{T}$.
\end{itemize}

\noindent
Every Teichmüller space has a 
Weil-Petersson metric. Does $\mathcal{T}$
carry a metric which is induced
by a Weil-Petersson metric? No.

Takhtajan-Teo ($\sim$2004):
$\mathcal{T}$ carries a Hilbert
manifold structure $\mathcal{T}_1$,
it has $\mathbb{R}$ connected components,
contains every Teichmüller space
and its metric (induced by the
Hilbert manifold structure)
induces a Weil-Petersson
metric.

\begin{definition}
\label{definition-wp-teichmuller}
$\mathcal{T}_0 \subseteq T_1$
is the connected component of $I_0$.
\end{definition}

\medskip\noindent
{\bf Properties.}

\begin{itemize}
\item $\mathcal{T}_0$ has a
unique homogeneous Kähler metric.
\item Kähler potential is $I$.
\item Among the dozens
of characterizations of this space,
we have:
$\Gamma$ is $\mathcal{T}$
if and only if $\Gamma$
is chord arch and arc-length
parametrization is $H^{3/2}$,
there exists $c$ such that
$$
d_{\text{ext}}(p,q)
\leq  d_{\text{int}}(p,q)
\leq cd_{\text{ext}}(p,q).
$$
\end{itemize}

\noindent
Here's another characterization 
of this space:

\begin{theorem}[Bishop]
\label{theorem-bishop-teichmmuller}
$\Phi \in \mathcal{T}_0$ if and only if
$\Gamma$ bounds a minimal surface
$\Sigma$ in $\mathbb{H}^3$
of finite topology and finite
renormalized area, which
is the integral of the determinant
of the second fundamental form
(and can be infinite).
\end{theorem}

\medskip\noindent
{\bf Reminder on anti de Sitter space.}

$$
\text{AdS}^{2,1}
=
\{(x_1,…,x_4)
: x_1^2+x_2^2-x_3^2-x_4^2=-1\}
/\{\pm  1\}.
$$
Also, define for any $M$
$$
\text{Det}(M):=
-\frac{1}{2}
\text{Tr}(J_0M J_0M^t)
$$
for
$$
J_0=
\begin{pmatrix}0&-1\\ 
1&0\end{pmatrix}.
$$
Then
$$
\begin{aligned}
\text{AdS}^{2,1}&=
\{M : \text{Det}M=0, M\neq 0\}
/\{\pm I_0\}
\end{aligned}
$$

\noindent
and
$$
\begin{aligned}
\delta_{\infty}\text{AdS}^{2,1}&=
\{M: \text{Det}(M)=0,M \neq 0\}\\
&=\{M:\text{rk}M=1\}/\mathbb{R}^*\\
&=\underbrace{\mathbb{R}P^1}_{\text{image}} 
\times \underbrace{\mathbb{R}P^1}
_{\text{kernel}}.
\end{aligned}
$$

\noindent
A homeomorphism 
$\varphi:\mathbb{R}P^1 \to \mathbb{R}P^1$
gives a graph $\Lambda_\varphi
\subset \partial_{\infty}\text{AdS}
^{2,1}$ which is spacelike.

Anti de Sitter space is a lot nicer
than hyperbolic space:

\begin{theorem}[Labourie-Toul-Wolf]
\label{theorem-lab-tou-wol}
Every spacelike curve $\Lambda$
in $\partial_{\infty}\text{AdS}^{2,1}$
bounds a unique complete maximal disk.
\end{theorem}

\begin{theorem}[Graham-et al.]
\label{theorem-graham-teich}
$\varphi$ is WP if and only if
$D_\varphi$ has finite renormalized area.
\end{theorem}

\section{On multiplicities in length spectra of semi-arithmetic surfaces}
\label{section-multiplicities-length-spectra-semi-arithmetic-surfaces}

\noindent
Misha Belolipetsky.

Ter 26 mai 2026, 15:30.

SALA 236.

\medskip\noindent 
{\bf Abstract.} Mostramos
que determinadas superfícies semi-aritméticas
de Riemann têm um crescimento exponencial das
multiplicidades médias em seu espectro de
comprimento. Anteriormente, o crescimento
exponencial das multiplicidades médias era
conhecido apenas para superfícies
aritméticas, embora os resultados
experimentais que indicavam a possibilidade
de exemplos não aritméticos já estivessem
disponíveis. Explicarei alguns detalhes
interessantes de nossa prova e indicarei
algumas questões em aberto relacionadas. A
palestra é baseada em um trabalho conjunto
recente com Gregory Cosac, Cayo Dória e
Gisele Teixeira Paula.

\medskip\noindent
Zoll metrics $\leftrightarrow$
systole has large mutiplicity.

\medskip\noindent
Our case $\leftrightarrow$
``large geodesics have
very large multiplicity''.

\medskip\noindent
Plan of the talk:
\begin{enumerate}
\item Introduction to semiarithmetic
surfaces.
\item 
\end{enumerate}


\medskip\noindent
Let $S=\mathbb{H}^2/\Gamma$
be a Riemann surface, i.e.
the hyperbolic plane modulo a 
Fuschian group, which is a finite discrete
subgroup of $\text{PSL}(2,\mathbb{R})$
of finite coarea.
We introduce the number $\mathcal{N}(\ell)$,
the number of closed geodesics on $S$
of length $\leq \ell$.

\begin{theorem}[Huber, Selberg]
\label{theorem-huber-selberg}
$\mathcal{N}(\ell)
\sim \frac{\text{exp}\ell}{\ell}$.
\end{theorem}

\noindent
This number is well-defined: i.e.
there are only finitely many closed
geodesics in a hyperbolic surface.

\begin{remark}
\label{remark-nlog}
$\mathcal{N}(\text{log} x)
\sim \frac{x}{\text{log} x}$.
\end{remark}

\noindent
The length set is
$$
\{0<\ell_1<\ell_2<…\}
$$
Also we denote
$$
\mathcal{N}'(\ell)
=\#\{n : \ell_n <\ell\}
$$
and $g(\ell)$ is the multiplicity of
$\ell_n$.

\begin{definition}
\label{definition-mean-multiplicity}
The {\it mean multiplicity} 
$\langle g(\ell)\rangle$ is a continuous
monotone function
$$
\mathcal{N}(\ell)
=\sum_{\ell_n<\ell}g(\ell_n)
\sim
\sum_{\ell_n \leq  \ell}
\langle f(\ell)\rangle
$$
\end{definition}

$$
\langle g(\ell)\rangle
\sim
\frac{\mathcal{N}(\ell)}{\mathcal{N}'
(\ell)}.
$$
\medskip\noindent
In a hyperbolic surface
$S = \mathbb{H}^2/\Gamma$,
closed geodesics correspond to
hyperbolic elements $\gamma \in \Gamma$,
$|\text{tr}\gamma|>2$.
Similarly, lengths correspond to traces.
We have
$$
\ell(\gamma)
=2 \text{arcosh}
\left(\frac{|\text{tr}\gamma|}{2}\right).
$$
$$
\mathcal{L}(\Gamma)
=\{|\text{tr}(\gamma)|:\gamma \in \Gamma,
\gamma\text{ hyp.}\}
$$
\begin{example}
\label{example-modular-group}
$\Gamma = \text{PSL}_2(\mathbb{Z})$,
$\mathcal{L}(\Gamma) \subset \mathbb{N}$
(in particular, discrete in $\mathbb{R}$).
\end{example}

$$
\#(\mathcal{L}(\Gamma)
\cap[2,T])<T.
$$
Each length $\ell$
produces a unique trace
$\leq  2 \text{cosh}\left(\frac{\ell}{2}
\right)$.
$$
\mathcal{N}'(\ell)
\leq 2 \text{cosh}(\ell/2)\leq 2e^{\ell/2},
$$
$$
\langle g(\ell)\rangle
\sim\frac{\mathcal{N}(\ell)}{
\mathcal{N}'(\ell)}
\geq \frac{\text{exp}(\ell/2)}{2\ell}
$$

\begin{theorem}[Luo-Sernak (1994)]
\label{theorem-luo-sarnak}
Instead of discretness
of $\mathcal{L}(\Gamma)$ in $\mathbb{R}$,
it is enough
$$
\# (\mathcal{L}(\Gamma)
\cap [N-1,N]
\leq \text{const.}
$$
for any $N \in \mathbb{N}$.
\end{theorem}

\noindent
This is called the
{\it bounded clustering property (BC)}.

\begin{example}
\label{example-triangle24oo}
$\Gamma=\Delta(2,4,\infty)$.
It is well-known this group is generated
by
$$
\begin{pmatrix}0&1\\ -1&0\end{pmatrix},
\qquad 
\begin{pmatrix}1&\sqrt{2}\\ 0&1\end{pmatrix}.
$$
$\mathcal{L}(\Gamma) \supset 2 \mathbb{N}
\cup \sqrt{2}\mathbb{N}$.
Non discrete in $\mathbb{R}$,
but satisfies B. C.
\end{example}

\begin{remark}
\label{remark-bc-holds}
BC holds for ``some''
arithmetic groups.
\end{remark}

\begin{itemize}
\item BC $\implies $ exponential group
of $\langle g(\ell)\rangle$.
\item All arithmetic Fuchsian groups have 
BC.
\item Conjecture: BC iff group is $\Gamma$
is arithmetic.
\item Non-compact case: Gerinska-Leuzinger
(2008).
\item Open in general.
\end{itemize}

\medskip\noindent
Let's try to define arithmetic 
and semi-arithmatic groups.

Let $\Gamma \subset \text{SL}_2(\mathbb{R})$
be a Fuschian group.
Consider the {\it invariant trace field}.
It's the field generated by the traces
of the elements in the group:
$$
K:=\mathbb{Q}\bigl(\{\text{tr}(\gamma^2):
\gamma \in \Gamma\}\bigr).
$$

Recall that a totally real field
is such that all complex embeddings are real.
More precisely:

\begin{definition}
\label{definition-totally-real-field}
A number field $K$ is {\it totally real} if 
every field embedding
$K \hookrightarrow \mathbb{C}$ has image 
contained in $\mathbb{R}$.
\end{definition}

\begin{definition}
\label{definition-semiarithmetic}
Let $\Gamma$ be a Fuchsian group
of finite covolume. $\Gamma$ is called
{\it semiarithmetic} is $K$ is totally
real and
$\text{tr}(\gamma) \in \mathcal{O}_K$
--- the integers of $K$.
\end{definition}

$$
\Gamma^{(2)}
\subset \text{SL}_2(\mathcal{O}_K)
\underbrace{\subset}_{\text{lattice}}
\text{SL}_2(K \otimes_\mathbb{Q}\mathbb{R})
=\text{SL}_2(\mathbb{R})^2
\times 
\underbrace{\text{SO}(3)^{[K:\mathbb{Q}]-r}}_{
\text{compact}}
$$
That is $\text{SL}_2(\mathbb{R})^r$
is a noncompact semisimple piece,
and $\text{SO}(3)^{[K:\mathbb{Q}]-r}$
is compact. Then we take projection,
which ammounts to ``modding out the
noncompact part''.

  $$
  \xymatrix{
  \mathrm{SL}_2(K \otimes_{\mathbb{Q}}
  \mathbb{R}) \ar[r]^-{\pi_{\infty_1}}
  \ar[dr]_{\pi_{\infty_2}} &
  \mathrm{SL}_2(\mathbb{R})^r
\times\text{SO}(3)^{[K:\mathbb{Q}]-r} \\
  & \mathrm{SL}_2(\mathbb{R})
  }
  $$

Now for $r = 1$ corresponds with 
$\Gamma$ being arithmetic.
For  $r>1$, is semi-arithmetic,
$r$ is the arithmetic dimension.
More precisely,

\begin{definition}
\label{definition-arithmetic-group}
Let $\Gamma$ be a Fuchsian group of finite
covolume, and let $K$ be its invariant trace
field. We say that $\Gamma$ is
{\it arithmetic} if $\Gamma^{(2)}$ is
commensurable with $\mathrm{SL}_2(\mathcal{O}_K)$.
\end{definition}

What is happening here is that one starts
with a group of matrices over the number
field $K$ and then views it inside the much
larger real Lie group
$\mathrm{SL}_2(K \otimes_\mathbb{Q}\mathbb{R})$.
After extending scalars, this ambient group
splits into one or more noncompact
$\mathrm{SL}_2(\mathbb{R})$ factors and a
product of compact factors. The projection
diagram is the way of choosing the real
factor that gives the actual action on the
hyperbolic plane, while the other factors are
auxiliary arithmetic data. In the case
$r=1$, there is only one noncompact factor,
and that is exactly the arithmetic case;
when $r>1$, the same construction gives the
semi-arithmetic situation.

\medskip\noindent
Now,
$$
\xymatrix{
\{|\text{tr}\gamma|:\gamma \in \Gamma^{(2)}\}
\ar[r]^{\sigma_i}\ar[dr]&  \text{unbounded}\\
&C[0,2]
}
$$
satisfies BC by the pigeon hole.

\begin{example}[Golden triangle]
\label{example-golden-triangle}
$$
\Gamma = \Delta(2,5,\infty)
=\langle \begin{pmatrix}0&-1\\
-1&0\end{pmatrix},
\begin{pmatrix}1&\gamma\\ 0&1\end{pmatrix}
\rangle,
$$
where $\gamma=\frac{1+\sqrt{5}}{2}$.
Then
$$
\Gamma\underbrace{\subset}_{\text{thin}}
\text{SL}_2(\mathbb{Z}[\sqrt{5}])
\underbrace{\subset}_{\text{arithm.}}
\text{SL}_2(\mathbb{R})^2.
$$
Here $K=\mathbb{Q}(\sqrt{5})$,
$\sigma_1=\text{id}:K \hookrightarrow 
\mathbb{R}$,
$\sigma_2:\sqrt{5} \to -\sqrt{5}$.
$\sigma_1(\Gamma) = \Gamma
\subset\text{SL}_2(\mathbb{R})$ - Fuschian
subgroup.
$\sigma_2(\Gamma)\subset
\text{SL}_2(\mathbb{R})$ - not discrete.
\end{example}

\noindent
And that's it! Since this was particularly
difficult to grasp, let us
include an AI-summary:

\medskip\noindent
{\bf On multiplicities in length spectra of
semiarithmetic surfaces.}
This was a talk by Misha Belolipetsky on
joint work with Gregory Cosac, Cayo Dória
and Gisele Teixeira Paula. The main result
is that certain semiarithmetic Riemann
surfaces have exponential growth of the
mean multiplicities in their length spectra.
Previously this phenomenon was known for
arithmetic surfaces; the point of the talk
was that it also occurs in some genuinely
semiarithmetic examples.

\medskip\noindent
The geometric object is a hyperbolic surface
$$
S=\mathbb{H}^2/\Gamma,
$$
where $\Gamma$ is a Fuchsian group, i.e. a
discrete subgroup of
$\operatorname{PSL}_2(\mathbb{R})$ of finite
coarea. We write $\mathcal{N}(\ell)$ for the
number of closed geodesics on $S$ of length
at most $\ell$. The Huber--Selberg theorem
says
$$
\mathcal{N}(\ell)\sim \frac{e^\ell}{\ell}.
$$
Thus the total number of closed geodesics
grows exponentially. The question of the talk
is subtler: if we list only the distinct
lengths
$$
0<\ell_1<\ell_2<\cdots,
$$
and write
$$
\mathcal{N}'(\ell)=\#\{n:\ell_n<\ell\},
$$
then how large is the average multiplicity
of a length? If $g(\ell_n)$ denotes the
number of closed geodesics with length
$\ell_n$, then morally
$$
\langle g(\ell)\rangle
\sim
\frac{\mathcal{N}(\ell)}{\mathcal{N}'(\ell)}.
$$
So the problem becomes: there are very many
closed geodesics, but how many different
lengths do they produce?

\medskip\noindent
The bridge from geometry to arithmetic is
given by traces. Closed geodesics correspond
to hyperbolic conjugacy classes in
$\Gamma$. If $\gamma\in\Gamma$ is
hyperbolic, then $|\operatorname{tr}\gamma|>2$
and
$$
\ell(\gamma)
=
2\operatorname{arccosh}
\left(\frac{|\operatorname{tr}\gamma|}{2}\right).
$$
Thus length is controlled by trace. In
particular, many different conjugacy classes
can have the same trace, hence the same
length. For example, in the modular group
$\operatorname{PSL}_2(\mathbb{Z})$, traces
are integers. Therefore the number of
possible trace values up to size
$2\cosh(\ell/2)$ is at most on the order of
$e^{\ell/2}$. Combining this with
Huber--Selberg gives
$$
\langle g(\ell)\rangle
\geq
\frac{e^\ell/\ell}{e^{\ell/2}}
=
\frac{e^{\ell/2}}{\ell}.
$$
This is the basic mechanism: if arithmetic
forces the trace set to be sparse, then many
geodesics must share the same length.

\medskip\noindent
Luo and Sarnak observed that one does not
need the trace set to be literally discrete
in $\mathbb{R}$. It is enough to have the
bounded clustering property:
$$
\#(\mathcal{L}(\Gamma)\cap[N-1,N])
\leq C
$$
for all integers $N$. Here
$\mathcal{L}(\Gamma)$ is the set of absolute
trace values of hyperbolic elements of
$\Gamma$. This condition says that traces
may cluster, but only in a uniformly bounded
way. Bounded clustering is enough to force
exponential growth of the average
multiplicity.

\medskip\noindent
This is where the commutative algebra enters.
Given a Fuchsian group $\Gamma$, its invariant
trace field is
$$
K=
\mathbb{Q}
\bigl(\{\operatorname{tr}(\gamma^2):
\gamma\in\Gamma\}\bigr).
$$
The semiarithmetic condition says that $K$ is
totally real and that the traces are
``integers'' in $K$, meaning
$$
\operatorname{tr}(\gamma)\in\mathcal{O}_K.
$$
Here $\mathcal{O}_K$ is the ring of integers
of the number field $K$. This is exactly the
notion of integrality from commutative
algebra: an element is integral over a ring
if it satisfies a monic polynomial with
coefficients in that ring. Thus
$$
\mathcal{O}_K
=
\{x\in K:
x\text{ satisfies a monic polynomial in }
\mathbb{Z}[t]\}.
$$
Equivalently, $\mathcal{O}_K$ is the
integral closure of $\mathbb{Z}$ inside
$K$. This was the important conceptual point:
these are not ordinary integers, but they
are the elements of $K$ that behave like
integers relative to $\mathbb{Z}$.

\medskip\noindent
The reason this helps is that algebraic
integers in a fixed number field form a
lattice-like object. Since $K$ is totally
real, it has several real embeddings
$$
\sigma_i:K\hookrightarrow\mathbb{R}.
$$
So one trace has several real shadows
$\sigma_i(\operatorname{tr}\gamma)$. One of
these embeddings gives the actual Fuchsian
group acting on $\mathbb{H}^2$, while the
other embeddings provide arithmetic
constraints. In the arithmetic case there is
only one noncompact real factor. In the
semiarithmetic case there may be more than
one noncompact factor, and this number is
called the arithmetic dimension. The proof
uses the fact that the traces lie in
$\mathcal{O}_K$ and that their conjugates are
controlled. By a pigeonhole/lattice-counting
argument this gives bounded clustering of
trace values.

\medskip\noindent
A guiding example is the golden triangle
group
$$
\Gamma=\Delta(2,5,\infty).
$$
Its trace field is
$$
K=\mathbb{Q}(\sqrt{5}),
$$
with two real embeddings:
$$
\sigma_1(\sqrt{5})=\sqrt{5},
\qquad
\sigma_2(\sqrt{5})=-\sqrt{5}.
$$
The first embedding gives the actual
Fuchsian group. The second embedding gives
another real realization, which is not
discrete. This is typical of the
semiarithmetic situation: the group still has
arithmetic trace data, but it is not
arithmetic in the classical sense.

\medskip\noindent
The main story of the talk can therefore be
summarized as follows. Hyperbolic geometry
turns closed geodesics into conjugacy classes.
The length of a closed geodesic is determined
by a trace. If the traces live inside the
integer ring $\mathcal{O}_K$ of a number
field, commutative algebra gives arithmetic
control over how those traces can be
distributed. This control implies bounded
clustering, and bounded clustering forces
large average multiplicities in the length
spectrum. The new theorem says that this
mechanism works not only for arithmetic
surfaces, but also for certain
semiarithmetic surfaces.

\section{Grothendieck--Teichm\"uller theory: from arithmetic geometry to arithmetic topology}
\label{section-grothendieck-teichmuller-arithmetic-topology}

\noindent
Marco Boggi.
Qua 27 mai 2026, 15:30.
SALA 228.

\medskip\noindent
{\bf Abstract.}
Grothendieck--Teichm\"uller theory originated
in arithmetic geometry through the study of
Galois actions on geometric \'etale
fundamental groups. In this talk, we argue
that arithmetic topology provides a natural
setting for applications of
Grothendieck--Teichm\"uller theory. More
specifically, we explain how our theory of
automorphisms of profinite mapping class
groups may lead to new results on profinite
rigidity questions for mapping class groups.

\medskip\noindent
{\bf Warning.}
These are live notes from the talk,
written with unusually heavy AI assistance
since that day I forgot my glasses.
Some theorem
statements may be imprecise; the goal is to
record the ideas and references from the day.\medskip\noindent

\medskip\noindent
The main upshot is that GT theory allows us
to pass the problem to a free group in three
generators, and that the talk aims to
contrast arithmetic and topological
properties.

A more personally significant 
upshot, explained next,
is that the étale fundamental
group of the moduli space of curves
is the profinite completion of
the fundamental group of any curve in 
the moduli, which makes sense since the
underlying topological space of all curves
in the moduli is the same.

\medskip\noindent
Historically, this is very much in the spirit
of Grothendieck's 1984 \emph{Esquisse d'un
programme}. Drinfeld then introduced the
Grothendieck--Teichm\"uller group in 1991,
making precise a framework that was already
expected to have arithmetic applications.

\medskip\noindent
The point is that $\mathcal M_{g,n}$ is
not just an abstract moduli space: it is
built from one fixed topological surface
$S_{g,n}$. A point of $\mathcal M_{g,n}$
is essentially a complex structure on this
same underlying surface.

\medskip\noindent
Thus the topological fundamental group of
the moduli space records how a family of
curves can move around and come back to the
same curve with its underlying surface
changed by a diffeomorphism. This gives the
mapping class group:
$$
\pi_1^{\mathrm{top}}
\bigl(\mathcal M_{g,n}(\mathbb C)\bigr)
\cong
\Gamma_{g,n}.
$$
Here
$$
\Gamma_{g,n}
=
\operatorname{Diff}^+(S_{g,n})/
\operatorname{Diff}_0(S_{g,n}).
$$

\medskip\noindent
The \'etale fundamental group is the
algebraic-geometric version of the
fundamental group. Instead of remembering
all covering spaces, it remembers finite
algebraic covers. Over $\mathbb C$, this
amounts to taking the profinite completion
of the usual fundamental group. Therefore
one gets
$$
\pi_1^{\acute{e}t}
\bigl(\mathcal M_{g,n}
\otimes_{\mathbb Q}\overline{\mathbb Q}\bigr)
\cong
\widehat{\Gamma}_{g,n}.
$$

\medskip\noindent
So the profinite mapping class group is not
an artificial object: it is literally the
\'etale fundamental group of the moduli
space of curves. This is why
Grothendieck--Teichm\"uller theory naturally
acts on profinite mapping class groups.

\medskip\noindent
The starting point is the \'etale
fundamental group of the projective line with
three points removed. Over $\mathbb{C}$, the
space
$$
\mathbb{P}^1_{\mathbb{C}} \setminus \{0,1,\infty\}
$$
is a sphere with three punctures, so its
topological fundamental group is a free group
on two generators:
$$
\pi_1^{\mathrm{top}}
\bigl(\mathbb{P}^1_{\mathbb{C}} \setminus
\{0,1,\infty\}\bigr)
\cong F_2.
$$
After profinite completion one gets
$$
\pi_1^{\acute{e}t}
\bigl(\mathbb{P}^1_{\overline{\mathbb{Q}}} \setminus
\{0,1,\infty\}\bigr)
\cong \widehat{F}_2,
$$
the free profinite group on two generators.

\medskip\noindent
Now let
$$
X=\mathbb{P}^1_{\mathbb{Q}} \setminus
\{0,1,\infty\}.
$$
There is an exact sequence of profinite
groups
$$
1\to
\pi_1^{\acute{e}t}(X_{\overline{\mathbb{Q}}})
\to
\pi_1^{\acute{e}t}(X)
\to
\operatorname{Gal}(\overline{\mathbb{Q}}/\mathbb{Q})
\to 1.
$$
Because the geometric fundamental group is
\(\widehat{F}_2\), the absolute Galois group
acts by outer automorphisms on a free
profinite group:
$$
\operatorname{Gal}(\overline{\mathbb{Q}}/\mathbb{Q})
\longrightarrow
\operatorname{Out}(\widehat{F}_2).
$$
This is one of the first places where
Grothendieck's vision becomes concrete.
The theorem says that this action is
faithful.

\medskip\noindent
In particular, the absolute Galois group can
be studied through its action on the
combinatorics of finite coverings of
\(\mathbb{P}^1\setminus\{0,1,\infty\}\).
This is the origin of dessins d'enfants.
A finite cover of
\(\mathbb{P}^1\setminus\{0,1,\infty\}\)
corresponds to a finite quotient of
\(\widehat{F}_2\), and the Galois action
permutes these quotients in a very rigid way.
That rigidity is what eventually leads to the
Grothendieck--Teichm\"uller group.

\medskip\noindent
The point of the talk is that a similar
picture appears in arithmetic topology.
Mapping class groups also act on profinite
fundamental groups of surfaces, and one can
study their automorphisms in a way that
mirrors the Galois action above. This
suggests profinite rigidity questions: to
what extent does the profinite completion of
a mapping class group remember the original
group? The Grothendieck--Teichm\"uller
framework provides a natural language for
these problems, and may lead to new rigidity
results.

\medskip\noindent
In the arithmetic case there are more
automorphisms than in the topological case,
and this is what makes the theory useful for
applications. The speaker suggested that
these methods should lead to a proof of a
failure of profinite rigidity: more
precisely, if \(d(S)\) denotes the complexity
of the surface, then there should exist
finite-index subgroups \(\Gamma\) and
\(\Gamma'\) of \(\Gamma(S)\) such that
$$
\widehat{\Gamma}\cong \widehat{\Gamma'}
\qquad\text{but}\qquad
\Gamma\not\cong \Gamma'.
$$

\medskip\noindent
For \(g=0\), the speaker can prove that every
pointwise inner automorphism of
\(\widehat{\Gamma}(S)\) is actually inner.
That is, if an automorphism acts on every
single element as conjugation, then the
automorphism itself is conjugation.

\medskip\noindent
As a corollary, the Grothendieck--Teichm\"uller
group acts faithfully on the conjugacy
classes of elements of the profinite
completion of the braid group
\(\widehat{B}_n\). This is interesting from
the topological point of view because of the
Alexander theorem, which characterizes
braids via their relationship with links.

\medskip\noindent
Another application mentioned in the talk is
Furusho's work on finite-type (Vassiliev)
knot theory. The GT action can be used to
organize knot-theoretic information in a way
that is compatible with the profinite braid
group viewpoint.

\medskip\noindent
Let $B$ denote the profinite mapping class
group under consideration. We write
\(\operatorname{Aut}^{\mathbb{I}_0}(B)\) for
the subgroup of \(\operatorname{Aut}(B)\)
consisting of those automorphisms that
preserve the procyclic subgroups generated by
the elements \(C_r\) arising from Dehn
twists. The centralizer of the natural image
of the Grothendieck--Teichm\"uller action in
\(\operatorname{Aut}^{\mathbb{I}_0}(B)\) is
what we call the Grothendieck--Teichm\"uller
group \(G(c,0)\).

\medskip\noindent
For a connected oriented surface
\(S_{g,n}\) of genus \(g\) with \(n\)
punctures and \(\chi(S_{g,n})<0\), the
mapping class group is
the group of orientation-preserving
diffeomorphisms modulo isotopy to the
identity. More precisely, if
\(\Gamma_{g,n}\) denotes the mapping class
group of \(S_{g,n}\), then
\(\Gamma_{g,n}=\mathrm{P}\Gamma(S_{g,n})\)
in the pure case.

\begin{theorem}\label{thm:mapping-class-moduli-space}
There is a natural isomorphism
$$
\Gamma_{g,n}\cong
\pi_1^{\mathrm{top}}
\bigl(\mathcal{M}_{g,n}(\mathbb{C})\bigr),
$$
where \(\mathcal{M}_{g,n}\) is the moduli
space of smooth genus \(g\) curves with
\(n\) marked points.
\end{theorem}

\begin{proof}
This is the classical identification of the
mapping class group with the orbifold
fundamental group of the moduli space. In the
case of ordered marked points, the pure
mapping class group corresponds to the
fundamental group of the pointed moduli
space.
\end{proof}

\medskip\noindent
There is also a standard exact sequence
relating the pure mapping class group to the
full mapping class group:
$$
1\to \widehat{\mathrm{P}\Gamma}_{g,n}
\to \widehat{\Gamma}_{g,n}
\to \Sigma_n
\to 1.
$$
At the discrete level this comes from the
action of the mapping class group on the set
of punctures, and the quotient records the
induced permutation of the punctures.

\begin{remark}
How can we study the Deligne--Mumford
boundary through points on a surface? This
appears to be a standard method discussed in
the talk.
\end{remark}

\medskip\noindent
The theorem implies that the \'etale
fundamental group of the moduli space is the
profinite completion of the mapping class
group:
$$
\pi_1^{\acute{e}t}
\bigl(\mathcal{M}_{g,n}\otimes_{\mathbb{Q}}\overline{\mathbb{Q}}\bigr)
\cong
\widehat{\Gamma}_{g,n}.
$$
In the pure case this is
\(\widehat{\mathrm{P}\Gamma}(S_{g,n})\).

\begin{theorem}[Hoshi--Mochizuki]\label{thm:hoshi-mochizuki-centralizer}
Let \(\Sigma_n\) act on the profinite
mapping class group in the natural way by
permuting the punctures. Then there is a
natural injection
$$
Z_{\operatorname{Out}^{\mathbb{I}_0}
(\widehat{\Gamma}_{g,n})}(\Sigma_n)
\hookrightarrow
Z_{\operatorname{Out}^{\mathbb{I}_0}
(\widehat{\Gamma}_{g,n+1})}(\Sigma_{n+1}),
$$
and this map is an isomorphism for \(n\ge 5\).
\end{theorem}

\begin{proof}
This is the centralizer comparison used in
the study of profinite mapping class groups
and their puncture-permutation symmetries.
The low-puncture cases require separate
analysis, while for \(n\ge 5\) the natural
stabilization map becomes an isomorphism.
\end{proof}

\begin{theorem}[Hoshi--Minamide--Mochizuki]\label{thm:hoshi-minamide-mochizuki}
There is a natural isomorphism
$$
\operatorname{Out}^{\mathbb{I}_0}
\bigl(\mathrm{P}\check\Gamma(S)\bigr)
\cong
\Sigma_n \times \widehat{\operatorname{GT}}.
$$
Equivalently, in the genus zero punctured
sphere case one may write
$$
\operatorname{Out}^{\mathbb{I}_0}
\bigl(\mathrm{P}\check\Gamma_{0,n}\bigr)
\cong
\Sigma_n \times \widehat{\operatorname{GT}}.
$$
\end{theorem}

\begin{proof}
This is the theorem that removes the extra
rigidity hypotheses and identifies the full
outer automorphism group, under the
\(\mathbb{I}_0\)-condition, as the product of
the puncture permutation group and the
profinite Grothendieck--Teichm\"uller group.
\end{proof}

\medskip\noindent
There is also a slogan in the genus-zero
setting:
$$
\operatorname{Out}\bigl(\widehat{\Gamma}_{0,n}\bigr)
\cong \widehat{\operatorname{GT}}
$$
for the appropriate completion and rigidity
condition encoded by the exponent on
\(\operatorname{Out}\). More generally, for
small genus \(g\leq 2\), the speaker proved a
corresponding isomorphism
$$
\operatorname{Out}^{\ast}\bigl(\widehat{\Gamma}_{g,n}\bigr)
\cong
\Sigma_n \times \widehat{\operatorname{GT}},
$$
where the superscript \(\ast\) indicates the
same rigidity condition as in the previous
theorems.



\section{Multiplicity of the fibers in
Lagrangian fibrations}
\label{section-multiplicity-fibers-
lagrangian-fibrations}

\noindent
Misha Verbitsky.
Ter 02 jun 2026, 15:30.

\medskip\noindent
Seja \(M\) um coletor hyperkahler de
holonomia m\'axima e
\(\pi : M \to \mathbf{P}^n\) uma fibração
lagrangiana. Uma fibra \(F\) de \(\pi\) \'
e m\'ultipla de multiplicidade \(p\) se um
pequeno disco transversal a \(F\) encontrar
uma fibra geral de \(\pi\) em \(p > 1\)
pontos. Provarei que as fibras m\'ultiplas
n\~ao ocorrem em codimens\~ao 1 se e
somente se o pullback de uma se\c{c}\~ao de
hiperplano for primitivo (indivis\'ivel) na
cohomologia de \(M\). Em seguida,
explicarei como usar essa observa\c{c}\~ao
para provar que n\~ao h\'a fibras m\'ultiplas
em codimens\~ao 1. Os argumentos se
baseiam em uma vers\~ao forte do
desaparecimento de Kodaira-Nakano, v\'alido
para feixes de linhas com curvatura
semipositiva, e na estabilidade do feixe
tangente. Este \'{e} um trabalho conjunto
com Ljudmila Kamenova.

\medskip\noindent
It's a mystery why any elliptic fibration
over a K3 surface has no multiple fibers.
We will explain how to generalize this result
to higher dimensions. By some
theorem, elliptic fibrations shall
be Lagrangian fibrations.

\medskip\noindent

\begin{theorem}[Calabi--Yau]
\label{theorem-calabi-yau}
A compact, K\"ahler, holomorphically
symplectic manifold admits a unique
hyperk\"ahler metric in any K\"ahler class.
\end{theorem}

\begin{definition}
\label{definition-hyperkahler-manifold}
For the rest of this talk, a hyperk\"ahler
manifold is a compact, K\"ahler,
holomorphically symplectic manifold.
\end{definition}

\begin{definition}
\label{definition-maximal-holonomy}
A hyperk\"ahler manifold \(M\) is called
maximal holonomy, or IHS, if
\(\pi_1(M)=0\) and \(H^{2,0}(M)=\mathbf{C}\).
\end{definition}

\medskip\noindent
\textbf{Bogomolov's decomposition.}
Any hyperk\"ahler manifold admits a finite
covering which is a product of a torus and
several hyperk\"ahler manifolds of maximal
holonomy.

\medskip\noindent
Further on, all hyperk\"ahler manifolds are
assumed to be of maximal holonomy.

\begin{theorem}[Matsushita]
\label{theorem-matsushita-lagrangian}
Let \(M\) be a hyperk\"ahler manifold of
maximal holonomy, and let
\(\pi : M \to X\) be a surjective
holomorphic map, with \(0 < \dim X < \dim M\).
Then \(\pi\) is a Lagrangian fibration
(that is, it has holomorphic Lagrangian
fibers).
\end{theorem}

\begin{theorem}[Hwang]
\label{theorem-hwang-smooth-base}
In these assumptions, \(X\) is biholomorphic
to \(\mathbf{C}P^n\) when it is smooth.
\end{theorem}

\medskip\noindent
The conjecture is that \(X\) is biholomorphic
to \(\mathbf{C}P^n\) when it is normal.

\begin{theorem}[Matsushita]
\label{theorem-matsushita-normal-base}
Let \(M\) be a hyperk\"ahler manifold of
maximal holonomy, and let
\(\pi : M \to X\) be a Lagrangian fibration,
with \(X\) normal. Then
\(H^*(X,\mathbf{Q}) \cong H^*(\mathbf{C}P^n,
\mathbf{Q})\).
\end{theorem}

\begin{remark}
\label{remark-general-fibers-abelian}
General fibers of \(\pi\) are Abelian
varieties (projective complex tori), by
Arnold-Liouville. Conversely, as shown by
Hwang-Weiss, any Lagrangian complex torus in
\(M\) is a fiber of a Lagrangian fibration.
\end{remark}

\begin{remark}
\label{remark-fibration}
A fibration \(\pi : M \to X\) is a proper
surjective holomorphic map of complex
manifolds. We shall always tacitly assume
that the fibration is proper and
equidimensional, that is, the irreducible
components of all fibers have dimension
\(\dim M - \dim X\). A special fiber of
\(\pi\) is a fiber containing a critical
point.
\end{remark}

\begin{definition}
\label{definition-multiplicity-special-fiber}
Let \(Z\) be an irreducible component of a
special fiber of a fibration
\(\pi : M \to X\), and \(z=\pi(Z)\). The
multiplicity of \(Z\) is the rank of
\(\pi^*(\mathcal{O}_X/\mathfrak{m}_z)\) in a
general point of \(Z\), where
\(\mathfrak{m}_z \subset \mathcal{O}_X\) is
the maximal ideal of \(z\).
\end{definition}

\begin{remark}
\label{remark-multiplicity-one-dimensional}
Suppose that \(\dim X=1\). Locally in \(M\),
around a general point of \(m \in Z\), there
is a coordinate system \(x_1,\ldots,x_n\),
such that \(\pi(x_1,\ldots,x_n)=x_1^d\).
In this case, \(d\) is multiplicity of
\(Z\).
\end{remark}

\begin{remark}
\label{remark-equivalent-definition}
Another definition of multiplicity is more
geometric, but it is equivalent to the one
given above; the equivalence is left as an
exercise.
\end{remark}

\begin{definition}
\label{definition-multiplicity-transversal-disk}
Let \(Z\) be an irreducible component of a
special fiber of a fibration
\(\pi : M \to X\), and \(z=\pi(Z)\).
Consider a small disc \(D \subset M\) of
dimension \(\dim X\) transversal to \(Z\)
and intersecting it in \(m\). The
multiplicity of \(Z\) in \(m \in M\) is the
number of intersection points between \(D\)
and a general fiber of \(\pi\) which is
sufficiently close to \(Z\).
\end{definition}

\begin{theorem}[Campana-Kamenova-Verbitsky]
\label{theorem-campana-kamenova-verbitsky}
Let \(\pi : M \to X\) be an Abelian
fibration (that is, a fibration with general
fiber a complex torus), and assume that
\(M\) is K\"ahler. Let \(\mu_i\) be the
multiplicity of irreducible components of
\(\pi^{-1}(z)\) for some \(z \in X\), and
\(\gcd(\mu_i)\) their greatest common
divisor. Then \(\gcd(\mu_i)=\min \mu_i\).
\end{theorem}

\begin{definition}
\label{definition-discriminant}
Let \(\pi : M \to X\) be a surjective
holomorphic map of complex manifolds, and
\(D \subset X\) its set of critical values,
which is known as the discriminant, or the
discriminant divisor (it has codimension
1). We say that \(\pi\) has no multiple
fibers in codimension 1 if for a general
point \(x \in D\), the fiber \(\pi^{-1}(x)\)
has a component with multiplicity 1.
\end{definition}

\begin{theorem}
\label{theorem-elliptic-k3-no-multiple-fibers}
Let \(\pi : M \to X\) be an elliptic
fibration on a K3 surface. Then \(\pi\) has
no multiple fibers.
\end{theorem}

\medskip\noindent
The main result today is a generalization of
this theorem.

\begin{theorem}
\label{theorem-main-no-multiple-fibers}
Let \(\pi : M \to X\) be a Lagrangian
fibration on a hyperk\"ahler manifold. Then
\(\pi\) has no multiple fibers in
codimension 1.
\end{theorem}

\begin{lemma}
\label{lemma-h2-torsion-free}
Let \(M\) be a hyperk\"ahler manifold of
maximal holonomy. Then \(H^2(M)\) is
torsion-free.
\end{lemma}

\begin{proof}
The universal coefficients formula gives the
exact sequence:
$$
0 \to \operatorname{Ext}_{\mathbb{Z}}^1
\bigl(H_1(M;\mathbb{Z}),\mathbb{Z}\bigr)
\to H^2(M;\mathbb{Z}) \to
\operatorname{Hom}_{\mathbb{Z}}
\bigl(H_2(M;\mathbb{Z}),\mathbb{Z}\bigr)
\to 0.
$$
Since \(H_1(M,\mathbb{Z})=0\) for a maximal
holonomy hyperk\"ahler manifold, this gives
an isomorphism
\[
H^2(M,\mathbb{Z}) =
\operatorname{Hom}_{\mathbb{Z}}
\bigl(H_2(M;\mathbb{Z}),\mathbb{Z}\bigr),
\]
hence the torsion vanishes.
\end{proof}

\begin{lemma}
\label{lemma-primitive-classes-lagrangian}
Let \(\pi : M \to \mathbf{C}P^n\) be a
Lagrangian fibration on a hyperk\"ahler
manifold of maximal holonomy, and
let \(H \subset \mathbf{C}P^n\) be a
hyperplane section. Then the following
assertions are equivalent.
\begin{itemize}
\item The homology class of
\(\pi^{-1}(H)\) is primitive.
\item The map \(\pi\) has no multiple
fibers in codimension 1.
\end{itemize}
\end{lemma}

\begin{proof}
Step 1: The implication (ii) \(\Rightarrow\)
(i) is proven later today using the ETMDPS
vanishing theorem. The converse implication:
let \(D_1\) be an irreducible component of
the discriminant \(D\); assume it has
multiplicity \(\mu\). We need to show that
the preimage of the hyperplane section \(H\)
is not primitive.

Step 2: Arguing ad absurdum, assume that
\(D_1\) is homologous to \(k[H]\), but
\([\pi^{-1}(H)]\) is primitive in
\(H^2(M,\mathbb{Z})\). Since
\([D_1] = k[H]\), \(\pi^{-1}(D_1)\) is the
zero set of a section
\(s \in H^0(\mathbf{C}P^n,\mathcal{O}(k))\).
The pullback of this section has multiplicity
\(\mu\), since \(\pi^*(\mathcal{O}(1))\) is
primitive, \(k\) is divisible by \(\mu\).
The universal coefficients formula implies
that \(H^2(M,\mathbb{Z})\) is torsion-free.
Therefore, \(\pi^*s = s_0^\mu\), where
\(s_0 \in H^0(M,\pi^*\mathcal{O}(k/\mu))\).
Since \(\pi^*\mathcal{O}(k/\mu)\) is
trivial on fibers of \(\pi\), the section
\(s_0\) is a pullback of a section
\(s_1 \in H^0(\mathbf{C}P^n,\mathcal{O}(k/\mu))\).
This is impossible, because \(s\) and hence
\(s_1\) vanishes on \(D_1\), which has
degree \(k\).
\end{proof}

\begin{theorem}[Fujiki]
\label{theorem-fujiki}
Let \(\eta \in H^2(M)\), and \(\dim M = 2n\),
where \(M\) is hyperk\"ahler. Then
\[
\int_M \eta^{2n} = cq(\eta,\eta)^n,
\]
for some primitive integer quadratic form
\(q\) on \(H^2(M,\mathbb{Z})\), and
\(c>0\) a rational number.
\end{theorem}

\begin{definition}
\label{definition-bbf-form}
This form is called the Bogomolov-Beauville-
Fujiki form. It is defined by the Fujiki's
relation uniquely, up to a sign. The sign is
determined from the following formula
\[
\lambda q(\eta,\eta)=
\int_X \eta \wedge \eta \wedge \Omega^{n-1}
\wedge \overline{\Omega}^{n-1}
- \frac{n-1}{2n}
\left(
\int_X \eta \wedge \Omega^{n-1}
\wedge \overline{\Omega}^{n}
\right)
\left(
\int_X \eta \wedge \Omega^{n}
\wedge \overline{\Omega}^{n-1}
\right)
\]
where \(\Omega\) is the holomorphic symplectic
form, and \(\lambda>0\).
\end{definition}

\begin{remark}
\label{remark-bbf-signature}
\(q\) has signature \((3,b_2-3)\). It is
negative definite on primitive forms, and
positive definite on \(\langle \Omega,
\overline{\Omega}, \omega \rangle\), where
\(\omega\) is a K\"ahler form.
\end{remark}

\begin{definition}
\label{definition-holomorphic-euler}
Let \(B\) be a holomorphic vector bundle (or
a coherent sheaf). The holomorphic Euler
characteristic is
\[
\chi(B) := \sum_i (-1)^i H^i(M,B).
\]
\end{definition}

\begin{theorem}[Riemann-Roch-Hirzebruch]
\label{theorem-rriemann-roch-hirzebruch}
Let \(M\) be a compact complex manifold, and
\(B\) a holomorphic vector bundle. Then
\(\chi(B)\) can be expressed through the
the Chern classes of \(TM\) and \(B\),
\[
\chi(B)=\int_M td(TM)\wedge ch(B)
\]
where \(td\) is the the Todd polynomial on
Chern classes of \(TM\),
\[
td(M)=1+\frac{1}{2}c_1+\frac{1}{12}(c_1^2+c_2)
+\frac{1}{24}c_1c_2+\frac{1}{720}
(-c_1^4+4c_1^2c_2+c_1c_3+3c_2^2-c_4)+\ldots
\]
and \(ch(B)\) its Chern character,
\[
ch(B)=1+c_1+\frac{1}{2}(c_1^2-2c_2)
+\frac{1}{6}(c_1^3-3c_1c_2+3c_3)+\ldots
\]
\end{theorem}

\noindent
The idea of the following theorem
is to look at $\text{SO}(H^2)$.
The function $\chi(L)$ is an isometry
of BBF form, which makes it a polynomial
on $q(v,v)$. Also, isometries of $H^2$
are correspond with automorphisms of $M$.

\begin{theorem}[Huybrechts]
\label{theorem-huybrechts}
Let \(M\) be a hyperk\"ahler manifold,
\(\dim_{\mathbf{C}} M = 2n\), and \(L\) a
holomorphic line bundle. Then
\[
\chi(L)=\sum a_i q(c_1(L))^i,
\]
where the coefficients \(a_i\) are constants
depending on the topology of \(M\).
\end{theorem}

\begin{proof}
Step 1: Let \(A^*\) be the subalgebra in
cohomology generated by \(H^2(M)\). Then
\[
A^{2i} \cong \operatorname{Sym}^i(H^2(M))
\]
up to the middle degree, and
\[
A^{n+i} \cong \operatorname{Sym}^{n-i}
\bigl(H^2(M)\bigr);
\]
there is an \(O(H^2(M))\)-action on
cohomology, and the multiplication is
\(O(H^2(M))\)-invariant (V., 1995).

Step 2: All Chern classes of \(TM\) are
\(O(H^2(M))\)-invariant, but there is only
one (up to a constant multiplier)
\(O(H^2(M))\)-invariant functional on
\(\operatorname{Sym}^{2i}(H^2(M))\). On the
class \(\eta^{2i} \in H^{4i}(M)\) this
functional takes value \(q(\eta,\eta)^i\).
Therefore, all \(L\)-dependent coefficients
in the Hirzebruch-Riemann-Roch formula for
\(\chi(L)\) are expressed through
\(q(c_1(L))\).
\end{proof}

\begin{lemma}[Corollary]
\label{corollary-huybrechts-chi}
Let \(L\) be a line bundle on a hyperk\"ahler
manifold \(M\), \(\dim_{\mathbf{C}} M = 2n\).
Assume that \(q(c_1(L))=0\). Then
\(\chi(L)=n+1\).
\end{lemma}

\begin{proof}
Indeed, \(\chi(L)=\chi(\mathcal{O}_M)=n+1\),
with the second equality implied by
Bochner's vanishing theorem.
\end{proof}

\begin{definition}
\label{definition-semipositive}
A real \((1,1)\)-form \(\eta\) on a complex
manifold \(M\) is called semipositive if
\(\eta(x,Ix)\geq 0\) for all real tangent
vectors \(x\).
\end{definition}

\medskip\noindent
The following theorem was rediscovered
several times during 1990-ies (Enoki,
Takegoshi, Morougane). Its most general form
(which we do not use) is due to Demailly,
Peternell and Schneider.

It is similar to Kodaira-Nakano
vanishing theorem: instead of having that
$H^i(L \otimes K)$ will vanish for ample $L$,
we have an estimate for how big it can be:

\begin{theorem}[ETMDPS]
\label{theorem-etmdps-vanishing}
Let \((M,I,\omega)\) be a compact K\"ahler
manifold, \(\dim_{\mathbf{C}} M = n\), \(K\)
its canonical bundle, and \(L\) a
holomorphic line bundle on \(M\) equipped
with a Hermitian metric \(h\). Assume that
the curvature \(\Theta\) of \(L\) is a
positive form on \(M\). Then the wedge
multiplication operator \(\eta \mapsto
\omega^i \wedge \eta\) induces a surjective
map
\[
H^0(\Omega^{n-i}_M \otimes L)
\xrightarrow{\ \omega^i\ }
H^i(K \otimes L).
\]
Here \(\omega\) is considered as an element
in \(H^1(\Omega^1_M)\), and multiplication by
\(\omega\) maps
\[
H^k(\Omega^{n-l}_M \otimes L)
to H^{k+1}(\Omega^{n-l+1}_M \otimes L).
\]
\end{theorem}

\begin{lemma}
\label{lemma-global-sections-vanish}
Let \(B\) be a vector bundle equipped with a
filtration \(0 = B_0 \subset B_1 \subset
\ldots \subset B_k = B\). Assume that
\(H^0(B_i/B_{i-1}) = 0\) for all \(i\).
Then \(H^0(B)=0\).
\end{lemma}

\begin{theorem}
\label{theorem-line-bundle-trivial-on-fibers}
Let \(M\) be a hyperk\"ahler manifold
admitting a Lagrangian fibration
\(\pi : M \to X\), and \(H\) a line bundle on
\(X\). Let \(L\) be a line bundle such that
\(L^{\otimes k} = \pi^*H\). Then \(L\) is
trivial on all smooth fibers of \(\pi\).
\end{theorem}

\begin{proof}
Step 1: Let \(F\) be a smooth fiber of
\(\pi\), which is an abelian variety by
Arnol'd-Liouville. Then \(T^*M|_F\) is an
extension of a trivial bundle \(TF\) with
another trivial bundle \(NF = T^*F\). For
any non-trivial line bundle
\(L \in \operatorname{Pic}_0(F)\), we have
\(H^0(L\otimes TF)=0\) and
\(H^0(L\otimes NF)=0\), which implies that
\(H^0(L\otimes T^*M|_F)=0\). Similarly one
\(H^0(L\otimes \wedge^k T^*M|_F)=0\)
(Lemma 2).
Step 2: Unless \(L\) is trivial on \(F\), we
have \(H^0(L\otimes \wedge^k T^*M|_F)=0\),
which implies that \(H^0(L\otimes \wedge^*M)
=0\). By Enoki-Mourugane-Takegoshi-Demailly-
Peternell-Schneider theorem, this implies
that \(H^i(L)=0\), hence \(\chi(L)=0\),
contradicting the formula \(\chi(L)=n+1\).
\end{proof}

Consider a fibration \(M \to \mathbf{C}P^n\)
with fiber \(T^n\). This gives an exact
sequence
\[
0 \to TF \to TM|_F \to \pi^*T\mathbf{C}P^n|_F
\to 0.
\]
The point is that when we have a
decomposition into trivial bundles, it is
very easy to control the sections.
The usage appears to be that a torsion
bundle tensored with a trivial bundle cannot
have sections.

\begin{proposition}
\label{proposition-fiberwise-monodromy-line-bundle}
Let \(M\) be a hyperk\"ahler manifold
admitting a Lagrangian fibration
\(\pi : M \to X\), and \(H\) a line bundle on
\(X\). Let \(L\) be a line bundle such
that \(L^{\otimes k} = \pi^*H\). Then \(L\)
admits a connection \(\nabla\) which is flat
on each restriction \(L|_F\) to the fiber of
\(\pi\).
\end{proposition}

\begin{proof}
Choose a constant metric \(h^k\) on
\(L^{\otimes k}|_F = \mathcal{O}_F\) and let
\(h\) be its \(k\)-th root, which is a metric
on \(L|_F\). Since \(h^k\) is constant on
each fiber, its curvature vanishes, and the
Chern connection \(\nabla\) associated with
\(h\) is also flat.
\end{proof}

\begin{definition}
\label{definition-fiberwise-monodromy}
Fiberwise monodromy of \(L\) is its
monodromy on the fibers of \(\pi\).
\end{definition}

\begin{remark}
\label{remark-pullback-monodromy}
Clearly, \(L = \pi^*L_0\) if and only if the
fiberwise monodromy of \(L\) on each fiber
is trivial.
\end{remark}

\begin{remark}
\label{remark-flat-connection}
Since \(\mathcal{O}(1)\) is flat on
\(\mathbf{C}^n = \mathbf{C}P^n \setminus H\),
we also obtain a flat connection on
\(L|_{\pi^{-1}(\mathbf{C}P^n\setminus H)}\).
The main theorem would follow if we prove
that the monodromy of this flat connection is
trivial on \(\pi^{-1}(z)\), where
\(z \in D\) is a general point of the
discriminant.
\end{remark}

\section{Classification of admissible
W-algebras of low growth}
\label{section-w-algebras-of-low-growth}

\noindent
Igor Alarcon Blatt. June 10, 2026.
Algebra Seminar, IMPA.

\medskip\noindent
{\bf Abstract.} 
Admissible W-algebras form an important class
of vertex algebras, closely related to affine
Kac-Moody algebras and conformal field
theory. A natural invariant of such algebras
is (the asymptotic behaviour of) their
normalised characters. We classify all
admissible W-algebras whose normalised
characters have asymptotic growth less than
1. The proof relies on the computation of the
relevant data together with a theorem of Kac
and Wakimoto, leading to a characterisation
in terms of Virasoro minimal models.

\medskip\noindent
Goal: classify a certain class of W-algebras.
W-algebras are a kind of vertex algebra.

[A lot of content missing.
Some ideas are: minimal models
of Virasoro algebras,
assymptotic data]

\medskip\noindent
Recall that partitions of $n$
correspond with nilpotent orbits.

\begin{theorem}[Arakawa, 2015]
\label{theorem-arakawa}
For $k$ an admissible level with
denominator $u$, there is a unique
nilpotent orbit $\mathbb{O}_n$
suh that $W_k(\mathfrak{g},f)$
is rational for $f \in \overline{
\mathbb{O}_n}$. The orbits $\mathbb{O}_n$
correspond to the partitions listed below.
\end{theorem}

\medskip\noindent
Upshot: use a lemma 
(Kac-Wakimoto) which asserts
that if the assymptotic growth
satisfies $>1$ then there is a direct
sum expression.

\end{document}
